% The IEEE Reference Guide, V 3.28.2025, page by page, with every example typeset by
% biblatex-ieee. Compare it with the original PDF, e.g., using diff-pdf.
% The pages are generated by ieee-reference-guide-replica.py.
\documentclass{article}
\usepackage[letterpaper,margin=1in]{geometry}
\usepackage{fontspec}
\setmainfont{texgyretermes}[Extension=.otf, UprightFont=*-regular, BoldFont=*-bold,
  ItalicFont=*-italic, BoldItalicFont=*-bolditalic]
\setsansfont{texgyreheros}[Extension=.otf, UprightFont=*-regular, BoldFont=*-bold,
  ItalicFont=*-italic, BoldItalicFont=*-bolditalic]
\usepackage{csquotes}
\usepackage[style=ieee,backend=biber]{biblatex}
\usepackage{xurl}
\urlstyle{same}
\usepackage{tikz}
\addbibresource{ieee-reference-guide.bib}
\pagestyle{empty}

% Like \fullcite, but ends the entry as the bibliography does (final period)
\makeatletter
\DeclareCiteCommand{\bibentry}
  {\defcounter{maxnames}{\blx@maxbibnames}}
  {\usedriver{\DeclareNameAlias{sortname}{default}}{\thefield{entrytype}}%
   \finentry}
  {}
  {}
\makeatother

% One page of the guide; coordinates are PDF points from the top left corner
\newenvironment{guidepage}
  {\null\begin{tikzpicture}[remember picture,overlay,x=1bp,y=-1bp,
     shift={(current page.north west)}]}
  {\end{tikzpicture}\newpage}
\newcommand*{\fR}{\rmfamily\mdseries\upshape}
\newcommand*{\fRB}{\rmfamily\bfseries\upshape}
\newcommand*{\fRI}{\rmfamily\mdseries\itshape}
\newcommand*{\fRBI}{\rmfamily\bfseries\itshape}
\newcommand*{\fS}{\sffamily\mdseries\upshape}
\newcommand*{\fSB}{\sffamily\bfseries\upshape}
\newcommand*{\fSI}{\sffamily\mdseries\itshape}
\newcommand*{\fSBI}{\sffamily\bfseries\itshape}
% \gT{x}{baseline}{text}: one line of the guide, made of \gS{font}{size}{color}{text}
\newcommand{\gT}[3]{\node[anchor=base west,inner sep=0pt] at (#1,#2) {#3};}
\newcommand{\gS}[4]{{\fontsize{#2}{#2}\csname f#1\endcsname\selectfont\color{#3}\relax #4}}
% \gB{x}{baseline}: a bullet as large as the guide's
\newcommand{\gB}[2]{\fill (#1+2.6,#2-3.4) circle (1.9);}
% \gE{x}{baseline}{width}{first line indent}{line distance}{key}{appended text}
\newcommand{\gE}[7]{\node[anchor=base west,inner sep=0pt] at (#1,#2)
  {\parbox[t]{#3bp}{\fontsize{10}{#5bp}\fR\selectfont\raggedright
    \hyphenpenalty=10000 \exhyphenpenalty=10000
    \setlength{\parindent}{#4bp}\indent\bibentry{#6}#7}};}
% \gEC{top}{bottom}: \gE clipped to the given part of the page (entry split over pages)
\newcommand{\gEC}[9]{\begin{scope}\clip (0,#1) rectangle (612,#2);
  \gE{#3}{#4}{#5}{#6}{#7}{#8}{#9}\end{scope}}
% Rules and shaded cells of the tables, underlines of links
\newcommand{\gL}[6]{\draw[#1,line width=#2bp] (#3,#4) -- (#5,#6);}
\newcommand{\gR}[6]{\draw[#1,line width=#2bp] (#3,#4) rectangle (#5,#6);}
\newcommand{\gF}[5]{\fill[#1] (#2,#3) rectangle (#4,#5);}

\begin{document}
% Generated from the PDF of the IEEE Reference Guide, V 3.28.2025, by ieee-reference-guide-replica.py.
% Do not edit; regenerate instead.
\definecolor{c000000}{HTML}{000000}
\definecolor{c0070C0}{HTML}{0070C0}
\definecolor{c1155CC}{HTML}{1155CC}
\definecolor{c222222}{HTML}{222222}
\definecolor{c333333}{HTML}{333333}
\definecolor{c434343}{HTML}{434343}
\definecolor{c4A86E8}{HTML}{4A86E8}
\definecolor{cEFEFEF}{HTML}{EFEFEF}
\begin{guidepage}% page 1
\gT{184.83}{106.07}{\gS{RB}{25}{c4A86E8}{REFERENCE GUIDE}}
\gT{72.00}{607.50}{\gS{RB}{13}{c000000}{IEEE Publication Operations }}
\gT{72.00}{622.44}{\gS{RB}{13}{c000000}{445 Hoes Lane }}
\gT{72.00}{637.39}{\gS{RB}{13}{c000000}{Piscataway, NJ 08854 USA }}
\gT{72.00}{650.48}{\gS{R}{11}{c000000}{V 3.28.2025}}
\gT{454.74}{652.34}{\gS{RB}{13}{c000000}{© 2025 IEEE }}
\end{guidepage}
\begin{guidepage}% page 2
\gT{536.00}{43.47}{\gS{R}{8}{c000000}{2 }}
\gT{238.32}{88.80}{\gS{RB}{18}{c000000}{Table of Contents}}
\gT{72.00}{122.15}{\gS{RB}{11}{c000000}{REFERENCE GUIDE}}
\gT{534.50}{122.15}{\gS{RB}{11}{c000000}{1 }}
\gT{72.00}{136.87}{\gS{RB}{10}{c000000}{I. Citing References}}
\gT{535.00}{136.87}{\gS{RB}{10}{c000000}{4}}
\gT{90.00}{151.37}{\gS{R}{10}{c000000}{A. References in Text}}
\gT{535.00}{151.37}{\gS{R}{10}{c000000}{4}}
\gT{90.00}{165.87}{\gS{R}{10}{c000000}{B. Citing Parts of a Reference}}
\gT{535.00}{165.87}{\gS{R}{10}{c000000}{4}}
\gT{72.00}{180.37}{\gS{RB}{10}{c000000}{II. Style}}
\gT{535.00}{180.37}{\gS{RB}{10}{c000000}{4}}
\gT{90.00}{194.86}{\gS{R}{10}{c000000}{A. Blogs }}
\gT{535.00}{194.86}{\gS{R}{10}{c000000}{5 }}
\gT{90.00}{209.36}{\gS{R}{10}{c000000}{B. Books}}
\gT{535.00}{209.36}{\gS{R}{10}{c000000}{5 }}
\gT{108.00}{223.86}{\gS{R}{10}{c000000}{Book}}
\gT{535.00}{223.86}{\gS{R}{10}{c000000}{5}}
\gT{108.00}{238.36}{\gS{R}{10}{c000000}{Books, Monographs (Online)}}
\gT{535.00}{238.36}{\gS{R}{10}{c000000}{5}}
\gT{108.00}{252.86}{\gS{R}{10}{c000000}{Book (Foreign Publication or Translated)}}
\gT{535.00}{252.86}{\gS{R}{10}{c000000}{6}}
\gT{108.00}{267.36}{\gS{R}{10}{c000000}{Book With Chapter Title}}
\gT{535.00}{267.36}{\gS{R}{10}{c000000}{6}}
\gT{108.00}{281.86}{\gS{R}{10}{c000000}{Book With Editor(s)}}
\gT{535.00}{281.86}{\gS{R}{10}{c000000}{6}}
\gT{108.00}{296.36}{\gS{R}{10}{c000000}{Book With Series Title, Volume Title, and Edition}}
\gT{535.00}{296.36}{\gS{R}{10}{c000000}{7}}
\gT{90.00}{310.86}{\gS{R}{10}{c000000}{C. Conferences and Conference Proceedings}}
\gT{535.00}{310.86}{\gS{R}{10}{c000000}{7}}
\gT{108.00}{325.36}{\gS{R}{10}{c000000}{Conference Paper (Paper Presented at a Conference)}}
\gT{535.00}{325.36}{\gS{R}{10}{c000000}{7}}
\gT{108.00}{339.85}{\gS{R}{10}{c000000}{Conference Proceedings in Print (Paper Presented at a Conference)}}
\gT{535.00}{339.85}{\gS{R}{10}{c000000}{8}}
\gT{108.00}{354.35}{\gS{R}{10}{c000000}{Conference Proceedings With DOI}}
\gT{535.00}{354.35}{\gS{R}{10}{c000000}{8}}
\gT{108.00}{368.85}{\gS{R}{10}{c000000}{Conference Proceedings With Editors}}
\gT{535.00}{368.85}{\gS{R}{10}{c000000}{8}}
\gT{108.00}{383.35}{\gS{R}{10}{c000000}{Conference Proceedings With Location}}
\gT{535.00}{383.35}{\gS{R}{10}{c000000}{9}}
\gT{108.00}{397.85}{\gS{R}{10}{c000000}{Conference Proceedings With Series Title, Volume Title, and Edition}}
\gT{535.00}{397.85}{\gS{R}{10}{c000000}{9}}
\gT{108.00}{412.35}{\gS{R}{10}{c000000}{Conference Paper Online}}
\gT{535.00}{412.35}{\gS{R}{10}{c000000}{9}}
\gT{108.00}{426.85}{\gS{R}{10}{c000000}{Conference Proceedings Online}}
\gT{535.00}{426.85}{\gS{R}{10}{c000000}{9}}
\gT{90.00}{441.35}{\gS{R}{10}{c000000}{D. Course}}
\gT{535.00}{441.35}{\gS{R}{10}{c000000}{9}}
\gT{108.00}{455.85}{\gS{R}{10}{c000000}{Coursepack}}
\gT{530.00}{455.85}{\gS{R}{10}{c000000}{10}}
\gT{90.00}{470.35}{\gS{R}{10}{c000000}{E. Datasets}}
\gT{530.00}{470.35}{\gS{R}{10}{c000000}{10}}
\gT{108.00}{484.84}{\gS{R}{10}{c000000}{Online Dataset Reference Using a DOI}}
\gT{530.38}{484.84}{\gS{R}{10}{c000000}{11}}
\gT{108.00}{499.34}{\gS{R}{10}{c000000}{Online Dataset Reference Using a DOI Resolver}}
\gT{530.38}{499.34}{\gS{R}{10}{c000000}{11}}
\gT{108.00}{513.84}{\gS{R}{10}{c000000}{Online Dataset Reference Using a Website Address}}
\gT{530.38}{513.84}{\gS{R}{10}{c000000}{11}}
\gT{90.00}{528.34}{\gS{R}{10}{c000000}{F. Handbooks}}
\gT{530.38}{528.34}{\gS{R}{10}{c000000}{11}}
\gT{90.00}{542.84}{\gS{R}{10}{c000000}{G. Lectures}}
\gT{530.00}{542.84}{\gS{R}{10}{c000000}{12}}
\gT{108.00}{557.34}{\gS{R}{10}{c000000}{Lecture Notes}}
\gT{530.00}{557.34}{\gS{R}{10}{c000000}{12}}
\gT{108.00}{571.84}{\gS{R}{10}{c000000}{Lecture Online}}
\gT{530.00}{571.84}{\gS{R}{10}{c000000}{12}}
\gT{90.00}{586.34}{\gS{R}{10}{c000000}{H. Legal Citations}}
\gT{530.00}{586.34}{\gS{R}{10}{c000000}{12}}
\gT{90.00}{600.84}{\gS{R}{10}{c000000}{I. Manuals}}
\gT{530.00}{600.84}{\gS{R}{10}{c000000}{12}}
\gT{108.00}{615.34}{\gS{R}{10}{c000000}{Manual (Print)}}
\gT{530.00}{615.34}{\gS{R}{10}{c000000}{13}}
\gT{108.00}{629.83}{\gS{R}{10}{c000000}{Manual (Online)}}
\gT{530.00}{629.83}{\gS{R}{10}{c000000}{13}}
\gT{90.00}{644.33}{\gS{R}{10}{c000000}{J. News Articles}}
\gT{530.00}{644.33}{\gS{R}{10}{c000000}{13 }}
\gT{108.00}{658.83}{\gS{R}{10}{c000000}{News Article (Print)}}
\gT{530.00}{658.83}{\gS{R}{10}{c000000}{13 }}
\gT{108.00}{673.33}{\gS{R}{10}{c000000}{News Article (Online)}}
\gT{530.00}{673.33}{\gS{R}{10}{c000000}{13}}
\gT{90.00}{687.83}{\gS{R}{10}{c000000}{K. Online Video (e.g., YouTube)}}
\gT{530.00}{687.83}{\gS{R}{10}{c000000}{14}}
\gT{90.00}{702.33}{\gS{R}{10}{c000000}{L. Patent}}
\gT{530.00}{702.33}{\gS{R}{10}{c000000}{14}}
\gT{108.00}{716.83}{\gS{R}{10}{c000000}{Patent Online}}
\gT{530.00}{716.83}{\gS{R}{10}{c000000}{14}}
\end{guidepage}
\begin{guidepage}% page 3
\gT{536.00}{43.47}{\gS{R}{8}{c000000}{3 }}
\gT{90.00}{84.34}{\gS{R}{10}{c000000}{M. Periodicals}}
\gT{530.00}{84.34}{\gS{R}{10}{c000000}{14}}
\gT{108.00}{98.83}{\gS{R}{10}{c000000}{Periodicals}}
\gT{530.00}{98.83}{\gS{R}{10}{c000000}{15}}
\gT{108.00}{113.33}{\gS{R}{10}{c000000}{Periodical With Article ID}}
\gT{530.00}{113.33}{\gS{R}{10}{c000000}{16}}
\gT{108.00}{127.83}{\gS{R}{10}{c000000}{Periodical With DOI}}
\gT{530.00}{127.83}{\gS{R}{10}{c000000}{16}}
\gT{108.00}{142.33}{\gS{R}{10}{c000000}{Periodical in Early Access}}
\gT{530.00}{142.33}{\gS{R}{10}{c000000}{16}}
\gT{108.00}{156.83}{\gS{R}{10}{c000000}{Periodical in Other Language}}
\gT{530.00}{156.83}{\gS{R}{10}{c000000}{17}}
\gT{108.00}{171.33}{\gS{R}{10}{c000000}{Periodicals Online}}
\gT{530.00}{171.33}{\gS{R}{10}{c000000}{17}}
\gT{108.00}{185.83}{\gS{R}{10}{c000000}{Virtual Journal}}
\gT{530.00}{185.83}{\gS{R}{10}{c000000}{17}}
\gT{90.00}{200.33}{\gS{R}{10}{c000000}{N. Reports}}
\gT{530.00}{200.33}{\gS{R}{10}{c000000}{17}}
\gT{108.00}{214.83}{\gS{R}{10}{c000000}{Report Online}}
\gT{530.00}{214.83}{\gS{R}{10}{c000000}{18}}
\gT{90.00}{229.33}{\gS{R}{10}{c000000}{O. Software}}
\gT{530.00}{229.33}{\gS{R}{10}{c000000}{18}}
\gT{90.00}{243.83}{\gS{R}{10}{c000000}{P. Standards}}
\gT{530.00}{243.83}{\gS{R}{10}{c000000}{20}}
\gT{108.00}{258.32}{\gS{R}{10}{c000000}{Standard Online}}
\gT{530.00}{258.32}{\gS{R}{10}{c000000}{20}}
\gT{90.00}{272.82}{\gS{R}{10}{c000000}{Q. Theses and Dissertations}}
\gT{530.00}{272.82}{\gS{R}{10}{c000000}{20}}
\gT{108.00}{287.32}{\gS{R}{10}{c000000}{Thesis Online}}
\gT{530.00}{287.32}{\gS{R}{10}{c000000}{20}}
\gT{90.00}{301.82}{\gS{R}{10}{c000000}{R. U.S. Government Documents}}
\gT{530.00}{301.82}{\gS{R}{10}{c000000}{21}}
\gT{108.00}{316.32}{\gS{R}{10}{c000000}{Government Online}}
\gT{530.00}{316.32}{\gS{R}{10}{c000000}{21}}
\gT{90.00}{330.82}{\gS{R}{10}{c000000}{S. Unpublished}}
\gT{530.00}{330.82}{\gS{R}{10}{c000000}{21}}
\gT{108.00}{345.32}{\gS{R}{10}{c000000}{Preprint arXiv}}
\gT{530.00}{345.32}{\gS{R}{10}{c000000}{21}}
\gT{90.00}{359.82}{\gS{R}{10}{c000000}{T. Websites}}
\gT{530.00}{359.82}{\gS{R}{10}{c000000}{22}}
\gT{72.00}{375.25}{\gS{RB}{11}{c000000}{III. Notes About Online References}}
\gT{529.01}{375.25}{\gS{RB}{11}{c000000}{22}}
\gT{72.00}{390.90}{\gS{RB}{11}{c000000}{IV. Useful Abbreviations in References}}
\gT{529.01}{390.90}{\gS{RB}{11}{c000000}{23}}
\gT{90.00}{405.61}{\gS{R}{10}{c000000}{A. Common Abbreviations of Words in References}}
\gT{530.00}{405.61}{\gS{R}{10}{c000000}{23}}
\gT{90.00}{420.11}{\gS{R}{10}{c000000}{B. List of Publishers}}
\gT{530.00}{420.11}{\gS{R}{10}{c000000}{28}}
\gT{90.00}{434.61}{\gS{R}{10}{c000000}{C. Abbreviations of Periodicals With Non-English Titles}}
\gT{530.00}{434.61}{\gS{R}{10}{c000000}{44}}
\end{guidepage}
\begin{guidepage}% page 4
\gT{536.00}{43.47}{\gS{R}{8}{c000000}{4 }}
\gT{244.11}{84.36}{\gS{RB}{14}{c0070C0}{I. C}\gS{RB}{9.8}{c0070C0}{ITING}\gS{RB}{14}{c0070C0}{ R}\gS{RB}{9.8}{c0070C0}{EFERENCES}}
\gT{90.00}{113.16}{\gS{RB}{14}{c000000}{A.}\hspace*{4.40bp}\gS{R}{14}{c000000}{References in Text }}
\gT{85.50}{132.43}{\gS{R}{10}{c000000}{References need not be cited in the text. When they are, they appear on the line, in square brackets, inside the }}
\gT{72.00}{144.43}{\gS{R}{10}{c000000}{punctuation. Grammatically, they may be treated as if they were nouns or footnote numbers, e.g.: }}
\gT{85.50}{168.43}{\gS{R}{10}{c000000}{According to [1]; as demonstrated in [2]; as shown by Brown [4], [5]; as mentioned earlier [2], [4], [5], [6], [7], }}
\gT{85.50}{180.43}{\gS{R}{10}{c000000}{[9]; Smith [4] and Brown and Jones [5]; Wood et al. [7] }}
\gT{72.00}{204.43}{\gS{R}{10}{c222222}{Reference ranges in text will not include an en dash. All references will be written out. For example, “[1]–[4]” will }}
\gT{72.00}{216.43}{\gS{R}{10}{c222222}{now be “[1], [2], [3], [4]”. }}
\gT{86.25}{240.43}{\gS{RB}{10}{c000000}{NOTE:}\gS{R}{10}{c000000}{ Use “et al.” when three or more names are given for a reference cited in the text.}}
\gT{90.00}{277.96}{\gS{RB}{14}{c000000}{B.}\hspace*{5.17bp}\gS{R}{14}{c000000}{Citing Parts of a Reference }}
\gT{86.25}{297.23}{\gS{R}{10}{c000000}{Check the reference list for “ibid.” or “op. cit.” These refer to a previous reference and should be eliminated from }}
\gT{72.00}{309.23}{\gS{R}{10}{c000000}{the reference section. In text, cite the earlier reference number and renumber the reference section accordingly. If the }}
\gT{72.00}{321.23}{\gS{R}{10}{c000000}{“ibid.” gives a new page number, or other information, use the following forms: }}
\gT{108.00}{338.23}{\gS{R}{10}{c000000}{[3, Thm. 1]; [3, Lemma 2]; [3, pp. 5–10]; [3, eq. (2)]; [3, Fig. 1]; [3, Appendix I]; [3, Sect. 4.5]; }}
\gT{108.00}{350.23}{\gS{R}{10}{c000000}{[3, Ch. 2, pp. 5–10]; [3, Algorithm 5]. }}
\gT{72.00}{367.23}{\gS{RB}{10}{c000000}{NOTE:}\gS{R}{10}{c000000}{ Editing of references may entail careful renumbering of references, as well as the citations in text.  }}
\gT{279.83}{416.84}{\gS{RB}{14}{c0070C0}{II. S}\gS{RB}{9.8}{c0070C0}{TYLE}}
\gB{90.00}{443.49}
\gT{108.00}{443.49}{\gS{R}{10}{c000000}{Reference numbers are set flush left and form a column of their own, hanging out beyond the body of the }}
\gT{108.00}{455.49}{\gS{R}{10}{c000000}{reference. The reference numbers are on the line, enclosed in square brackets.  }}
\gB{90.00}{468.04}
\gT{108.00}{468.04}{\gS{R}{10}{c000000}{In all references, the given name of the author or editor is abbreviated to the initial only and precedes the }}
\gT{108.00}{480.04}{\gS{R}{10}{c000000}{last name. Do not use commas around Jr., Sr., and III in names: Michael Smith Jr.; Ray Barnett Sr.; Lucas }}
\gT{108.00}{492.04}{\gS{R}{10}{c000000}{Molignaro III.  }}
\gB{90.00}{504.59}
\gT{108.00}{504.59}{\gS{R}{10}{c000000}{IEEE publications must list names of all authors, up to six names. If there are more than six names listed, }}
\gT{108.00}{516.59}{\gS{R}{10}{c000000}{use the primary author’s name followed by “et al.” For non-IEEE publications, “et al.” may be used if }}
\gT{108.00}{528.59}{\gS{R}{10}{c000000}{additional names are not provided.  }}
\gB{90.00}{541.14}
\gT{108.00}{541.14}{\gS{R}{10}{c222222}{All references, except those ending with URLs, will end with a period, including those with a DOI. }}
\gT{108.00}{553.14}{\gS{R}{10}{c222222}{However, if a reference contains both a DOI or accessed date, as well as a URL, the DOI or accessed date }}
\gT{108.00}{565.14}{\gS{R}{10}{c222222}{will be placed first, followed by a period, and then the URL.  }}
\gB{90.00}{577.69}
\gT{108.00}{577.69}{\gS{R}{10}{c000000}{Note that when citing IEEE Transactions, if the issue number or month is not available, research IEEE }}
\gT{108.00}{589.69}{\gS{R}{10}{c000000}{Xplore to see if the publication carries these details. Note that not all IEEE publications include issue }}
\gT{108.00}{601.69}{\gS{R}{10}{c000000}{numbers and months.  }}
\gB{90.00}{614.24}
\gT{108.00}{614.24}{\gS{R}{10}{c000000}{Note that early access articles should use the “Date of Publication.”  }}
\gB{90.00}{626.79}
\gT{108.00}{626.79}{\gS{R}{10}{c000000}{Articles published within a volume should use the volume and issue publication date.  }}
\gB{90.00}{639.34}
\gT{108.00}{639.34}{\gS{R}{10}{c000000}{All references must include a date that comprises at least the year of publication (or year of the accessed }}
\gT{108.00}{651.34}{\gS{R}{10}{c000000}{date if applicable). If no date is provided, use the abbreviation "(n.d.)" in the same position where the date }}
\gT{108.00}{663.34}{\gS{R}{10}{c000000}{would normally be displayed in the reference, preceded and followed by periods, e.g.: }}
\gE{108.00}{687.84}{432.00}{0.00}{13.80}{style:nodate}{}
\end{guidepage}
\begin{guidepage}% page 5
\gT{536.00}{43.47}{\gS{R}{8}{c000000}{5 }}
\gB{90.00}{81.38}
\gT{108.00}{81.38}{\gS{R}{10}{c000000}{For any reference that cites two months for the same issue, the two months should be separated by a slash }}
\gT{108.00}{93.38}{\gS{R}{10}{c000000}{(e.g., Jul./Aug.), followed by the year of publication.  }}
\gB{90.00}{105.93}
\gT{108.00}{105.93}{\gS{R}{10}{c000000}{References may not include all information; please obtain and include relevant information.  }}
\gB{90.00}{118.48}
\gT{108.00}{118.48}{\gS{R}{10}{c000000}{Do not combine references. There must be only one reference with each number.  }}
\gB{90.00}{131.03}
\gT{108.00}{131.03}{\gS{R}{10}{c000000}{If there is a URL included with the print reference, it can be included at the end of the reference (see style }}
\gT{108.00}{143.03}{\gS{R}{10}{c000000}{for online references). }}
\gT{72.00}{189.27}{\gS{R}{14}{c000000}{A. Blogs  }}
\gT{72.00}{215.31}{\gS{RB}{12}{c333333}{Blog}}
\gT{93.75}{231.39}{\gS{RI}{10}{c000000}{Basic Format: }}
\gB{85.50}{245.82}
\gT{103.50}{245.82}{\gS{R}{10}{c000000}{J. K. Author, “Title of the post,” }\gS{RI}{10}{c000000}{Title of the Blog, }\gS{R}{10}{c000000}{Month, Day, Year. [Online]. Available: http://www.url.com}}
\gT{93.75}{261.27}{\gS{RI}{10}{c000000}{Examples: }}
\gB{90.00}{275.70}
\gE{108.00}{275.70}{432.00}{0.00}{13.00}{blog:weissberger}{}
\gB{90.00}{328.10}
\gE{108.00}{328.10}{432.00}{0.00}{13.00}{blog:owens}{}
\gT{72.00}{389.23}{\gS{R}{14}{c000000}{B. Books }}
\gT{72.00}{412.85}{\gS{RB}{12}{c333333}{Book }}
\gT{93.75}{428.94}{\gS{RI}{10}{c000000}{Basic Format:  }}
\gB{85.50}{443.37}
\gT{94.50}{443.37}{\gS{R}{10}{c000000}{J. K. Author, “Title of chapter in the book,” in }\gS{RI}{10}{c000000}{Title of Published Book, x}\gS{R}{10}{c000000}{th ed. City of Publisher, (only U.S. }}
\gT{94.50}{455.37}{\gS{R}{10}{c000000}{State), Country: Abbrev. of Publisher, year, ch. }\gS{RI}{10}{c000000}{x}\gS{R}{10}{c000000}{, sect. }\gS{RI}{10}{c000000}{x}\gS{R}{10}{c000000}{, pp. }\gS{RI}{10}{c000000}{xxx–xxx. }}
\gT{93.75}{470.82}{\gS{RI}{10}{c000000}{Examples:  }}
\gB{85.50}{485.25}
\gE{94.50}{485.25}{445.50}{0.00}{13.80}{book:klaus}{}
\gB{85.50}{497.80}
\gE{94.50}{497.80}{447.46}{0.00}{12.00}{book:stein}{}
\gB{85.50}{522.35}
\gE{94.50}{522.35}{447.76}{0.00}{12.00}{book:myer}{}
\gB{85.50}{546.90}
\gE{94.50}{546.90}{447.52}{0.00}{12.00}{book:abramowitz}{}
\gB{85.50}{571.45}
\gE{94.50}{571.45}{447.85}{0.00}{12.00}{book:moore}{}
\gB{85.50}{596.00}
\gE{94.50}{596.00}{447.44}{0.00}{12.00}{book:westinghouse}{}
\gB{85.50}{620.55}
\gE{94.50}{620.55}{447.28}{0.00}{12.00}{book:histace}{}
\gT{72.00}{658.31}{\gS{RB}{12}{c333333}{Books, Monographs (Online) }}
\gT{93.75}{674.40}{\gS{RI}{10}{c000000}{Basic Format:  }}
\gB{85.50}{685.83}
\gT{93.00}{685.83}{\gS{R}{10}{c000000}{J. K. Author, “Title of chapter in the book,” in }\gS{RI}{10}{c000000}{Title of Published Book}\gS{R}{10}{c000000}{, xth ed. City of Publisher, State, }}
\gT{93.00}{697.83}{\gS{R}{10}{c000000}{Country: Abbrev. of Publisher, year, ch. x, sect. x, pp. xxx–xxx. [Online]. Available: http://www.web.com  }}
\gT{94.50}{713.28}{\gS{RI}{10}{c000000}{Examples:  }}
\end{guidepage}
\begin{guidepage}% page 6
\gT{536.00}{43.47}{\gS{R}{8}{c000000}{6 }}
\gB{85.50}{81.38}
\gE{94.50}{81.38}{447.58}{0.00}{12.00}{book:young}{}
\gB{85.50}{117.93}
\gE{94.50}{117.93}{447.95}{0.00}{12.00}{book:terahertz}{}
\gB{85.50}{142.48}
\gE{94.50}{142.48}{447.50}{0.00}{12.00}{book:kurland}{}
\gT{72.00}{180.25}{\gS{RB}{12}{c333333}{Book (Foreign Publication or Translated) }}
\gT{93.75}{196.33}{\gS{RI}{10}{c000000}{Basic Format (when referencing a foreign publication that has not been translated): }}
\gB{86.25}{210.76}
\gT{93.75}{210.76}{\gS{R}{10}{c000000}{J. K. Author, “Title of chapter in the book,” in }\gS{RI}{10}{c000000}{Title of Published Book}\gS{R}{10}{c000000}{, X. Editor, Ed., }\gS{RI}{10}{c000000}{x}\gS{R}{10}{c000000}{th ed. City of Publisher, }}
\gT{93.75}{222.76}{\gS{R}{10}{c000000}{State (only U.S.), Country: Abbrev. of Publisher (in Language), year, ch. }\gS{RI}{10}{c000000}{x}\gS{R}{10}{c000000}{, sect. }\gS{RI}{10}{c000000}{x}\gS{R}{10}{c000000}{, pp. }\gS{RI}{10}{c000000}{xxx–xxx. }}
\gT{94.50}{238.21}{\gS{RI}{10}{c000000}{Examples: }}
\gB{86.25}{252.64}
\gE{93.75}{252.64}{446.25}{0.00}{13.80}{book:ichiro}{}
\gB{86.25}{265.19}
\gE{93.75}{265.19}{448.66}{0.00}{12.00}{book:pontryagin}{}
\gT{94.50}{311.19}{\gS{RI}{10}{c000000}{Basic Format (when referencing a translated publication): }}
\gB{86.25}{323.74}
\gT{93.75}{323.74}{\gS{R}{10}{c000000}{J. K. Author, “Title of chapter in the book,” in }\gS{RI}{10}{c000000}{Title of Published Book}\gS{R}{10}{c000000}{, trans. J. K. Translator. City of }}
\gT{93.75}{335.24}{\gS{R}{10}{c000000}{Publisher, State (only U.S.), Country: Abbrev. of Publisher, year, ch. }\gS{RI}{10}{c000000}{x}\gS{R}{10}{c000000}{, sect.}\gS{RI}{10}{c000000}{ x}\gS{R}{10}{c000000}{, pp. }\gS{RI}{10}{c000000}{xxx–xxx}\gS{R}{10}{c000000}{. }}
\gT{94.50}{351.74}{\gS{RI}{10}{c000000}{Examples:}}
\gB{86.25}{366.28}
\gE{93.75}{366.28}{448.23}{0.00}{12.00}{book:dessalles}{}
\gB{86.25}{389.32}
\gE{93.75}{389.32}{446.25}{0.00}{13.80}{book:menchu}{}
\gT{72.00}{414.08}{\gS{RB}{12}{c333333}{Book With Chapter Title }}
\gT{93.75}{430.17}{\gS{RI}{10}{c000000}{Basic Format: }}
\gB{86.25}{444.60}
\gT{93.75}{444.60}{\gS{R}{10}{c000000}{J. K. Author, “Title of chapter in the book,” in }\gS{RI}{10}{c000000}{Title of Published Book}\gS{R}{10}{c000000}{, X. Editor, Ed., City of Publisher, State }}
\gT{93.75}{456.60}{\gS{R}{10}{c000000}{(only U.S.), Country: Abbrev. of Publisher, year, ch. }\gS{RI}{10}{c000000}{x}\gS{R}{10}{c000000}{, sect. }\gS{RI}{10}{c000000}{x}\gS{R}{10}{c000000}{, pp. }\gS{RI}{10}{c000000}{xxx–xxx. }}
\gT{93.75}{472.05}{\gS{RI}{10}{c000000}{Examples: }}
\gB{86.25}{486.48}
\gE{93.75}{486.48}{448.18}{0.00}{12.00}{book:ogura}{}
\gB{86.25}{511.02}
\gE{93.75}{511.02}{448.73}{0.00}{12.00}{book:li}{}
\gT{72.00}{563.79}{\gS{RB}{12}{c333333}{Book With Editor(s) }}
\gT{93.75}{579.87}{\gS{RI}{10}{c000000}{Basic Format: }}
\gB{86.25}{594.30}
\gT{93.75}{594.30}{\gS{R}{10}{c000000}{J. K. Author, “Title of chapter in the book,” in }\gS{RI}{10}{c000000}{Title of Published Book}\gS{R}{10}{c000000}{, X. Editor, Ed., City of Publisher, State }}
\gT{93.75}{606.31}{\gS{R}{10}{c000000}{(only U.S.), Country: Abbrev. of Publisher, year, pp. }\gS{RI}{10}{c000000}{xxx–xxx. }}
\gB{86.25}{618.85}
\gT{93.75}{618.85}{\gS{R}{10}{c000000}{X. Editor, Ed. }\gS{RI}{10}{c000000}{Title of Published Book}\gS{R}{10}{c000000}{. City of Publisher, State (only U.S.), Country: Abbrev. of Publisher, year.}}
\gT{93.75}{634.30}{\gS{RI}{10}{c000000}{Examples: }}
\gB{86.25}{648.73}
\gE{93.75}{648.73}{448.21}{0.00}{12.00}{book:stein}{}
\gB{86.25}{673.28}
\gE{93.75}{673.28}{448.34}{0.00}{12.00}{book:bennett}{}
\gB{86.25}{709.83}
\gEC{0}{714.00}{93.75}{709.83}{448.00}{0.00}{13.80}{book:leonard}{}
\end{guidepage}
\begin{guidepage}% page 7
\gF{cEFEFEF}{102.75}{336.86}{196.50}{351.86}
\gF{cEFEFEF}{196.50}{336.86}{313.50}{351.86}
\gF{cEFEFEF}{313.50}{336.86}{416.25}{351.86}
\gF{cEFEFEF}{416.25}{336.86}{509.25}{351.86}
\gF{cEFEFEF}{102.75}{366.86}{196.50}{381.86}
\gF{cEFEFEF}{196.50}{366.86}{313.50}{381.86}
\gF{cEFEFEF}{313.50}{366.86}{416.25}{381.86}
\gF{cEFEFEF}{416.25}{366.86}{509.25}{381.86}
\gF{cEFEFEF}{102.75}{396.86}{196.50}{411.86}
\gF{cEFEFEF}{196.50}{396.86}{313.50}{411.86}
\gF{cEFEFEF}{313.50}{396.86}{416.25}{411.86}
\gF{cEFEFEF}{416.25}{396.86}{509.25}{411.86}
\gF{cEFEFEF}{102.75}{426.86}{196.50}{441.86}
\gF{cEFEFEF}{196.50}{426.86}{313.50}{441.86}
\gF{cEFEFEF}{313.50}{426.86}{416.25}{441.86}
\gF{cEFEFEF}{416.25}{426.86}{509.25}{441.86}
\gF{cEFEFEF}{102.75}{456.86}{196.50}{471.86}
\gF{cEFEFEF}{196.50}{456.86}{313.50}{471.86}
\gF{cEFEFEF}{313.50}{456.86}{416.25}{471.86}
\gF{cEFEFEF}{416.25}{456.86}{509.25}{471.86}
\gF{cEFEFEF}{102.75}{486.86}{196.50}{501.86}
\gF{cEFEFEF}{196.50}{486.86}{313.50}{501.86}
\gF{cEFEFEF}{313.50}{486.86}{416.25}{501.86}
\gF{cEFEFEF}{416.25}{486.86}{509.25}{501.86}
\gL{c000000}{1.00}{102.50}{336.00}{102.50}{502.00}
\gL{c000000}{1.00}{196.50}{336.00}{196.50}{502.00}
\gL{c000000}{1.00}{313.50}{336.00}{313.50}{502.00}
\gL{c000000}{1.00}{416.50}{336.00}{416.50}{502.00}
\gL{c000000}{1.00}{509.50}{336.00}{509.50}{502.00}
\gL{c000000}{1.00}{102.00}{336.50}{510.00}{336.50}
\gL{c000000}{1.00}{102.00}{351.50}{510.00}{351.50}
\gL{c000000}{1.00}{102.00}{366.50}{510.00}{366.50}
\gL{c000000}{1.00}{102.00}{381.50}{510.00}{381.50}
\gL{c000000}{1.00}{102.00}{396.50}{510.00}{396.50}
\gL{c000000}{1.00}{102.00}{411.50}{510.00}{411.50}
\gL{c000000}{1.00}{102.00}{426.50}{510.00}{426.50}
\gL{c000000}{1.00}{102.00}{441.50}{510.00}{441.50}
\gL{c000000}{1.00}{102.00}{456.50}{510.00}{456.50}
\gL{c000000}{1.00}{102.00}{471.50}{510.00}{471.50}
\gL{c000000}{1.00}{102.00}{486.50}{510.00}{486.50}
\gL{c000000}{1.00}{102.00}{501.50}{510.00}{501.50}
\gT{536.00}{43.47}{\gS{R}{8}{c000000}{7 }}
\gEC{70.83}{792}{93.75}{67.03}{448.00}{0.00}{13.80}{book:leonard}{}
\gT{72.00}{109.60}{\gS{RB}{12}{c333333}{Book With Series Title, Volume Title, and Edition }}
\gT{93.75}{125.68}{\gS{RI}{10}{c000000}{Basic Format: }}
\gB{86.25}{140.11}
\gT{93.75}{140.11}{\gS{R}{10}{c000000}{J. K. Author, “Title of chapter in the book,” in }\gS{RI}{10}{c000000}{Title of Published Book }\gS{R}{10}{c000000}{(Title of the Series), X. Editor, Ed., }\gS{RI}{10}{c000000}{x}\gS{R}{10}{c000000}{th }}
\gT{93.75}{152.11}{\gS{R}{10}{c000000}{ed. City of Publisher, State (only U.S.), Country: Abbrev. of Publisher, year, ch. }\gS{RI}{10}{c000000}{x}\gS{R}{10}{c000000}{, sect. }\gS{RI}{10}{c000000}{x}\gS{R}{10}{c000000}{, pp. }\gS{RI}{10}{c000000}{xxx–xxx. }}
\gT{93.75}{167.56}{\gS{RI}{10}{c000000}{Examples:}}
\gB{86.25}{181.99}
\gE{93.75}{181.99}{448.04}{0.00}{12.00}{book:feistauer}{}
\gB{86.25}{206.54}
\gE{93.75}{206.54}{448.34}{0.00}{12.00}{book:abramowitz}{}
\gB{86.25}{231.09}
\gE{93.75}{231.09}{448.59}{0.00}{12.00}{book:moore}{}
\gT{72.00}{283.33}{\gS{R}{14}{c000000}{C. Conferences and Conference Proceedings  }}
\gT{72.00}{298.19}{\gS{R}{10}{c000000}{The general form for citing conference proceedings is to list the author and title of the paper, followed by the name }}
\gT{72.00}{310.19}{\gS{R}{10}{c000000}{of the conference using the following abbreviations as well as those found in the “Common Abbreviations of Words }}
\gT{72.00}{322.19}{\gS{R}{10}{c000000}{in References” list. }}
\gT{108.75}{347.56}{\gS{R}{10}{c000000}{Annals  }}
\gT{202.50}{347.56}{\gS{R}{10}{c000000}{Ann.  }}
\gT{319.50}{347.56}{\gS{R}{10}{c000000}{Proceedings  }}
\gT{422.25}{347.56}{\gS{R}{10}{c000000}{Proc.  }}
\gT{108.75}{362.56}{\gS{R}{10}{c000000}{Annual  }}
\gT{202.50}{362.56}{\gS{R}{10}{c000000}{Annu.  }}
\gT{319.50}{362.56}{\gS{R}{10}{c000000}{Record  }}
\gT{422.25}{362.56}{\gS{R}{10}{c000000}{Rec.  }}
\gT{108.75}{377.56}{\gS{R}{10}{c000000}{Colloquium  }}
\gT{202.50}{377.56}{\gS{R}{10}{c000000}{Colloq.  }}
\gT{319.50}{377.56}{\gS{R}{10}{c000000}{Symposium  }}
\gT{422.25}{377.56}{\gS{R}{10}{c000000}{Symp.  }}
\gT{108.75}{392.56}{\gS{R}{10}{c000000}{Conference  }}
\gT{202.50}{392.56}{\gS{R}{10}{c000000}{Conf.  }}
\gT{319.50}{392.56}{\gS{R}{10}{c000000}{Technical Digest  }}
\gT{422.25}{392.56}{\gS{R}{10}{c000000}{Tech. Dig.  }}
\gT{108.75}{407.56}{\gS{R}{10}{c000000}{Congress  }}
\gT{202.50}{407.56}{\gS{R}{10}{c000000}{Congr.  }}
\gT{319.50}{407.56}{\gS{R}{10}{c000000}{Technical Paper  }}
\gT{422.25}{407.56}{\gS{R}{10}{c000000}{Tech. Paper  }}
\gT{108.75}{422.56}{\gS{R}{10}{c000000}{Convention  }}
\gT{202.50}{422.56}{\gS{R}{10}{c000000}{Conv.  }}
\gT{319.50}{422.56}{\gS{R}{10}{c000000}{Workshop }}
\gT{422.25}{422.56}{\gS{R}{10}{c000000}{Workshop }}
\gT{108.75}{437.56}{\gS{R}{10}{c000000}{Digest  }}
\gT{202.50}{437.56}{\gS{R}{10}{c000000}{Dig.  }}
\gT{319.50}{437.56}{\gS{R}{10}{c000000}{First  }}
\gT{422.25}{437.56}{\gS{R}{10}{c000000}{1st  }}
\gT{108.75}{452.56}{\gS{R}{10}{c000000}{Exposition  }}
\gT{202.50}{452.56}{\gS{R}{10}{c000000}{Expo.  }}
\gT{319.50}{452.56}{\gS{R}{10}{c000000}{Second  }}
\gT{422.25}{452.56}{\gS{R}{10}{c000000}{2nd  }}
\gT{108.75}{467.56}{\gS{R}{10}{c000000}{International  }}
\gT{202.50}{467.56}{\gS{R}{10}{c000000}{Int.  }}
\gT{319.50}{467.56}{\gS{R}{10}{c000000}{Third  }}
\gT{422.25}{467.56}{\gS{R}{10}{c000000}{3rd  }}
\gT{108.75}{482.56}{\gS{R}{10}{c000000}{Meeting  }}
\gT{202.50}{482.56}{\gS{R}{10}{c000000}{Meeting  }}
\gT{319.50}{482.56}{\gS{R}{10}{c000000}{Fourth/nth...  }}
\gT{422.25}{482.56}{\gS{R}{10}{c000000}{4th/nth…  }}
\gT{108.75}{497.56}{\gS{R}{10}{c000000}{National  }}
\gT{202.50}{497.56}{\gS{R}{10}{c000000}{Nat.  }}
\gT{72.00}{522.94}{\gS{R}{10}{c000000}{Write out all the remaining words, but omit most articles and prepositions like “of the” and “on.” That is, }}
\gT{72.00}{534.94}{\gS{RI}{10}{c000000}{Proceedings of the 1996 Robotics and Automation Conference }\gS{R}{10}{c000000}{becomes }\gS{RI}{10}{c000000}{Proc. 1996 Robot. Automat. Conf. }\gS{R}{10}{c000000}{If an }}
\gT{72.00}{546.94}{\gS{R}{10}{c000000}{ordinal number is in the conference name, use the numerical form instead of spelling it out (e.g., “1st” instead of }}
\gT{72.00}{558.94}{\gS{R}{10}{c000000}{“First”). Include the location if given. For U.S. locations, “USA” must be included after city and state. }}
\gT{72.00}{570.94}{\gS{RB}{10}{c000000}{NOTE: }\gS{R}{10}{c000000}{All published conference or proceedings papers have page numbers.  }}
\gT{72.00}{596.20}{\gS{RB}{12}{c333333}{Conference Paper (Paper Presented at a Conference) }}
\gT{93.75}{612.29}{\gS{RI}{10}{c000000}{Basic Format: }}
\gB{86.25}{626.72}
\gT{93.75}{626.72}{\gS{R}{10}{c000000}{J. K. Author, “Title of paper,” presented at the Abbreviated Name of Conf., City of Conf., Abbrev. State, }}
\gT{93.75}{638.71}{\gS{R}{10}{c000000}{Country, Month and day(s), year, Paper number. }}
\gT{93.75}{654.16}{\gS{RI}{10}{c000000}{Examples: }}
\gB{86.25}{668.60}
\gE{93.75}{668.60}{448.74}{0.00}{12.00}{conf:caratelli}{}
\gB{86.25}{693.14}
\gE{93.75}{693.14}{448.00}{0.00}{12.00}{conf:kreifeldt}{}
\end{guidepage}
\begin{guidepage}% page 8
\gT{536.00}{43.47}{\gS{R}{8}{c000000}{8 }}
\gB{86.25}{81.38}
\gE{93.75}{81.38}{448.18}{0.00}{12.00}{conf:juette}{}
\gB{86.25}{105.93}
\gE{93.75}{105.93}{448.66}{0.00}{12.00}{conf:arrillaga}{}
\gT{72.00}{146.70}{\gS{RB}{12}{c333333}{Conference Proceedings in Print (Paper Presented at a Conference) }}
\gT{93.75}{162.78}{\gS{RI}{10}{c000000}{Basic Format: }}
\gB{86.25}{174.21}
\gT{93.75}{174.21}{\gS{R}{10}{c000000}{J. K. Author, “Title of paper,” in }\gS{RI}{10}{c000000}{Abbreviated Name of Conf.}\gS{R}{10}{c000000}{, (location of conference is optional), (Month and }}
\gT{93.75}{186.21}{\gS{R}{10}{c000000}{day(s) if provided) year, pp.}\gS{RI}{10}{c000000}{ xxx-xxx.}}
\gT{93.75}{201.66}{\gS{RI}{10}{c000000}{Examples: }}
\gB{86.25}{216.09}
\gE{93.75}{216.09}{448.60}{0.00}{12.00}{conf:amador}{}
\gB{86.25}{240.64}
\gE{93.75}{240.64}{448.31}{0.00}{11.00}{conf:faulhaber}{ \textit{*** If the year is given in the conference title, it may be omitted from the end of the }\textit{reference as shown here.***}}
\gB{86.25}{274.52}
\gE{93.75}{274.52}{448.28}{0.00}{12.00}{conf:bingulac}{}
\gB{86.25}{299.07}
\gE{93.75}{299.07}{448.30}{0.00}{12.00}{conf:doyle}{}
\gB{86.25}{323.62}
\gE{93.75}{323.62}{446.25}{0.00}{12.00}{conf:meadow}{}
\gB{86.25}{348.17}
\gE{93.75}{348.17}{448.70}{0.00}{12.00}{conf:parks}{}
\gB{86.25}{372.72}
\gE{93.75}{372.72}{446.25}{0.00}{13.80}{conf:hsia}{}
\gT{72.00}{401.48}{\gS{RB}{12}{c333333}{Conference Proceedings With DOI }}
\gT{93.75}{417.57}{\gS{RI}{10}{c000000}{Basic Format: }}
\gB{86.25}{429.00}
\gT{93.75}{429.00}{\gS{R}{10}{c000000}{J. K. Author, “Title of paper,” in }\gS{RI}{10}{c000000}{Abbreviated Name of Conf.}\gS{R}{10}{c000000}{, (location of conference is optional), year, pp.}}
\gT{93.75}{441.00}{\gS{RI}{10}{c000000}{xxx–xxx}\gS{R}{10}{c000000}{, doi: }\gS{RI}{10}{c000000}{xxx. }}
\gT{93.75}{456.45}{\gS{RI}{10}{c000000}{Examples: }}
\gB{86.25}{470.88}
\gE{93.75}{470.88}{448.48}{0.00}{12.00}{conf:veruggio}{}
\gB{86.25}{495.43}
\gE{93.75}{495.43}{448.35}{0.00}{12.00}{conf:zhao}{}
\gT{72.00}{548.19}{\gS{RB}{12}{c333333}{Conference Proceedings With Editors }}
\gT{93.75}{564.28}{\gS{RI}{10}{c000000}{Basic Format: }}
\gB{86.25}{575.71}
\gT{93.75}{575.71}{\gS{R}{10}{c000000}{J. K. Author, “Title of paper,” in }\gS{RI}{10}{c000000}{Abbreviated Name of Conf.}\gS{R}{10}{c000000}{, X. Editor, Ed. (location of conference is optional), }}
\gT{93.75}{587.71}{\gS{R}{10}{c000000}{year, pp.}\gS{RI}{10}{c000000}{ xxx-xxx. }}
\gT{93.75}{603.16}{\gS{RI}{10}{c000000}{Examples: }}
\gB{86.25}{617.59}
\gE{93.75}{617.59}{448.60}{0.00}{12.00}{conf:amadoreditors}{}
\gB{86.25}{654.14}
\gE{93.75}{654.14}{448.23}{0.00}{12.00}{conf:armstrong}{}
\gB{86.25}{690.69}
\gE{93.75}{690.69}{448.23}{0.00}{12.00}{conf:funes}{}
\end{guidepage}
\begin{guidepage}% page 9
\gT{536.00}{43.47}{\gS{R}{8}{c000000}{9 }}
\gT{72.00}{82.60}{\gS{RB}{12}{c333333}{Conference Proceedings With Location }}
\gT{93.75}{98.68}{\gS{RI}{10}{c000000}{Basic Format: }}
\gB{86.25}{110.11}
\gT{93.75}{110.11}{\gS{R}{10}{c000000}{J. K. Author, “Title of paper,” in }\gS{RI}{10}{c000000}{Abbreviated Name of Conf.}\gS{R}{10}{c000000}{, City, State, Country, year, pp.}\gS{RI}{10}{c000000}{ xxx–xxx. }}
\gT{93.75}{125.56}{\gS{RI}{10}{c000000}{Examples: }}
\gB{86.25}{139.99}
\gE{93.75}{139.99}{448.56}{0.00}{12.00}{conf:carmichael}{}
\gB{86.25}{176.54}
\gE{93.75}{176.54}{448.60}{0.00}{12.00}{conf:sarkar}{}
\gT{72.00}{217.31}{\gS{RB}{12}{c333333}{Conference Proceedings With Series Title, Volume Title, and Edition }}
\gT{93.75}{233.39}{\gS{RI}{10}{c000000}{Basic Format: }}
\gB{86.25}{244.82}
\gT{93.75}{244.82}{\gS{R}{10}{c000000}{J. K. Author, “Title of paper,” in }\gS{RI}{10}{c000000}{Abbreviated Name of Conf.}\gS{R}{10}{c000000}{ in Volume Title, in Series Title, ed., year, pp.}}
\gT{93.75}{255.27}{\gS{RI}{10}{c000000}{xxx-xxx. }}
\gT{94.50}{269.60}{\gS{RI}{10}{c000000}{Example: }}
\gB{86.25}{284.03}
\gE{93.75}{284.03}{448.65}{0.00}{12.00}{conf:amadorseries}{}
\gT{72.00}{333.80}{\gS{RB}{12}{c333333}{Conference Paper Online }}
\gT{93.75}{349.88}{\gS{RI}{10}{c000000}{Basic Format: }}
\gB{86.25}{361.31}
\gT{93.75}{361.31}{\gS{R}{10}{c000000}{J. K. Author. (Date). Title. Presented at Abbreviated Conf. title. [Type of Medium]. Available: site/path/file }}
\gT{94.50}{376.76}{\gS{RI}{10}{c000000}{Examples: }}
\gB{86.25}{391.19}
\gE{93.75}{391.19}{448.73}{0.00}{12.00}{conf:process}{}
\gB{86.25}{427.74}
\gE{93.75}{427.74}{448.11}{0.00}{12.00}{conf:taylor}{}
\gB{86.25}{464.29}
\gE{93.75}{464.29}{448.43}{0.00}{12.00}{conf:chandrasekaran}{}
\gT{72.00}{514.05}{\gS{RB}{12}{c333333}{Conference Proceedings Online }}
\gT{93.75}{530.14}{\gS{RI}{10}{c000000}{Basic Format: }}
\gB{86.25}{541.57}
\gT{93.75}{541.57}{\gS{R}{10}{c000000}{J. K. Author, “Title of paper,” in }\gS{RI}{10}{c000000}{Abbreviated Name of Conf.}\gS{R}{10}{c000000}{, (location of conference is optional), year, pp.}}
\gT{93.75}{553.57}{\gS{RI}{10}{c000000}{xxx-xxx}\gS{R}{10}{c000000}{. [Online]. Available: http://www.url.com}}
\gT{93.75}{569.02}{\gS{RI}{10}{c000000}{Examples: }}
\gB{86.25}{583.45}
\gE{93.75}{583.45}{448.73}{0.00}{12.00}{conf:schubert}{}
\gB{86.25}{620.00}
\gE{93.75}{620.00}{448.51}{0.00}{12.00}{conf:yanamadala}{}
\gT{72.00}{684.24}{\gS{R}{14}{c000000}{D. Course }}
\gT{93.75}{699.55}{\gS{RI}{10}{c000000}{Basic Format: }}
\gT{93.75}{715.43}{\gS{R}{10}{c000000}{Name of University. (Year). Title of course. [Online]. Available: URL}}
\end{guidepage}
\begin{guidepage}% page 10
\gL{c1155CC}{1.00}{72.00}{328.50}{258.00}{328.50}
\gL{c1155CC}{1.00}{72.00}{656.50}{381.00}{656.50}
\gT{532.01}{43.47}{\gS{R}{8}{c000000}{10 }}
\gT{93.75}{79.28}{\gS{RI}{10}{c000000}{Example: }}
\gB{86.25}{93.71}
\gE{93.75}{93.71}{448.64}{0.00}{12.00}{course:argosy}{}
\gT{72.00}{134.48}{\gS{RB}{12}{c333333}{Coursepack }}
\gT{93.75}{150.56}{\gS{RI}{10}{c000000}{Basic Format: }}
\gB{86.25}{164.99}
\gT{93.75}{164.99}{\gS{R}{10}{c000000}{J. K. Instructor. Title of coursepack. (Semester). Title of course. University/Publisher location: }}
\gT{93.75}{176.99}{\gS{R}{10}{c000000}{University/Publisher name.}}
\gT{93.75}{192.44}{\gS{RI}{10}{c000000}{Example: }}
\gB{86.25}{206.87}
\gE{93.75}{206.87}{448.34}{0.00}{12.00}{course:oden}{}
\gT{72.00}{281.11}{\gS{R}{14}{c000000}{E. Datasets  }}
\gT{72.00}{302.96}{\gS{RB}{10}{c000000}{FORCE11 Data Citation Principles}}
\gT{72.00}{314.96}{\gS{R}{10}{c000000}{IEEE }}
\gT{105.90}{314.96}{\gS{R}{10}{c000000}{follows }}
\gT{147.16}{314.96}{\gS{R}{10}{c000000}{the }}
\gT{170.87}{314.96}{\gS{R}{10}{c000000}{FORCE11 }}
\gT{224.58}{314.96}{\gS{R}{10}{c000000}{Data }}
\gT{254.96}{314.96}{\gS{R}{10}{c000000}{Citation }}
\gT{298.67}{314.96}{\gS{R}{10}{c000000}{Principles. }}
\gT{352.65}{314.96}{\gS{R}{10}{c000000}{These }}
\gT{388.02}{314.96}{\gS{R}{10}{c000000}{principles }}
\gT{438.95}{314.96}{\gS{R}{10}{c000000}{are }}
\gT{462.65}{314.96}{\gS{R}{10}{c000000}{as }}
\gT{482.48}{314.96}{\gS{R}{10}{c000000}{follows }}
\gT{523.73}{314.96}{\gS{R}{10}{c000000}{(per }}
\gT{72.00}{326.96}{\gS{R}{10}{c1155CC}{https://www.force11.org/datacitationprinciples}\gS{R}{10}{c000000}{): }}
\gT{90.00}{351.47}{\gS{R}{10}{c000000}{1)}\gS{RB}{10}{c000000}{ }\gS{RB}{10}{c000000}{Importance:}\gS{R}{10}{c000000}{ Data should be considered legitimate, citable products of research. Data citations should be }}
\gT{108.00}{363.47}{\gS{R}{10}{c000000}{accorded the same importance in the scholarly record as citations of other research objects, such as }}
\gT{108.00}{375.47}{\gS{R}{10}{c000000}{publications. }}
\gT{90.00}{387.98}{\gS{R}{10}{c000000}{2)}\gS{RB}{10}{c000000}{ }\gS{RB}{10}{c000000}{Credit and Attribution:}\gS{R}{10}{c000000}{ Data citations should facilitate giving scholarly credit and normative and legal }}
\gT{108.00}{399.98}{\gS{R}{10}{c000000}{attribution to all contributors to the data, recognizing that a single style or mechanism of attribution may }}
\gT{108.00}{411.98}{\gS{R}{10}{c000000}{not be applicable to all data. }}
\gT{90.00}{424.48}{\gS{R}{10}{c000000}{3)}\gS{RB}{10}{c000000}{ }\gS{RB}{10}{c000000}{Evidence:}\gS{R}{10}{c000000}{ In scholarly literature, whenever and wherever a claim relies upon data, the corresponding data }}
\gT{108.00}{436.48}{\gS{R}{10}{c000000}{should be cited. }}
\gT{90.00}{448.99}{\gS{R}{10}{c000000}{4)}\gS{RB}{10}{c000000}{ }\gS{RB}{10}{c000000}{Unique Identification:}\gS{R}{10}{c000000}{ A data citation should include a persistent method for identification that is machine }}
\gT{108.00}{460.99}{\gS{R}{10}{c000000}{actionable, globally unique, and widely used by a community. }}
\gT{90.00}{473.49}{\gS{R}{10}{c000000}{5)}\gS{RB}{10}{c000000}{ }\gS{RB}{10}{c000000}{Access:}\gS{R}{10}{c000000}{ Data citations should facilitate access to the data themselves and to such associated metadata, }}
\gT{108.00}{485.49}{\gS{R}{10}{c000000}{documentation, code, and other materials, as are necessary for both humans and machines to make }}
\gT{108.00}{497.49}{\gS{R}{10}{c000000}{informed use of the referenced data. }}
\gT{90.00}{510.00}{\gS{R}{10}{c000000}{6)}\gS{RB}{10}{c000000}{ }\gS{RB}{10}{c000000}{Persistence:}\gS{R}{10}{c000000}{ Unique identifiers, and metadata describing the data, and its disposition, should persist—even }}
\gT{108.00}{522.00}{\gS{R}{10}{c000000}{beyond the lifespan of the data they describe. }}
\gT{90.00}{534.51}{\gS{R}{10}{c000000}{7)}\gS{RB}{10}{c000000}{ }\gS{RB}{10}{c000000}{Specificity and Verifiability:}\gS{R}{10}{c000000}{ Data citations should facilitate identification of, access to, and verification of }}
\gT{108.00}{546.51}{\gS{R}{10}{c000000}{the specific data that support a claim. Citations or citation metadata should include information about }}
\gT{108.00}{558.51}{\gS{R}{10}{c000000}{provenance and fixity sufficient to facilitate verifying that the specific timeslice, version and/or granular }}
\gT{108.00}{570.51}{\gS{R}{10}{c000000}{portion of data retrieved subsequently is the same as was originally cited. }}
\gT{90.00}{583.01}{\gS{R}{10}{c000000}{8)}\gS{RB}{10}{c000000}{ }\gS{RB}{10}{c000000}{Interoperability and Flexibility:}\gS{R}{10}{c000000}{ Data citation methods should be sufficiently flexible to accommodate the }}
\gT{108.00}{595.01}{\gS{R}{10}{c000000}{variant practices among communities, but should not differ so much that they compromise interoperability }}
\gT{108.00}{607.01}{\gS{R}{10}{c000000}{of data citation practices across communities. }}
\gT{72.00}{631.01}{\gS{RI}{10}{c000000}{Note:}\gS{R}{10}{c000000}{ Reproduced (including references therein) from M. Martone, Ed. “Data Citation Synthesis Group: Joint }}
\gT{72.00}{643.01}{\gS{R}{10}{c000000}{Declaration of Data Citation Principles}\gS{RI}{10}{c000000}{.}\gS{R}{10}{c000000}{” San Diego, CA, USA: FORCE11, 2014. [Online]. Available: }}
\gT{72.00}{655.01}{\gS{R}{10}{c1155CC}{https://www.force11.org/group/joint-declaration-data-citation-principles-final}}
\gT{72.00}{679.01}{\gS{R}{10}{c000000}{When providing reference details as much information as available should be provided (italics denote that metadata }}
\gT{72.00}{691.01}{\gS{R}{10}{c000000}{would benefit from the details, although it is not required for reference listing):}}
\gB{90.00}{715.56}
\gT{108.00}{715.56}{\gS{RB}{10}{c000000}{Published mixed data/software package:}\gS{R}{10}{c000000}{ author(s), name, unique identifier, location/repository, release }}
\end{guidepage}
\begin{guidepage}% page 11
\gT{532.30}{43.47}{\gS{R}{8}{c000000}{11 }}
\gT{108.00}{80.83}{\gS{R}{10}{c000000}{date,}\gS{RI}{10}{c000000}{ license, description, keywords. }}
\gT{72.00}{104.83}{\gS{R}{10}{c000000}{The essential components of a citation to a dataset are the following: Author names of each individual or }}
\gT{72.00}{116.83}{\gS{R}{10}{c000000}{organizational entity responsible for the creation of the dataset; the date published or disseminated [year and full }}
\gT{72.00}{128.83}{\gS{R}{10}{c000000}{date (if available)]; the complete title of the dataset, including the edition or version number, if applicable; publisher }}
\gT{72.00}{140.83}{\gS{R}{10}{c000000}{and/or distributor; and electronic location or identifier (URL or DOI if applicable). Append the date retrieved if the }}
\gT{72.00}{152.83}{\gS{R}{10}{c000000}{title and locator are not specific to the exact instance of the data you used.  }}
\gT{93.75}{175.28}{\gS{RI}{10}{c000000}{Basic Format:  }}
\gB{85.50}{189.71}
\gT{94.50}{189.71}{\gS{R}{10}{c000000}{Author, Date, Year. “Title of Dataset,” distributed by Publisher/Distributor, http://url.com (or if DOI is used, }}
\gT{94.50}{201.71}{\gS{R}{10}{c000000}{end with a period) }}
\gT{93.75}{217.16}{\gS{RI}{10}{c000000}{Example:  }}
\gB{85.50}{231.59}
\gE{94.50}{231.59}{447.87}{0.00}{12.00}{data:ansolabehere}{}
\gT{72.00}{276.01}{\gS{RB}{12}{c333333}{Online Dataset Reference Using a DOI }}
\gT{93.75}{292.09}{\gS{RI}{10}{c000000}{Basic Format: }}
\gB{86.25}{306.52}
\gT{93.75}{306.52}{\gS{R}{10}{c000000}{ Author, Date, “Title of Dataset,” Source, doi: }\gS{RI}{10}{c000000}{xxx}\gS{R}{10}{c000000}{. }}
\gT{93.75}{321.97}{\gS{RI}{10}{c000000}{Example: }}
\gB{86.25}{336.40}
\gE{93.75}{336.40}{448.25}{0.00}{12.00}{data:hhsdoi}{}
\gT{72.00}{389.17}{\gS{RB}{12}{c333333}{Online Dataset Reference Using a DOI Resolver }}
\gT{93.75}{405.25}{\gS{RI}{10}{c000000}{Basic Format: }}
\gB{86.25}{419.68}
\gT{93.75}{419.68}{\gS{R}{10}{c000000}{Author, Date, “Title of Dataset,” Source, doi: URL. }}
\gT{93.75}{435.13}{\gS{RI}{10}{c000000}{Example: }}
\gB{86.25}{449.56}
\gE{93.75}{449.56}{448.25}{0.00}{12.00}{data:hhsdoi}{}
\gT{72.00}{514.33}{\gS{RB}{12}{c333333}{Online Dataset Reference Using a Website Address }}
\gT{93.75}{530.41}{\gS{RI}{10}{c000000}{Basic Format: }}
\gB{86.25}{544.84}
\gT{93.75}{544.84}{\gS{R}{10}{c000000}{Author, Date, “Title of Dataset,” Source. [Online]. Available: http://www.url.com }}
\gT{93.75}{560.29}{\gS{RI}{10}{c000000}{Example: }}
\gB{86.25}{574.72}
\gE{93.75}{574.72}{448.25}{0.00}{12.00}{data:hhsurl}{}
\gT{72.00}{650.96}{\gS{R}{14}{c000000}{F. Handbooks  }}
\gT{93.75}{664.27}{\gS{RI}{10}{c000000}{Basic Format:  }}
\gB{85.50}{678.70}
\gT{94.50}{678.70}{\gS{RI}{10}{c000000}{Name of Manual/Handbook}\gS{R}{10}{c000000}{, }\gS{RI}{10}{c000000}{x }\gS{R}{10}{c000000}{ed., Abbrev. Name of Co., City of Co., Abbrev. State, year, pp. }\gS{RI}{10}{c000000}{xxx-xxx. }}
\gT{93.75}{694.15}{\gS{RI}{10}{c000000}{Examples:  }}
\gB{85.50}{708.58}
\gEC{0}{713.00}{94.50}{708.58}{447.70}{0.00}{13.80}{manual:transmission}{}
\end{guidepage}
\begin{guidepage}% page 12
\gT{532.01}{43.47}{\gS{R}{8}{c000000}{12 }}
\gEC{70.83}{792}{94.50}{67.03}{447.70}{0.00}{13.80}{manual:transmission}{}
\gB{85.50}{93.38}
\gE{94.50}{93.38}{445.50}{0.00}{13.80}{manual:motorola}{}
\gB{85.50}{105.93}
\gE{94.50}{105.93}{447.52}{0.00}{12.00}{manual:rca}{}
\gT{72.00}{158.17}{\gS{R}{14}{c000000}{G. Lectures  }}
\gT{72.00}{174.80}{\gS{RB}{12}{c333333}{Lecture Notes }}
\gT{93.75}{190.88}{\gS{RI}{10}{c000000}{Basic Format: }}
\gB{86.25}{205.31}
\gT{93.75}{205.31}{\gS{R}{10}{c000000}{J. K. Author. (Year). Title of lecture [Type of Medium]. Available: URL}}
\gT{93.75}{230.09}{\gS{RI}{10}{c000000}{Examples: }}
\gB{86.25}{244.52}
\gE{93.75}{244.52}{446.25}{0.00}{13.80}{lecture:barney}{}
\gB{86.25}{257.07}
\gE{93.75}{257.07}{448.43}{0.00}{12.00}{lecture:yardish}{}
\gT{72.00}{297.83}{\gS{RB}{12}{c333333}{Lecture Online }}
\gT{93.75}{313.92}{\gS{RI}{10}{c000000}{Basic Format: }}
\gB{86.25}{328.35}
\gT{93.75}{328.35}{\gS{R}{10}{c000000}{University name. (year). Title of lecture. [Type of Medium]. Available: URL}}
\gT{93.75}{343.80}{\gS{RI}{10}{c000000}{Example: }}
\gB{86.25}{358.23}
\gE{93.75}{358.23}{448.64}{0.00}{12.00}{lecture:argosy}{}
\gT{72.00}{410.47}{\gS{R}{14}{c000000}{H. }\gS{R}{14}{c000000}{Legal Citations }}
\gT{72.00}{435.74}{\gS{RI}{10}{c000000}{Note: }\gS{R}{10}{c000000}{In text, italicize the names of court cases. }}
\gT{72.00}{456.69}{\gS{RB}{10}{c000000}{U.S. Supreme Court Decisions }}
\gB{90.00}{470.19}
\gE{108.00}{470.19}{432.00}{0.00}{13.80}{legal:olmstead}{}
\gB{90.00}{482.74}
\gE{108.00}{482.74}{432.00}{0.00}{13.80}{legal:olmsteaddissent}{}
\gT{72.00}{502.74}{\gS{RB}{10}{c000000}{Lower-Court Decisions}}
\gB{90.00}{516.24}
\gE{108.00}{516.24}{434.30}{0.00}{12.00}{legal:councilman2003}{}
\gB{90.00}{540.79}
\gE{108.00}{540.79}{432.00}{0.00}{12.00}{legal:councilman2004}{}
\gT{72.00}{574.79}{\gS{RB}{10}{c000000}{U.S. Laws }}
\gB{90.00}{587.33}
\gE{108.00}{587.33}{432.00}{0.00}{13.80}{legal:uscode}{}
\gB{90.00}{599.89}
\gE{108.00}{599.89}{434.14}{0.00}{12.00}{legal:patriot}{}
\gB{90.00}{624.44}
\gE{108.00}{624.44}{433.85}{0.00}{12.00}{legal:omnibus}{}
\gT{72.00}{688.68}{\gS{R}{14}{c000000}{I. }\gS{R}{14}{c000000}{Manuals }}
\gT{93.75}{705.39}{\gS{RI}{10}{c000000}{Basic Format: }}
\end{guidepage}
\begin{guidepage}% page 13
\gT{532.01}{43.47}{\gS{R}{8}{c000000}{13 }}
\gB{86.25}{81.38}
\gT{93.75}{81.38}{\gS{R}{10}{c000000}{J. K. Author (or Abbrev. Name of Co., City of Co. Abbrev. State, Country). }\gS{RI}{10}{c000000}{Name of Manual/Handbook}\gS{R}{10}{c000000}{, }\gS{RI}{10}{c000000}{x}\gS{R}{10}{c000000}{ ed. }}
\gT{93.75}{93.38}{\gS{R}{10}{c000000}{(year). Accessed: Date. [Online]. Available: http://www.url.com}}
\gT{93.75}{108.83}{\gS{RI}{10}{c000000}{Examples: }}
\gB{86.25}{123.26}
\gE{93.75}{123.26}{448.67}{0.00}{12.00}{manual:breimann}{}
\gB{86.25}{147.81}
\gE{93.75}{147.81}{448.56}{0.00}{12.00}{manual:kuhn}{}
\gB{86.25}{172.36}
\gE{93.75}{172.36}{448.25}{0.00}{12.00}{manual:antcom}{}
\gT{72.00}{213.13}{\gS{RB}{12}{c333333}{Manual (Print) }}
\gT{93.75}{229.21}{\gS{RI}{10}{c000000}{Basic Format: }}
\gB{86.25}{243.64}
\gT{93.75}{243.64}{\gS{RI}{10}{c000000}{Name of Manual/Handbook}\gS{R}{10}{c000000}{, }\gS{RI}{10}{c000000}{x}\gS{R}{10}{c000000}{ ed., Abbrev. Name of Co., City of Co., Abbrev. State, Country, year, pp. }}
\gT{93.75}{254.09}{\gS{RI}{10}{c000000}{xxx–xxx. }}
\gT{93.75}{268.42}{\gS{RI}{10}{c000000}{Examples:}}
\gB{86.25}{282.85}
\gE{93.75}{282.85}{448.45}{0.00}{12.00}{manual:transmission}{}
\gB{86.25}{307.40}
\gE{93.75}{307.40}{446.25}{0.00}{13.80}{manual:motorola}{}
\gB{86.25}{319.95}
\gE{93.75}{319.95}{448.27}{0.00}{12.00}{manual:rca}{}
\gB{86.25}{344.50}
\gE{93.75}{344.50}{448.20}{0.00}{12.00}{manual:microsoft}{}
\gT{72.00}{385.26}{\gS{RB}{12}{c333333}{Manual (Online) }}
\gT{93.75}{401.35}{\gS{RI}{10}{c000000}{Basic Format: }}
\gB{86.25}{415.78}
\gT{93.75}{415.78}{\gS{R}{10}{c000000}{J. K. Author (or Abbrev. Name of Co., City of Co. Abbrev. State, Country). }\gS{RI}{10}{c000000}{Name of Manual/Handbook}\gS{R}{10}{c000000}{, }\gS{RI}{10}{c000000}{x}\gS{R}{10}{c000000}{ ed. }}
\gT{93.75}{427.78}{\gS{R}{10}{c000000}{(year). Accessed: Date. [Online]. Available: http://www.url.com}}
\gT{93.75}{443.23}{\gS{RI}{10}{c000000}{Examples: }}
\gB{86.25}{457.66}
\gE{93.75}{457.66}{448.67}{0.00}{12.00}{manual:breimann}{}
\gB{86.25}{482.21}
\gE{93.75}{482.21}{446.25}{0.00}{13.80}{manual:kuhn}{}
\gT{72.00}{522.45}{\gS{R}{14}{c000000}{J. News Articles }}
\gT{72.00}{546.68}{\gS{RB}{12}{c434343}{News Article (Print) }}
\gT{93.75}{561.56}{\gS{RI}{10}{c000000}{Basic Format: }}
\gB{85.50}{575.99}
\gT{103.50}{575.99}{\gS{R}{10}{c000000}{J. K. Author, “Title of the article,” }\gS{RI}{10}{c000000}{Title of the News Source, }\gS{R}{10}{c000000}{Month, Day, Year.}}
\gT{93.75}{591.44}{\gS{RI}{10}{c000000}{Examples: }}
\gB{90.00}{605.87}
\gE{108.00}{605.87}{432.00}{0.00}{13.80}{news:markoff}{}
\gB{90.00}{619.41}
\gE{108.00}{619.41}{432.00}{0.00}{13.80}{news:johnson}{}
\gT{72.00}{644.78}{\gS{RB}{12}{c434343}{News Article (Online) }}
\gT{93.75}{659.65}{\gS{RI}{10}{c000000}{Basic Format: }}
\gB{85.50}{674.08}
\gT{103.50}{674.08}{\gS{R}{10}{c000000}{J. K. Author, “Title of the article,” }\gS{RI}{10}{c000000}{Title of the News Source, }\gS{R}{10}{c000000}{Month, Day, Year. [Online]. Available: }}
\gT{103.50}{686.08}{\gS{R}{10}{c000000}{http://www.url.com}}
\gT{93.75}{701.53}{\gS{RI}{10}{c000000}{Examples: }}
\end{guidepage}
\begin{guidepage}% page 14
\gT{532.01}{43.47}{\gS{R}{8}{c000000}{14 }}
\gB{90.00}{81.38}
\gE{108.00}{81.38}{432.00}{0.00}{13.00}{news:clark}{}
\gB{90.00}{107.87}
\gE{108.00}{107.87}{432.00}{0.00}{13.00}{news:knight}{}
\gT{72.00}{175.01}{\gS{R}{14}{c000000}{K. }\gS{R}{14}{c000000}{Online Video (e.g., YouTube) }}
\gT{72.00}{198.28}{\gS{RI}{10}{c000000}{Basic Format:}}
\gB{90.00}{210.83}
\gT{108.00}{210.83}{\gS{R}{10}{c000000}{Video Owner/Creator, Location (if available). }\gS{RI}{10}{c000000}{Title of Video: In Initial Caps.}\gS{R}{10}{c000000}{ (Release date). Accessed: }}
\gT{108.00}{222.83}{\gS{R}{10}{c000000}{Month Day, Year. [Online Video]. Available: http://URL.onlinevideo.org}}
\gT{108.00}{233.28}{\gS{RI}{10}{c000000}{Examples:}}
\gB{90.00}{244.71}
\gE{108.00}{244.71}{434.33}{0.00}{12.00}{video:mtaonline}{}
\gB{90.00}{269.25}
\gE{108.00}{269.25}{433.97}{0.00}{12.00}{video:doane}{}
\gT{72.00}{321.50}{\gS{R}{14}{c000000}{L. }\gS{R}{14}{c000000}{Patent}}
\gT{93.75}{339.76}{\gS{R}{10}{c000000}{Retain or request the day of the month when referencing a patent. }}
\gT{93.75}{354.76}{\gS{RB}{10}{c000000}{NOTE:}\gS{R}{10}{c000000}{ Use “issued date” if several dates are given. }}
\gT{93.75}{377.21}{\gS{RI}{10}{c000000}{Basic Format: }}
\gB{86.25}{391.64}
\gT{93.75}{391.64}{\gS{R}{10}{c000000}{J. K. Author, “Title of patent,” U.S. Patent }\gS{RI}{10}{c000000}{x xxx xxx}\gS{R}{10}{c000000}{, Abbrev. Month, day, year. }}
\gB{86.25}{404.19}
\gT{93.75}{404.19}{\gS{R}{10}{c000000}{J. K. Author, “Title of patent,” Country Patent }\gS{RI}{10}{c000000}{xxx}\gS{R}{10}{c000000}{, Abbrev. Month, day, year. }}
\gT{93.75}{419.64}{\gS{RI}{10}{c000000}{Examples: }}
\gB{86.25}{434.07}
\gE{93.75}{434.07}{446.25}{0.00}{13.80}{patent:wilkinson}{}
\gB{86.25}{446.62}
\gE{93.75}{446.62}{448.30}{0.00}{12.00}{patent:mei}{}
\gB{86.25}{471.17}
\gE{93.75}{471.17}{446.25}{0.00}{13.80}{patent:black}{}
\gB{86.25}{483.72}
\gE{93.75}{483.72}{448.29}{0.00}{12.00}{patent:voinigescu}{}
\gB{86.25}{508.27}
\gE{93.75}{508.27}{448.57}{0.00}{12.00}{patent:klionovski}{}
\gT{72.00}{546.04}{\gS{RB}{12}{c333333}{Patent Online }}
\gT{93.75}{562.12}{\gS{RI}{10}{c000000}{Basic Format: }}
\gB{86.25}{576.55}
\gT{93.75}{576.55}{\gS{R}{10}{c000000}{Name of the invention, by inventor’s name. (year, month day). Patent Number [Type of medium]. Available: }}
\gT{93.75}{588.55}{\gS{R}{10}{c000000}{site/path/file }}
\gT{93.75}{604.00}{\gS{RI}{10}{c000000}{Examples: }}
\gB{86.25}{618.43}
\gE{93.75}{618.43}{448.48}{0.00}{12.00}{patent:brooks}{}
\gB{86.25}{642.98}
\gE{93.75}{642.98}{448.14}{0.00}{12.00}{patent:williams}{}
\gT{72.00}{695.22}{\gS{R}{14}{c000000}{M. }\gS{R}{14}{c000000}{Periodicals }}
\end{guidepage}
\begin{guidepage}% page 15
\gT{532.01}{43.47}{\gS{R}{8}{c000000}{15 }}
\gT{72.00}{82.60}{\gS{RB}{12}{c333333}{Periodicals}}
\gT{72.00}{100.74}{\gS{R}{10}{c000000}{Refer to the section “Common Abbreviations of Words in References” for non-IEEE titles. Note that periodical titles }}
\gT{72.00}{112.23}{\gS{R}{10}{c000000}{of only one word should not be abbreviated, but fully spelled out, e.g., }\gS{RI}{10}{c000000}{Science }\gS{R}{10}{c000000}{and }\gS{RI}{10}{c000000}{Nature. }}
\gT{72.00}{135.23}{\gS{R}{10}{c000000}{Prior to 1988, the volume number of IEEE Transactions/Journals carried the acronym of the journal. For example, an }}
\gT{72.00}{147.23}{\gS{R}{10}{c000000}{issue of the IEEE T}\gS{R}{8}{c000000}{RANSACTIONS ON}\gS{R}{10}{c000000}{ A}\gS{R}{8}{c000000}{UTOMATIC}\gS{R}{10}{c000000}{ C}\gS{R}{8}{c000000}{ONTROL}\gS{R}{10}{c000000}{ would read: }\gS{RI}{10}{c000000}{IEEE Trans. Autom. Control, }\gS{R}{10}{c000000}{vol. AC-26, }}
\gT{72.00}{159.23}{\gS{R}{10}{c000000}{no. 1, pp. 1–34, Jan. 1981. When referencing IEEE Transactions, both the issue number and month are included }}
\gT{72.00}{171.23}{\gS{R}{10}{c000000}{upon verification of frequency and starting month. DOIs are included, when provided by the author. }}
\gT{72.00}{195.23}{\gS{RB}{10}{c000000}{NOTE: }\gS{R}{10}{c000000}{The only exception to this rule is P}\gS{R}{8}{c000000}{ROCEEDINGS OF THE}\gS{R}{10}{c000000}{ I}\gS{R}{8}{c000000}{EEE}\gS{R}{10}{c000000}{, which never carried an acronym on the }}
\gT{72.00}{207.23}{\gS{R}{10}{c000000}{masthead. }}
\gT{72.00}{229.68}{\gS{RI}{10}{c000000}{Basic Format: }}
\gB{90.00}{241.11}
\gT{99.00}{241.11}{\gS{R}{10}{c000000}{J. K. Author, “Name of paper,” }\gS{RI}{10}{c000000}{Abbrev. Title of Periodical}\gS{R}{10}{c000000}{, vol. }\gS{RI}{10}{c000000}{x, }\gS{R}{10}{c000000}{no. }\gS{RI}{10}{c000000}{x, }\gS{R}{10}{c000000}{pp}\gS{RI}{10}{c000000}{. xxx-xxx, }\gS{R}{10}{c000000}{Abbrev. Month, year.  }}
\gB{90.00}{253.66}
\gT{99.00}{253.66}{\gS{R}{10}{c000000}{J. K. Author, “Name of paper,” }\gS{RI}{10}{c000000}{Abbrev. Title of Periodical}\gS{R}{10}{c000000}{, vol. }\gS{RI}{10}{c000000}{x, }\gS{R}{10}{c000000}{no. }\gS{RI}{10}{c000000}{x, }\gS{R}{10}{c000000}{pp}\gS{RI}{10}{c000000}{. xxx-xxx, }\gS{R}{10}{c000000}{Abbrev. Month, year, doi: }}
\gT{99.00}{265.66}{\gS{RI}{10}{c000000}{xxx}\gS{R}{10}{c000000}{. }}
\gT{72.00}{288.11}{\gS{RI}{10}{c000000}{Examples: }}
\gB{90.00}{299.54}
\gE{99.00}{299.54}{443.14}{0.00}{12.00}{per:chiampi}{}
\gB{90.00}{336.09}
\gE{99.00}{336.09}{442.76}{0.00}{12.00}{per:ito}{}
\gB{90.00}{360.64}
\gE{99.00}{360.64}{442.83}{0.00}{12.00}{per:fardel}{}
\gB{90.00}{385.19}
\gE{99.00}{385.19}{443.49}{0.00}{12.00}{per:zhang}{}
\gB{90.00}{409.74}
\gE{99.00}{409.74}{442.83}{0.00}{12.00}{per:preu}{}
\gB{90.00}{434.29}
\gE{99.00}{434.29}{443.07}{0.00}{12.00}{per:azodolmolky}{}
\gB{90.00}{458.84}
\gE{99.00}{458.84}{443.42}{0.00}{12.00}{per:eriksson}{}
\gB{90.00}{483.39}
\gE{99.00}{483.39}{441.00}{0.00}{13.80}{per:aronowitz}{}
\gB{90.00}{495.94}
\gE{99.00}{495.94}{443.16}{0.00}{12.00}{per:lavrova}{}
\gB{90.00}{532.49}
\gE{99.00}{532.49}{443.11}{0.00}{12.00}{per:wigner}{}
\gB{90.00}{557.04}
\gE{99.00}{557.04}{443.42}{0.00}{10.00}{per:miller}{ \textit{*** Always use this }\textit{style when the paper has been accepted or scheduled for a future publication but has not been published as }\textit{early access, i.e., do not use “to appear in.”***  }}
\gB{90.00}{588.25}
\gE{99.00}{588.25}{443.47}{0.00}{11.67}{per:vatta}{ \textit{***Always state if the article is early access and include its online version }\textit{date and DOI. The DOI is essential as it will not change, though the date may.*** }}
\gB{90.00}{634.13}
\gE{99.00}{634.13}{442.86}{0.00}{11.00}{per:kim}{ \textit{*** Always use this style when the }\textit{paper has not yet been accepted or scheduled for publication, i.e., do not use “to appear in.”*** }}
\gB{90.00}{656.01}
\gE{99.00}{656.01}{443.42}{0.00}{12.00}{per:rafferty}{}
\gB{90.00}{680.56}
\gE{99.00}{680.56}{443.36}{0.00}{12.00}{per:kopyt}{}
\gB{90.00}{705.11}
\gEC{0}{709.00}{99.00}{705.11}{443.23}{0.00}{13.80}{per:brunschwiler}{}
\end{guidepage}
\begin{guidepage}% page 16
\gT{532.01}{43.47}{\gS{R}{8}{c000000}{16 }}
\gEC{70.83}{792}{99.00}{67.03}{443.23}{0.00}{13.80}{per:brunschwiler}{}
\gB{90.00}{105.38}
\gE{99.00}{105.38}{443.37}{0.00}{10.67}{per:wu}{ \textit{***Authors may refer to this journal as Proc. IEE, but the abbreviation must be as listed above. Proc. IEE is }\textit{printed in the U.K. and must not be confused with the Proc. IEEE.*** }}
\gB{90.00}{148.59}
\gE{99.00}{148.59}{441.00}{0.00}{13.80}{per:specialissue}{}
\gT{72.00}{172.09}{\gS{RBI}{10}{c000000}{Article Referred to in the Same Issue:  }}
\gB{90.00}{184.64}
\gE{108.00}{184.64}{433.76}{0.00}{12.00}{per:aslip}{}
\gT{72.00}{208.64}{\gS{RB}{10}{c000000}{NOTE: }\gS{R}{10}{c000000}{Handle it exactly as any other reference with no difference.  }}
\gT{72.00}{240.05}{\gS{RB}{12}{c333333}{Periodical With Article ID }}
\gT{72.00}{258.19}{\gS{RB}{10}{c000000}{NOTE: }\gS{R}{10}{c000000}{Be aware that an article ID may erroneously be placed as page number in source files. If a page number }}
\gT{72.00}{269.69}{\gS{R}{10}{c000000}{looks “wrong” [e.g., there is only one long number, the range is not sequential, or there is a volume number with }}
\gT{72.00}{281.19}{\gS{R}{10}{c000000}{page range of 1–XX (with XX being a number)], this may be an indication that this is  an article ID and the author }}
\gT{72.00}{292.69}{\gS{R}{10}{c000000}{should be queried. }}
\gT{93.75}{318.63}{\gS{RI}{10}{c000000}{Basic Format: }}
\gB{86.25}{333.06}
\gT{93.75}{333.06}{\gS{R}{10}{c000000}{J. K. Author, “Name of paper,” }\gS{RI}{10}{c000000}{Abbrev. Title of Periodical}\gS{R}{10}{c000000}{, vol.}\gS{RI}{10}{c000000}{ x, }\gS{R}{10}{c000000}{no.}\gS{RI}{10}{c000000}{ x, }\gS{R}{10}{c000000}{Abbrev. month, year, Art. no. }\gS{RI}{10}{c000000}{xxx}\gS{R}{10}{c000000}{. }}
\gT{93.75}{348.51}{\gS{RI}{10}{c000000}{Examples: }}
\gB{86.25}{362.94}
\gE{93.75}{362.94}{448.74}{0.00}{12.00}{per:zhang}{}
\gB{86.25}{387.49}
\gE{93.75}{387.49}{448.05}{0.00}{12.00}{per:fardel}{}
\gB{86.25}{412.04}
\gE{93.75}{412.04}{448.08}{0.00}{12.00}{per:preu}{}
\gB{86.25}{436.59}
\gE{93.75}{436.59}{448.21}{0.00}{12.00}{per:kuang}{}
\gB{86.25}{461.14}
\gE{93.75}{461.14}{448.34}{0.00}{12.00}{per:yang}{}
\gT{72.00}{501.91}{\gS{RB}{12}{c333333}{Periodical With DOI }}
\gT{93.75}{517.99}{\gS{RI}{10}{c000000}{Basic Format: }}
\gB{86.25}{532.42}
\gT{93.75}{532.42}{\gS{R}{10}{c000000}{J. K. Author, “Name of paper,” }\gS{RI}{10}{c000000}{Abbrev. Title of Periodical}\gS{R}{10}{c000000}{, vol.}\gS{RI}{10}{c000000}{ x, }\gS{R}{10}{c000000}{no.}\gS{RI}{10}{c000000}{ x, }\gS{R}{10}{c000000}{pp}\gS{RI}{10}{c000000}{. xxx–xxx, }\gS{R}{10}{c000000}{Abbrev. month, year, doi: }}
\gT{93.75}{544.42}{\gS{RI}{10}{c000000}{xxx}\gS{R}{10}{c000000}{. }}
\gT{93.75}{559.87}{\gS{RI}{10}{c000000}{Example: }}
\gB{86.25}{574.30}
\gE{93.75}{574.30}{448.39}{0.00}{12.00}{per:chiampidoi}{}
\gT{72.00}{622.07}{\gS{RB}{12}{c333333}{Periodical in Early Access }}
\gT{72.00}{639.70}{\gS{RB}{10}{c000000}{NOTE:}\gS{R}{10}{c000000}{ Once an article is in the form of advanced online access at the publisher, cite this version, and not the arXiv }}
\gT{72.00}{651.70}{\gS{R}{10}{c000000}{version. The online access is the final version of record. Include the title of the journal, the date of record, and the }}
\gT{72.00}{663.70}{\gS{R}{10}{c000000}{doi. The doi is essential in the reference information, as this will not change. }}
\gT{93.75}{679.15}{\gS{RI}{10}{c000000}{Basic Format: }}
\gB{86.25}{693.58}
\gT{93.75}{693.58}{\gS{R}{10}{c000000}{J. K. Author, “Name of paper,” }\gS{RI}{10}{c000000}{Abbrev. Title of Periodical.,}\gS{R}{10}{c000000}{ Abbrev. month, year, doi: xxx. }}
\gT{93.75}{709.03}{\gS{RI}{10}{c000000}{Example: }}
\end{guidepage}
\begin{guidepage}% page 17
\gT{532.01}{43.47}{\gS{R}{8}{c000000}{17 }}
\gB{86.25}{81.38}
\gE{93.75}{81.38}{446.25}{0.00}{12.00}{per:vatta}{}
\gT{72.00}{129.15}{\gS{RB}{12}{c333333}{Periodical in Other Language }}
\gT{93.75}{145.23}{\gS{RI}{10}{c000000}{Basic Format: }}
\gB{86.25}{159.66}
\gT{93.75}{159.66}{\gS{R}{10}{c000000}{J. K. Author, “Name of paper,” (in Language), }\gS{RI}{10}{c000000}{Abbrev. Title of Periodical}\gS{R}{10}{c000000}{, vol.}\gS{RI}{10}{c000000}{ x, }\gS{R}{10}{c000000}{no.}\gS{RI}{10}{c000000}{ x, }\gS{R}{10}{c000000}{Abbrev. month, year, }}
\gT{93.75}{171.66}{\gS{R}{10}{c000000}{Art. no. }\gS{RI}{10}{c000000}{xxx}\gS{R}{10}{c000000}{. }}
\gT{93.75}{187.11}{\gS{RI}{10}{c000000}{Examples: }}
\gB{86.25}{201.54}
\gE{93.75}{201.54}{448.36}{0.00}{12.00}{per:wigner}{}
\gB{86.25}{226.09}
\gE{93.75}{226.09}{448.69}{0.00}{12.00}{per:lavrova}{}
\gT{72.00}{278.86}{\gS{RB}{12}{c333333}{Periodicals Online }}
\gT{93.75}{294.94}{\gS{RI}{10}{c000000}{Basic Format:  }}
\gB{85.50}{306.37}
\gT{93.00}{306.37}{\gS{R}{10}{c000000}{J. K. Author, “Name of paper,” }\gS{RI}{10}{c000000}{Abbrev. Title of Periodical}\gS{R}{10}{c000000}{, vol. x, no. x, pp. xxx–xxx, Abbrev. Month, year. }}
\gT{93.00}{318.37}{\gS{R}{10}{c000000}{Accessed: Month, Day, Year, doi: 10.1109.XXX.123456. [Online]. Available: site/path/file  }}
\gT{94.50}{333.82}{\gS{RI}{10}{c000000}{Examples:  }}
\gB{85.50}{348.25}
\gE{93.00}{348.25}{449.47}{0.00}{12.00}{per:risk}{}
\gB{85.50}{384.80}
\gE{94.50}{384.80}{447.86}{0.00}{12.00}{per:kopyt}{}
\gT{72.00}{427.56}{\gS{RB}{12}{c333333}{Virtual Journal }}
\gT{93.75}{443.65}{\gS{RI}{10}{c000000}{Basic Format: }}
\gB{86.25}{458.08}
\gT{93.75}{458.08}{\gS{R}{10}{c000000}{Name(s) of Ed(s)., “Title of Issue,” in }\gS{RI}{10}{c000000}{Title of Journal}\gS{R}{10}{c000000}{, Abbrev. month year. [Online]. Available: URL }}
\gT{93.75}{473.53}{\gS{RI}{10}{c222222}{Examples:}}
\gB{86.25}{487.96}
\gE{93.75}{487.96}{448.12}{0.00}{12.00}{per:smith}{}
\gB{86.25}{524.51}
\gE{93.75}{524.51}{448.32}{0.00}{12.00}{per:robertson}{}
\gB{86.25}{549.06}
\gE{93.75}{549.06}{448.36}{0.00}{12.00}{per:rfid}{}
\gT{72.00}{613.30}{\gS{R}{14}{c000000}{N. }\gS{R}{14}{c000000}{Reports  }}
\gT{72.00}{629.57}{\gS{R}{10}{c000000}{The general form for citing technical reports is to place the name and location of the company or institution after the }}
\gT{72.00}{641.57}{\gS{R}{10}{c000000}{author and title and to give the report number and date at the end of the reference.  }}
\gT{93.75}{661.35}{\gS{RI}{10}{c000000}{Basic Format:  }}
\gB{85.50}{675.78}
\gT{94.50}{675.78}{\gS{R}{10}{c000000}{J. K. Author, “Title of report,” Abbrev. Name of Co., City of Co., Abbrev. State, Country, Rep. }\gS{RI}{10}{c000000}{xxx}\gS{R}{10}{c000000}{, year.  }}
\gT{93.75}{691.23}{\gS{RI}{10}{c000000}{Examples:  }}
\gB{85.50}{705.66}
\gEC{0}{710.00}{94.50}{705.66}{447.95}{0.00}{13.80}{report:reber}{}
\end{guidepage}
\begin{guidepage}% page 18
\gL{c1155CC}{1.00}{72.00}{589.50}{365.00}{589.50}
\gT{532.01}{43.47}{\gS{R}{8}{c000000}{18 }}
\gEC{70.83}{792}{94.50}{67.03}{447.95}{0.00}{13.80}{report:reber}{}
\gB{85.50}{93.38}
\gE{94.50}{93.38}{447.67}{0.00}{12.00}{report:davis}{}
\gB{85.50}{117.93}
\gE{94.50}{117.93}{447.97}{0.00}{12.00}{report:haskell}{}
\gB{85.50}{142.48}
\gE{94.50}{142.48}{447.95}{0.00}{12.00}{report:brusberg}{}
\gB{85.50}{179.03}
\gE{94.50}{179.03}{447.54}{0.00}{12.00}{report:diament}{}
\gT{72.00}{216.80}{\gS{RB}{12}{c333333}{Report Online }}
\gT{85.50}{234.43}{\gS{R}{10}{c000000}{Ensure a year is included and add the URL to the end of the reference. }}
\gT{93.75}{254.88}{\gS{RI}{10}{c000000}{Basic Format: }}
\gB{86.25}{269.31}
\gT{93.75}{269.31}{\gS{R}{10}{c000000}{J. K. Author, “Title of report,” Company, City, State, Country, Rep. no., (optional: vol./issue), Date. Accessed: }}
\gT{93.75}{281.31}{\gS{R}{10}{c000000}{Date. [Online]. Available: site/path/file }}
\gT{93.75}{296.76}{\gS{RI}{10}{c000000}{Examples: }}
\gB{86.25}{311.19}
\gE{93.75}{311.19}{448.28}{0.00}{12.00}{report:hijmans}{}
\gB{86.25}{335.74}
\gE{93.75}{335.74}{448.29}{0.00}{12.00}{report:teralyzer}{}
\gB{86.25}{360.29}
\gE{93.75}{360.29}{448.35}{0.00}{12.00}{report:zhao}{}
\gB{86.25}{384.84}
\gE{93.75}{384.84}{448.58}{0.00}{12.00}{report:bom}{}
\gB{86.25}{409.39}
\gE{93.75}{409.39}{448.71}{0.00}{12.00}{report:geobasis}{}
\gB{86.25}{445.94}
\gE{93.75}{445.94}{448.02}{0.00}{12.00}{report:kagaku}{}
\gB{86.25}{470.49}
\gE{93.75}{470.49}{448.31}{0.00}{12.00}{report:apple}{}
\gB{86.25}{495.04}
\gE{93.75}{495.04}{448.48}{0.00}{12.00}{report:urazhdin}{}
\gT{72.00}{547.28}{\gS{R}{14}{c000000}{O. }\gS{R}{14}{c000000}{Software }}
\gT{72.00}{563.55}{\gS{RB}{10}{c000000}{FORCE11 Software Citation Principles }}
\gT{72.00}{576.50}{\gS{R}{10}{c000000}{IEEE }}
\gT{104.40}{576.50}{\gS{R}{10}{c000000}{follows }}
\gT{144.91}{576.50}{\gS{R}{10}{c000000}{the FORCE11 Software Citation Principles. These principles are as follows (per }}
\gT{72.00}{588.50}{\gS{R}{10}{c1155CC}{http://www.force11.org/sites/default/files/software-citation-principles.pdf}\gS{R}{10}{c000000}{): }}
\gT{90.00}{609.38}{\gS{S}{10}{c000000}{1)}\gS{RB}{10}{c000000}{ }\gS{RB}{10}{c000000}{Importance: }\gS{R}{10}{c000000}{Software should be considered a legitimate and citable product of research. Software citations }}
\gT{108.00}{621.37}{\gS{R}{10}{c000000}{should be accorded the same importance in the scholarly record as citations of other research products, such }}
\gT{108.00}{633.37}{\gS{R}{10}{c000000}{as publications and data; they should be included in the metadata of the citing work, for example in the }}
\gT{108.00}{645.37}{\gS{R}{10}{c000000}{reference list of a journal article, and should not be omitted or separated. Software should be cited on the }}
\gT{108.00}{657.37}{\gS{R}{10}{c000000}{same basis as any other research product such as a paper or a book, that is, authors should cite the }}
\gT{108.00}{669.37}{\gS{R}{10}{c000000}{appropriate set of software products just as they cite the appropriate set of papers. }}
\gT{90.00}{682.25}{\gS{S}{10}{c000000}{2)}\gS{RB}{10}{c000000}{ }\gS{RB}{10}{c000000}{Credit and Attribution: }\gS{R}{10}{c000000}{Software citations should facilitate giving scholarly credit and normative and }}
\gT{108.00}{694.25}{\gS{R}{10}{c000000}{legal attribution to all contributors to the software, recognizing that a single style or mechanism of }}
\gT{108.00}{706.25}{\gS{R}{10}{c000000}{attribution may not be applicable to all software. }}
\end{guidepage}
\begin{guidepage}% page 19
\gT{532.01}{43.47}{\gS{R}{8}{c000000}{19 }}
\gT{90.00}{81.71}{\gS{S}{10}{c000000}{3)}\gS{RB}{10}{c000000}{ }\gS{RB}{10}{c000000}{Unique Identification: }\gS{R}{10}{c000000}{A software citation should include a method for identification that is machine }}
\gT{108.00}{93.71}{\gS{R}{10}{c000000}{actionable, globally unique, interoperable, and recognized by at least a community of the corresponding }}
\gT{108.00}{105.71}{\gS{R}{10}{c000000}{domain experts, and preferably by general public researchers. }}
\gT{90.00}{118.58}{\gS{S}{10}{c000000}{4)}\gS{RB}{10}{c000000}{ }\gS{RB}{10}{c000000}{Persistence:}\gS{R}{10}{c000000}{ Unique identifiers and metadata describing the software and its disposition should }}
\gT{108.00}{130.58}{\gS{R}{10}{c000000}{persist—even beyond the lifespan of the software they describe. }}
\gT{90.00}{143.46}{\gS{S}{10}{c000000}{5)}\gS{RB}{10}{c000000}{ }\gS{RB}{10}{c000000}{Accessibility:}\gS{R}{10}{c000000}{ Software citations should facilitate access to the software itself and to its associated }}
\gT{108.00}{155.46}{\gS{R}{10}{c000000}{metadata, documentation, data, and other materials necessary for both humans and machines to make }}
\gT{108.00}{167.46}{\gS{R}{10}{c000000}{informed use of the referenced software. }}
\gT{90.00}{180.34}{\gS{S}{10}{c000000}{6)}\gS{RB}{10}{c000000}{ }\gS{RB}{10}{c000000}{Specificity: }\gS{R}{10}{c000000}{Software citations should facilitate identification of, and access to, the specific version of }}
\gT{108.00}{192.34}{\gS{R}{10}{c000000}{software that was used. Software identification should be as specific as necessary, such as using version }}
\gT{108.00}{204.34}{\gS{R}{10}{c000000}{numbers, revision numbers, or variants such as platforms. }}
\gT{72.00}{224.34}{\gS{R}{10}{c000000}{***When providing reference details as much information as noted below should be provided (italics denote that }}
\gT{72.00}{236.34}{\gS{R}{10}{c000000}{metadata would benefit from the details, although it is not required for reference listing): }}
\gT{90.00}{257.22}{\gS{S}{10}{c000000}{1)}\gS{RB}{10}{c000000}{ }\gS{RB}{10}{c000000}{Use software for papers:}\gS{R}{10}{c000000}{ author(s), software name, location/repository, version number, release date, }}
\gT{108.00}{269.21}{\gS{R}{10}{c000000}{unique identifier. }}
\gT{90.00}{282.09}{\gS{S}{10}{c000000}{2)}\gS{RB}{10}{c000000}{ }\gS{RB}{10}{c000000}{Use software in/with new software: }\gS{R}{10}{c000000}{author(s), software name, location/repository, version number, release }}
\gT{108.00}{294.09}{\gS{R}{10}{c000000}{date, unique identifier, }\gS{RI}{10}{c000000}{software licensure, description. }}
\gT{90.00}{306.09}{\gS{R}{10}{c000000}{3)}\gS{R}{10}{c000000}{ }\gS{R}{10}{c000000}{Contribute to software: author(s), software name, location/repository, version number, release date, unique }}
\gT{108.00}{318.09}{\gS{R}{10}{c000000}{identifier, }\gS{RI}{10}{c000000}{contributor role, software license, description. }}
\gT{90.00}{330.97}{\gS{S}{10}{c000000}{4)}\gS{RB}{10}{c000000}{ }\gS{RB}{10}{c000000}{Determine use/citations of software: }\gS{R}{10}{c000000}{software name, unique identifier, indexed citations. }}
\gT{90.00}{343.85}{\gS{S}{10}{c000000}{5)}\gS{RB}{10}{c000000}{ }\gS{RB}{10}{c000000}{Get credit for software development: }\gS{R}{10}{c000000}{author(s), software name, location/repository, release date, unique }}
\gT{108.00}{355.85}{\gS{R}{10}{c000000}{identifier, }\gS{RI}{10}{c000000}{contributor role, indexed citations. }}
\gT{90.00}{368.72}{\gS{S}{10}{c000000}{6)}\gS{RB}{10}{c000000}{ }\gS{RB}{10}{c000000}{“Reproduce analysis”:}\gS{R}{10}{c000000}{ software name, location/repository, version number, release date, unique identifier, }}
\gT{108.00}{379.17}{\gS{RI}{10}{c000000}{software license, description. }}
\gT{90.00}{390.93}{\gS{S}{10}{c000000}{7)}\gS{RB}{10}{c000000}{ }\gS{RB}{10}{c000000}{Find software to implement task:}\gS{R}{10}{c000000}{ author(s), software name, location/repository, location/repository, }}
\gT{108.00}{402.93}{\gS{R}{10}{c000000}{unique identifier, indexed citations, }\gS{RI}{10}{c000000}{software license, description, keywords. }}
\gT{90.00}{415.81}{\gS{S}{10}{c000000}{8)}\gS{RB}{10}{c000000}{ }\gS{RB}{10}{c000000}{Publish software paper:}\gS{R}{10}{c000000}{ author(s), software name, location/repository, version number, release date, }}
\gT{108.00}{427.81}{\gS{R}{10}{c000000}{unique identifier. }}
\gT{90.00}{440.68}{\gS{S}{10}{c000000}{9)}\gS{RB}{10}{c000000}{ }\gS{RB}{10}{c000000}{Publish papers that cite software: }\gS{R}{10}{c000000}{author(s), software name, location/repository, version number, release }}
\gT{108.00}{452.69}{\gS{R}{10}{c000000}{date, unique identifier, indexed citations. }}
\gT{90.00}{465.56}{\gS{S}{10}{c000000}{10)}\gS{RB}{10}{c000000}{ Build catalog of software: }\gS{R}{10}{c000000}{author(s), software name, location/repository, version number, release date, }}
\gT{108.00}{477.56}{\gS{R}{10}{c000000}{unique identifier, indexed citations, }\gS{RI}{10}{c000000}{software license, description, keywords. }}
\gT{90.00}{490.44}{\gS{S}{10}{c000000}{11)}\gS{RB}{10}{c000000}{ Build software catalog/registry: }\gS{R}{10}{c000000}{author(s), software name, location/repository, unique identifier, }}
\gT{108.00}{500.89}{\gS{RI}{10}{c000000}{description, keywords. }}
\gT{90.00}{512.64}{\gS{S}{10}{c000000}{12)}\gS{RB}{10}{c000000}{ Show scientific impact of holdings: }\gS{R}{10}{c000000}{software name, unique identifier, indexed citations. }}
\gT{90.00}{525.52}{\gS{S}{10}{c000000}{13)}\gS{RB}{10}{c000000}{ Show how funded software has been used: }\gS{R}{10}{c000000}{software name, unique identifier, indexed citations. }}
\gT{90.00}{538.40}{\gS{S}{10}{c000000}{14)}\gS{RB}{10}{c000000}{ Evaluate contributions of researcher: }\gS{R}{10}{c000000}{author(s), unique identifier, release date, indexed citations, }}
\gT{108.00}{548.85}{\gS{RI}{10}{c000000}{contributor role.  }}
\gT{90.00}{560.61}{\gS{S}{10}{c000000}{15)}\gS{RB}{10}{c000000}{ Store software entry:}\gS{R}{10}{c000000}{ author(s), software name, unique identifier, location/repository, release date, indexed }}
\gT{108.00}{572.61}{\gS{R}{10}{c000000}{citations, }\gS{RI}{10}{c000000}{description, keywords. }}
\gT{90.00}{585.48}{\gS{S}{10}{c000000}{16)}\gS{RB}{10}{c000000}{ Published mixed data/software package: }\gS{R}{10}{c000000}{author(s), software name, unique identifier, location/repository, }}
\gT{108.00}{597.48}{\gS{R}{10}{c000000}{release date, }\gS{RI}{10}{c000000}{software license, description, keywords. }}
\gT{72.00}{617.48}{\gS{RI}{10}{c000000}{Note:}\gS{R}{10}{c000000}{ Reproduced (including references therein) from A. M. Smith, D. S. Katz, K. E. Niemeyer, and the }}
\gT{72.00}{629.48}{\gS{R}{10}{c000000}{FORCE11 Software Citation Working Group. “Software citation principles,” }\gS{RI}{10}{c000000}{PeerJ Comp. Sci.,}\gS{R}{10}{c000000}{ vol. 2, 2016, }}
\gT{72.00}{641.48}{\gS{R}{10}{c000000}{Article no. e86, doi: 10.7717/peerj-cs.86 }}
\gT{94.50}{663.93}{\gS{RI}{10}{c000000}{Basic Format: }}
\gB{90.00}{675.36}
\gT{99.00}{675.36}{\gS{R}{10}{c000000}{J. K. Author. }\gS{RI}{10}{c000000}{Title of Software}\gS{R}{10}{c000000}{. (Date). Repository or Archive. (version or year). Publisher Name. Accessed: }}
\gT{99.00}{687.36}{\gS{R}{10}{c000000}{Date (when applicable). [Type of Medium]. Global Persistent Identifier. Available: site/path/file }}
\gT{94.50}{705.81}{\gS{RI}{10}{c000000}{Examples: }}
\end{guidepage}
\begin{guidepage}% page 20
\gT{532.01}{43.47}{\gS{R}{8}{c000000}{20 }}
\gB{90.00}{81.38}
\gE{99.00}{81.38}{443.34}{0.00}{12.00}{sw:antcom}{}
\gB{90.00}{105.93}
\gE{99.00}{105.93}{443.43}{0.00}{12.00}{sw:arning}{}
\gB{90.00}{130.48}
\gE{99.00}{130.48}{441.00}{0.00}{13.80}{sw:msdn}{}
\gT{72.00}{170.72}{\gS{R}{14}{c000000}{P. }\gS{R}{14}{c000000}{Standards  }}
\gT{93.75}{185.44}{\gS{RI}{10}{c000000}{Basic Format:  }}
\gB{85.50}{199.87}
\gT{94.50}{199.87}{\gS{RI}{10}{c000000}{Title of Standard}\gS{R}{10}{c000000}{, Standard number, Corporate author, location, date.  }}
\gB{85.50}{212.42}
\gT{94.50}{212.42}{\gS{RI}{10}{c000000}{Title of Standard}\gS{R}{10}{c000000}{, Standard number, date.  }}
\gT{93.75}{227.87}{\gS{RI}{10}{c000000}{Examples:  }}
\gB{85.50}{242.30}
\gE{94.50}{242.30}{447.63}{0.00}{12.00}{std:itu}{}
\gB{85.50}{278.85}
\gE{94.50}{278.85}{445.50}{0.00}{13.80}{std:ieee308}{}
\gB{85.50}{291.40}
\gE{94.50}{291.40}{445.50}{0.00}{13.80}{std:ansi}{}
\gT{72.00}{317.16}{\gS{RB}{12}{c333333}{Standard Online }}
\gT{93.75}{333.25}{\gS{RI}{10}{c000000}{Basic Format: }}
\gB{86.25}{347.68}
\gT{93.75}{347.68}{\gS{RI}{10}{c000000}{Title of Standard}\gS{R}{10}{c000000}{, Standard number, Corporate author, location, date. [Online]. Available: http://www.url.com }}
\gB{86.25}{360.23}
\gT{93.75}{360.23}{\gS{RI}{10}{c000000}{Title of Standard}\gS{R}{10}{c000000}{, Standard number, date. [Online]. Available: http://www.url.com }}
\gT{93.75}{375.68}{\gS{RI}{10}{c000000}{Example: }}
\gB{86.25}{390.11}
\gE{93.75}{390.11}{448.52}{0.00}{12.00}{std:nerc}{}
\gT{72.00}{441.85}{\gS{R}{14}{c000000}{Q. }\gS{R}{14}{c000000}{Theses and Dissertations  }}
\gT{93.75}{456.56}{\gS{RI}{10}{c000000}{Basic Format:  }}
\gB{85.50}{471.00}
\gT{94.50}{471.00}{\gS{R}{10}{c000000}{J. K. Author, “Title of thesis,” M.S. thesis, Abbrev. Dept., Abbrev. Univ., City of Univ., Abbrev. State, year.  }}
\gB{85.50}{485.14}
\gT{94.50}{485.14}{\gS{R}{10}{c000000}{J. K. Author, “Title of dissertation,” Ph.D. dissertation, Abbrev. Dept., Abbrev. Univ., City of Univ., Abbrev. }}
\gT{94.50}{497.14}{\gS{R}{10}{c000000}{State, year.  }}
\gT{93.75}{512.59}{\gS{RI}{10}{c000000}{Examples:  }}
\gB{85.50}{527.02}
\gE{94.50}{527.02}{447.82}{0.00}{12.00}{thesis:williams}{}
\gB{85.50}{551.57}
\gE{94.50}{551.57}{447.80}{0.00}{12.00}{thesis:kawasaki}{}
\gB{85.50}{576.12}
\gE{94.50}{576.12}{447.75}{0.00}{11.50}{thesis:amer}{ \textit{*** The state abbreviation is }\textit{omitted if the name of the university includes the state name, i.e., “Univ. California, Berkeley.”*** }}
\gB{85.50}{610.00}
\gE{94.50}{610.00}{445.50}{0.00}{13.80}{thesis:becle}{}
\gT{72.00}{635.27}{\gS{RB}{12}{c333333}{Thesis Online }}
\gT{93.75}{651.35}{\gS{RI}{10}{c000000}{Basic Format: }}
\gB{86.25}{665.78}
\gT{93.75}{665.78}{\gS{R}{10}{c000000}{J. K. Author, “Title of thesis,” M.S. thesis, Abbrev. Dept., Abbrev. Univ., City of Univ., Abbrev. State, Country, }}
\gT{93.75}{677.78}{\gS{R}{10}{c000000}{year. [Online]. Available: http://www.url.com }}
\gB{86.25}{690.33}
\gT{93.75}{690.33}{\gS{R}{10}{c000000}{J. K. Author, “Title of dissertation,” Ph.D. dissertation, Abbrev. Dept., Abbrev. Univ., City of Univ., Abbrev. }}
\gT{93.75}{702.33}{\gS{R}{10}{c000000}{State, Country, year. [Online]. Available: http://www.url.com }}
\gB{86.25}{714.88}
\gT{93.75}{714.88}{\gS{R}{10}{c000000}{J. K. Author, “Title of theis,” Ph.D. thesis, Abbrev. Dept., Abbrev. Univ., City of Univ., Abbrev. State, Country, }}
\end{guidepage}
\begin{guidepage}% page 21
\gT{532.01}{43.47}{\gS{R}{8}{c000000}{21 }}
\gT{93.75}{80.83}{\gS{R}{10}{c000000}{year. [Online]. Available: http://www.url.com }}
\gT{94.50}{96.28}{\gS{RI}{10}{c000000}{Examples: }}
\gB{86.25}{110.71}
\gE{93.75}{110.71}{446.25}{0.00}{12.00}{thesis:jensen}{}
\gB{86.25}{135.26}
\gE{93.75}{135.26}{446.25}{0.00}{12.00}{thesis:bordaouin}{}
\gB{86.25}{159.81}
\gE{93.75}{159.81}{446.25}{0.00}{12.00}{thesis:schwartz}{}
\gT{72.00}{210.81}{\gS{RI}{10}{c000000}{Note:}\gS{RBI}{10}{c000000}{ }\gS{R}{10}{c000000}{Defer to the use of “thesis” and “dissertation” as provided by the authors; these are different depending on the }}
\gT{72.00}{222.81}{\gS{R}{10}{c000000}{degree and the academic institution and should not be changed based on degree level. }}
\gT{72.00}{263.05}{\gS{R}{14}{c000000}{R. }\gS{R}{14}{c000000}{U.S. Government Documents}}
\gT{93.75}{279.77}{\gS{RI}{10}{c000000}{Basic Format:  }}
\gB{85.50}{294.20}
\gT{94.50}{294.20}{\gS{R}{10}{c000000}{Legislative body. Number of Congress, Session. (year, month day). }\gS{RI}{10}{c000000}{Number of bill or resolution, Title}\gS{R}{10}{c000000}{. [Type of }}
\gT{94.50}{306.20}{\gS{R}{10}{c000000}{medium]. Available: site/path/file  }}
\gT{93.75}{321.65}{\gS{RI}{10}{c000000}{Example:  }}
\gB{85.50}{336.08}
\gE{94.50}{336.08}{447.75}{0.00}{12.00}{gov:house}{}
\gT{72.00}{373.84}{\gS{RB}{12}{c333333}{Government Online }}
\gT{93.75}{389.93}{\gS{RI}{10}{c000000}{Basic Format: }}
\gB{86.25}{404.36}
\gT{93.75}{404.36}{\gS{R}{10}{c000000}{Legislative body. Number of Congress, Session. (year, month day). }\gS{RI}{10}{c000000}{Number of bill or resolution}\gS{R}{10}{c000000}{, }\gS{RI}{10}{c000000}{Title}\gS{R}{10}{c000000}{. [Type of }}
\gT{93.75}{416.36}{\gS{R}{10}{c000000}{medium]. Available: site/path/file }}
\gT{93.75}{431.81}{\gS{RI}{10}{c000000}{Example: }}
\gB{86.25}{446.24}
\gE{93.75}{446.24}{448.50}{0.00}{12.00}{gov:house}{}
\gT{72.00}{497.98}{\gS{R}{14}{c000000}{S. }\gS{R}{14}{c000000}{Unpublished }}
\gT{86.25}{514.24}{\gS{R}{10}{c000000}{The two most common types of unpublished references are private communications and papers still in }}
\gT{72.00}{526.24}{\gS{R}{10}{c000000}{preparation. }}
\gT{93.75}{541.69}{\gS{RI}{10}{c000000}{Basic Format: }}
\gB{86.25}{556.12}
\gT{93.75}{556.12}{\gS{R}{10}{c000000}{J. K. Author, private communication, Abbrev. Month, year. }}
\gB{86.25}{568.67}
\gT{93.75}{568.67}{\gS{R}{10}{c000000}{J. K. Author, “Title of paper,” unpublished. }}
\gT{94.50}{584.13}{\gS{RI}{10}{c000000}{Examples: }}
\gB{86.25}{600.56}
\gE{93.75}{600.56}{446.25}{0.00}{13.80}{unpub:harrison}{}
\gB{86.25}{616.10}
\gE{93.75}{616.10}{446.25}{0.00}{13.80}{unpub:smith}{}
\gB{86.25}{628.65}
\gE{93.75}{628.65}{448.47}{0.00}{12.00}{unpub:brahms}{}
\gT{85.50}{666.09}{\gS{RB}{11}{c333333}{Preprint arXiv }}
\gT{94.50}{682.81}{\gS{RI}{10}{c000000}{Basic Format: }}
\gB{86.25}{694.35}
\gT{93.75}{694.35}{\gS{R}{10}{c000000}{J. K. Author, “Title of paper,” }\gS{R}{10}{c000000}{year, }\gS{RI}{10}{c000000}{arXiv number.}}
\gT{94.50}{709.80}{\gS{RI}{10}{c000000}{Examples: }}
\end{guidepage}
\begin{guidepage}% page 22
\gT{532.01}{43.47}{\gS{R}{8}{c000000}{22 }}
\gB{86.25}{81.38}
\gE{93.75}{81.38}{448.48}{0.00}{12.00}{unpub:urazhdin}{}
\gT{72.00}{132.62}{\gS{R}{14}{c000000}{T. }\gS{R}{14}{c000000}{Websites }}
\gT{72.00}{160.38}{\gS{R}{10}{c000000}{The most basic entry for a website consists of the author name(s), page title, website title, web address, and date }}
\gT{72.00}{172.38}{\gS{R}{10}{c000000}{accessed: }}
\gT{72.00}{196.38}{\gS{R}{10}{c000000}{First Name Initial(s) Last Name. “Page Title.” Website Title. Date Accessed. [Online]. Available: Web Address. }}
\gE{72.00}{220.38}{470.39}{0.00}{12.00}{web:smith}{}
\gT{72.00}{256.38}{\gS{R}{10}{c000000}{The first author’s name should be first initial(s) and then the last name and a period. Titles and affiliations associated }}
\gT{72.00}{268.38}{\gS{R}{10}{c000000}{with the author should be omitted. Do not use commas to precede a suffix, such as a roman numeral or Jr./Sr., after }}
\gT{72.00}{280.38}{\gS{R}{10}{c000000}{the author’s given name.  }}
\gT{72.00}{304.38}{\gS{R}{10}{c000000}{For a page with two or more authors, list them in the order as they appear on the website. Separate author names by }}
\gT{72.00}{316.38}{\gS{R}{10}{c000000}{a comma. }}
\gE{72.00}{340.38}{472.86}{0.00}{12.00}{web:smithdoe}{}
\gT{72.00}{376.38}{\gS{R}{10}{c000000}{Next, place the title of the online article or the title of the webpage itself. Then, the text “Accessed:” and the date on }}
\gT{72.00}{388.38}{\gS{R}{10}{c000000}{which you accessed the website (written in the format of “month day, year”). The text “[Online.] Available:” is then }}
\gT{72.00}{400.38}{\gS{R}{10}{c000000}{inserted with the complete URL after the colon.  }}
\gT{72.00}{424.38}{\gS{R}{10}{c000000}{For social media websites, please use the following style. }}
\gE{72.00}{448.38}{468.00}{0.00}{12.00}{web:costanzo}{}
\gT{196.84}{558.41}{\gS{RB}{14}{c0070C0}{III. N}\gS{RB}{9.8}{c0070C0}{OTES}\gS{RB}{14}{c0070C0}{ A}\gS{RB}{9.8}{c0070C0}{BOUT}\gS{RB}{14}{c0070C0}{ O}\gS{RB}{9.8}{c0070C0}{NLINE}\gS{RB}{14}{c0070C0}{ R}\gS{RB}{9.8}{c0070C0}{EFERENCES}}
\gT{72.00}{574.68}{\gS{R}{10}{c000000}{The guidelines for citing electronic information as offered here are in modified illustration of the adaptation by the }}
\gT{72.00}{586.68}{\gS{R}{10}{c000000}{International Standards Organization (ISO) documentation system and the American Psychological Association }}
\gT{72.00}{598.68}{\gS{R}{10}{c000000}{style. }}
\gT{86.25}{620.68}{\gS{R}{10}{c000000}{There are various options for including a URL in a reference. Follow the order of details and whether a period }}
\gT{72.00}{632.68}{\gS{R}{10}{c000000}{needs to be placed at the end of the reference. The style of the accessed date is “Accessed: Abbrev. month and day, }}
\gT{72.00}{644.68}{\gS{R}{10}{c000000}{year.” The placement of the accessed date within the reference should match how it is provided in the final }}
\gT{72.00}{656.68}{\gS{R}{10}{c000000}{author-submitted version. }}
\gT{86.25}{678.68}{\gS{R}{10}{c000000}{URLs are not hyperlinked in the proof.  }}
\gB{106.50}{701.23}
\gT{124.50}{701.23}{\gS{R}{10}{c000000}{Accessed date. DOI. [Online]. Available: URL (no period at end) }}
\gB{106.50}{713.78}
\gT{124.50}{713.78}{\gS{R}{10}{c000000}{Accessed date. [Online]. Available: URL (no period at end) }}
\end{guidepage}
\begin{guidepage}% page 23
\gF{cEFEFEF}{72.00}{482.17}{289.50}{497.17}
\gL{c000000}{1.00}{78.00}{494.50}{85.00}{494.50}
\gL{c000000}{1.00}{72.50}{482.00}{72.50}{708.00}
\gL{c000000}{1.00}{182.50}{497.00}{182.50}{708.00}
\gL{c000000}{1.00}{289.50}{482.00}{289.50}{708.00}
\gL{c000000}{1.00}{72.00}{482.50}{290.00}{482.50}
\gL{c000000}{1.00}{72.00}{497.50}{290.00}{497.50}
\gL{c000000}{1.00}{72.00}{512.50}{290.00}{512.50}
\gL{c000000}{1.00}{72.00}{527.50}{290.00}{527.50}
\gL{c000000}{1.00}{72.00}{542.50}{290.00}{542.50}
\gL{c000000}{1.00}{72.00}{557.50}{290.00}{557.50}
\gL{c000000}{1.00}{72.00}{572.50}{290.00}{572.50}
\gL{c000000}{1.00}{72.00}{587.50}{290.00}{587.50}
\gL{c000000}{1.00}{72.00}{602.50}{290.00}{602.50}
\gL{c000000}{1.00}{72.00}{617.50}{290.00}{617.50}
\gL{c000000}{1.00}{72.00}{632.50}{290.00}{632.50}
\gL{c000000}{1.00}{72.00}{647.50}{290.00}{647.50}
\gL{c000000}{1.00}{72.00}{662.50}{290.00}{662.50}
\gL{c000000}{1.00}{72.00}{677.50}{290.00}{677.50}
\gL{c000000}{1.00}{72.00}{692.50}{290.00}{692.50}
\gL{c000000}{1.00}{72.00}{707.50}{290.00}{707.50}
\gL{c000000}{1.00}{324.50}{470.00}{324.50}{711.00}
\gL{c000000}{1.00}{434.50}{470.00}{434.50}{711.00}
\gL{c000000}{1.00}{541.50}{470.00}{541.50}{711.00}
\gL{c000000}{1.00}{324.00}{470.50}{542.00}{470.50}
\gL{c000000}{1.00}{324.00}{485.50}{542.00}{485.50}
\gL{c000000}{1.00}{324.00}{500.50}{542.00}{500.50}
\gL{c000000}{1.00}{324.00}{515.50}{542.00}{515.50}
\gL{c000000}{1.00}{324.00}{530.50}{542.00}{530.50}
\gL{c000000}{1.00}{324.00}{545.50}{542.00}{545.50}
\gL{c000000}{1.00}{324.00}{560.50}{542.00}{560.50}
\gL{c000000}{1.00}{324.00}{575.50}{542.00}{575.50}
\gL{c000000}{1.00}{324.00}{590.50}{542.00}{590.50}
\gL{c000000}{1.00}{324.00}{605.50}{542.00}{605.50}
\gL{c000000}{1.00}{324.00}{620.50}{542.00}{620.50}
\gL{c000000}{1.00}{324.00}{635.50}{542.00}{635.50}
\gL{c000000}{1.00}{324.00}{650.50}{542.00}{650.50}
\gL{c000000}{1.00}{324.00}{665.50}{542.00}{665.50}
\gL{c000000}{1.00}{324.00}{680.50}{542.00}{680.50}
\gL{c000000}{1.00}{324.00}{695.50}{542.00}{695.50}
\gL{c000000}{1.00}{324.00}{710.50}{542.00}{710.50}
\gT{532.01}{43.47}{\gS{R}{8}{c000000}{23 }}
\gB{106.50}{81.38}
\gT{124.50}{81.38}{\gS{R}{10}{c000000}{Accessed date. DOI (add period at end). }}
\gB{106.50}{93.93}
\gT{124.50}{93.93}{\gS{R}{10}{c000000}{DOI. URL (no period at end) }}
\gB{106.50}{106.48}
\gT{124.50}{106.48}{\gS{R}{10}{c000000}{(No accessed date), DOI (add period at end unless URL). }}
\gB{106.50}{119.03}
\gT{124.50}{119.03}{\gS{R}{10}{c000000}{URL (no period at end)    }}
\gT{86.25}{139.48}{\gS{RI}{10}{c000000}{Guidelines for Breaking URLs:}}
\gB{104.25}{150.91}
\gT{122.25}{150.91}{\gS{R}{10}{c000000}{Break after slash, double slash, or period. }}
\gB{104.25}{163.46}
\gT{122.25}{163.46}{\gS{R}{10}{c000000}{Break “before” the hyphen that is part of an address, but do not break after; do not add hyphens or }}
\gT{122.25}{175.46}{\gS{R}{10}{c000000}{spaces; do not let addresses hyphenate. }}
\gB{104.25}{188.01}
\gT{122.25}{188.01}{\gS{R}{10}{c000000}{Break “before” a tilde (\textasciitilde{}), a hyphen, an underscore (\_), a question mark, or a percent (\%) symbol. }}
\gB{104.25}{200.56}
\gT{122.25}{200.56}{\gS{R}{10}{c000000}{Break before or after an equal sign or an ampersand (follow the same rule for the “at” (@) symbol).}}
\gT{187.63}{259.89}{\gS{RB}{14}{c0070C0}{IV. U}\gS{RB}{9.8}{c0070C0}{SEFUL}\gS{RB}{14}{c0070C0}{ A}\gS{RB}{9.8}{c0070C0}{BBREVIATIONS}\gS{RB}{14}{c0070C0}{ }\gS{RB}{9.8}{c0070C0}{IN}\gS{RB}{14}{c0070C0}{ R}\gS{RB}{9.8}{c0070C0}{EFERENCES}}
\gT{72.00}{287.66}{\gS{R}{10}{c000000}{Tips for abbreviating periodical titles in references: }}
\gB{90.00}{300.21}
\gT{108.00}{300.21}{\gS{R}{10}{c222222}{Words ending in "-ology" can be ended after the "-ol."  }}
\gT{108.00}{312.21}{\gS{R}{10}{c222222}{e.g., }\gS{RI}{10}{c222222}{Gastroenterology }\gS{R}{10}{c222222}{--> }\gS{RI}{10}{c222222}{Gastrooceanenterol}\gS{R}{10}{c222222}{.; }\gS{RI}{10}{c222222}{Endocrinology }\gS{R}{10}{c222222}{--> }\gS{RI}{10}{c222222}{Endocrinol.}}
\gB{90.00}{324.76}
\gT{108.00}{324.76}{\gS{R}{10}{c222222}{Words ending in "-graphy" can be ended after the "-gr." }}
\gT{108.00}{336.76}{\gS{R}{10}{c222222}{e.g., }\gS{RI}{10}{c222222}{Oceanography --> Oceanogr.; Crystallography --> Crystallgr.}}
\gB{90.00}{349.31}
\gT{108.00}{349.31}{\gS{R}{10}{c222222}{Compound words can be ended using the abbreviation of the last word. }}
\gT{108.00}{361.31}{\gS{R}{10}{c222222}{e.g., }\gS{RI}{10}{c222222}{Bioengineering --> Bioeng.; Nanobioscience --> Nanobiosci. }}
\gB{90.00}{373.86}
\gT{108.00}{373.86}{\gS{R}{10}{c222222}{Some abbreviations may apply to more than one word, e.g., "Mathematical" and "Mathematics" are both }}
\gT{108.00}{385.86}{\gS{R}{10}{c222222}{abbreviated as "Math.";  "Medical" and "Medicine" are both abbreviated as "Med." }}
\gB{90.00}{398.41}
\gT{108.00}{398.41}{\gS{R}{10}{c222222}{If an abbreviation is not included and cannot be abbreviated according to the general guidelines above, }}
\gT{108.00}{410.41}{\gS{R}{10}{c222222}{please spell out the word. }}
\gT{85.50}{451.74}{\gS{RB}{14}{c000000}{A.}\hspace*{4.40bp}\gS{R}{14}{c000000}{Common Abbreviations of Words in References }}
\gT{78.00}{493.63}{\gS{RB}{10}{c000000}{A }}
\gT{78.00}{508.63}{\gS{R}{10}{c000000}{Abstracts }}
\gT{188.25}{508.63}{\gS{R}{10}{c000000}{Abstr. }}
\gT{78.00}{523.63}{\gS{R}{10}{c000000}{Academy }}
\gT{188.25}{523.63}{\gS{R}{10}{c000000}{Acad. }}
\gT{78.00}{538.63}{\gS{R}{10}{c000000}{Accelerator }}
\gT{188.25}{538.63}{\gS{R}{10}{c000000}{Accel. }}
\gT{78.00}{553.63}{\gS{R}{10}{c000000}{Acoustics }}
\gT{188.25}{553.63}{\gS{R}{10}{c000000}{Acoust. }}
\gT{78.00}{568.63}{\gS{R}{10}{c000000}{Active }}
\gT{188.25}{568.63}{\gS{R}{10}{c000000}{Act. }}
\gT{78.00}{583.63}{\gS{R}{10}{c000000}{Administration }}
\gT{188.25}{583.63}{\gS{R}{10}{c000000}{Admin. }}
\gT{78.00}{598.63}{\gS{R}{10}{c000000}{Administrative }}
\gT{188.25}{598.63}{\gS{R}{10}{c000000}{Administ. }}
\gT{78.00}{613.63}{\gS{R}{10}{c000000}{Advanced }}
\gT{188.25}{613.63}{\gS{R}{10}{c000000}{Adv. }}
\gT{78.00}{628.63}{\gS{R}{10}{c000000}{Aeronautics }}
\gT{188.25}{628.63}{\gS{R}{10}{c000000}{Aeronaut. }}
\gT{78.00}{643.63}{\gS{R}{10}{c000000}{Aerospace }}
\gT{188.25}{643.63}{\gS{R}{10}{c000000}{Aerosp. }}
\gT{78.00}{658.63}{\gS{R}{10}{c000000}{Affective }}
\gT{188.25}{658.63}{\gS{R}{10}{c000000}{Affect. }}
\gT{78.00}{673.63}{\gS{R}{10}{c000000}{Africa, African }}
\gT{188.25}{673.63}{\gS{R}{10}{c000000}{Afr. }}
\gT{78.00}{688.63}{\gS{R}{10}{c000000}{Aircraft }}
\gT{188.25}{688.63}{\gS{R}{10}{c000000}{Aircr. }}
\gT{78.00}{703.63}{\gS{R}{10}{c000000}{Algebraic }}
\gT{188.25}{703.63}{\gS{R}{10}{c000000}{Algebr. }}
\gT{330.00}{482.14}{\gS{R}{10}{c000000}{American }}
\gT{440.25}{482.14}{\gS{R}{10}{c000000}{Amer. }}
\gT{330.00}{497.14}{\gS{R}{10}{c000000}{Analysis }}
\gT{440.25}{497.14}{\gS{R}{10}{c000000}{Anal. }}
\gT{330.00}{512.14}{\gS{R}{10}{c000000}{Annals }}
\gT{440.25}{512.14}{\gS{R}{10}{c000000}{Ann. }}
\gT{330.00}{527.14}{\gS{R}{10}{c000000}{Annual }}
\gT{440.25}{527.14}{\gS{R}{10}{c000000}{Annu. }}
\gT{330.00}{542.14}{\gS{R}{10}{c000000}{Apparatus }}
\gT{440.25}{542.14}{\gS{R}{10}{c000000}{App. }}
\gT{330.00}{557.14}{\gS{R}{10}{c000000}{Applications }}
\gT{440.25}{557.14}{\gS{R}{10}{c000000}{Appl. }}
\gT{330.00}{572.14}{\gS{R}{10}{c000000}{Applied }}
\gT{440.25}{572.14}{\gS{R}{10}{c000000}{Appl. }}
\gT{330.00}{587.14}{\gS{R}{10}{c000000}{Approximate }}
\gT{440.25}{587.14}{\gS{R}{10}{c000000}{Approx. }}
\gT{330.00}{602.14}{\gS{R}{10}{c000000}{Architecture }}
\gT{440.25}{602.14}{\gS{R}{10}{c000000}{Archit. }}
\gT{330.00}{617.14}{\gS{R}{10}{c000000}{Archive(s) }}
\gT{440.25}{617.14}{\gS{R}{10}{c000000}{Arch. }}
\gT{330.00}{632.14}{\gS{R}{10}{c000000}{Article }}
\gT{440.25}{632.14}{\gS{R}{10}{c000000}{Art. }}
\gT{330.00}{647.14}{\gS{R}{10}{c000000}{Artificial }}
\gT{440.25}{647.14}{\gS{R}{10}{c000000}{Artif. }}
\gT{330.00}{662.14}{\gS{R}{10}{c000000}{Assembly }}
\gT{440.25}{662.14}{\gS{R}{10}{c000000}{Assem. }}
\gT{330.00}{677.14}{\gS{R}{10}{c000000}{Association }}
\gT{440.25}{677.14}{\gS{R}{10}{c000000}{Assoc. }}
\gT{330.00}{692.14}{\gS{R}{10}{c000000}{Astronomy }}
\gT{440.25}{692.14}{\gS{R}{10}{c000000}{Astron. }}
\gT{330.00}{707.14}{\gS{R}{10}{c000000}{Astronautics }}
\gT{440.25}{707.14}{\gS{R}{10}{c000000}{Astronaut. }}
\end{guidepage}
\begin{guidepage}% page 24
\gF{cEFEFEF}{72.00}{222.00}{289.50}{237.00}
\gF{cEFEFEF}{72.00}{432.00}{289.50}{447.00}
\gF{cEFEFEF}{324.00}{207.00}{541.50}{222.00}
\gF{cEFEFEF}{324.00}{447.00}{541.50}{462.00}
\gL{c000000}{1.00}{78.00}{234.50}{85.00}{234.50}
\gL{c000000}{1.00}{78.00}{444.50}{85.00}{444.50}
\gL{c000000}{1.00}{330.00}{219.50}{337.00}{219.50}
\gL{c000000}{1.00}{330.00}{459.50}{337.00}{459.50}
\gL{c000000}{1.00}{72.50}{72.00}{72.50}{717.00}
\gL{c000000}{1.00}{182.50}{72.00}{182.50}{222.00}
\gL{c000000}{1.00}{182.50}{237.00}{182.50}{432.00}
\gL{c000000}{1.00}{182.50}{447.00}{182.50}{717.00}
\gL{c000000}{1.00}{289.50}{72.00}{289.50}{717.00}
\gL{c000000}{1.00}{72.00}{72.50}{290.00}{72.50}
\gL{c000000}{1.00}{72.00}{87.50}{290.00}{87.50}
\gL{c000000}{1.00}{72.00}{102.50}{290.00}{102.50}
\gL{c000000}{1.00}{72.00}{117.50}{290.00}{117.50}
\gL{c000000}{1.00}{72.00}{132.50}{290.00}{132.50}
\gL{c000000}{1.00}{72.00}{147.50}{290.00}{147.50}
\gL{c000000}{1.00}{72.00}{162.50}{290.00}{162.50}
\gL{c000000}{1.00}{72.00}{177.50}{290.00}{177.50}
\gL{c000000}{1.00}{72.00}{192.50}{290.00}{192.50}
\gL{c000000}{1.00}{72.00}{207.50}{290.00}{207.50}
\gL{c000000}{1.00}{72.00}{222.50}{290.00}{222.50}
\gL{c000000}{1.00}{72.00}{237.50}{290.00}{237.50}
\gL{c000000}{1.00}{72.00}{252.50}{290.00}{252.50}
\gL{c000000}{1.00}{72.00}{267.50}{290.00}{267.50}
\gL{c000000}{1.00}{72.00}{282.50}{290.00}{282.50}
\gL{c000000}{1.00}{72.00}{297.50}{290.00}{297.50}
\gL{c000000}{1.00}{72.00}{312.50}{290.00}{312.50}
\gL{c000000}{1.00}{72.00}{327.50}{290.00}{327.50}
\gL{c000000}{1.00}{72.00}{342.50}{290.00}{342.50}
\gL{c000000}{1.00}{72.00}{357.50}{290.00}{357.50}
\gL{c000000}{1.00}{72.00}{372.50}{290.00}{372.50}
\gL{c000000}{1.00}{72.00}{387.50}{290.00}{387.50}
\gL{c000000}{1.00}{72.00}{402.50}{290.00}{402.50}
\gL{c000000}{1.00}{72.00}{417.50}{290.00}{417.50}
\gL{c000000}{1.00}{72.00}{432.50}{290.00}{432.50}
\gL{c000000}{1.00}{72.00}{447.50}{290.00}{447.50}
\gL{c000000}{1.00}{72.00}{462.50}{290.00}{462.50}
\gL{c000000}{1.00}{72.00}{477.50}{290.00}{477.50}
\gL{c000000}{1.00}{72.00}{492.50}{290.00}{492.50}
\gL{c000000}{1.00}{72.00}{507.50}{290.00}{507.50}
\gL{c000000}{1.00}{72.00}{522.50}{290.00}{522.50}
\gL{c000000}{1.00}{72.00}{537.50}{290.00}{537.50}
\gL{c000000}{1.00}{72.00}{552.50}{290.00}{552.50}
\gL{c000000}{1.00}{72.00}{567.50}{290.00}{567.50}
\gL{c000000}{1.00}{72.00}{582.50}{290.00}{582.50}
\gL{c000000}{1.00}{72.00}{597.50}{290.00}{597.50}
\gL{c000000}{1.00}{72.00}{612.50}{290.00}{612.50}
\gL{c000000}{1.00}{72.00}{627.50}{290.00}{627.50}
\gL{c000000}{1.00}{72.00}{642.50}{290.00}{642.50}
\gL{c000000}{1.00}{72.00}{657.50}{290.00}{657.50}
\gL{c000000}{1.00}{72.00}{672.50}{290.00}{672.50}
\gL{c000000}{1.00}{72.00}{687.50}{290.00}{687.50}
\gL{c000000}{1.00}{72.00}{702.50}{290.00}{702.50}
\gL{c000000}{1.00}{72.00}{717.50}{290.00}{717.50}
\gL{c000000}{1.00}{324.50}{72.00}{324.50}{717.00}
\gL{c000000}{1.00}{434.50}{72.00}{434.50}{207.00}
\gL{c000000}{1.00}{434.50}{222.00}{434.50}{447.00}
\gL{c000000}{1.00}{434.50}{462.00}{434.50}{717.00}
\gL{c000000}{1.00}{541.50}{72.00}{541.50}{717.00}
\gL{c000000}{1.00}{324.00}{72.50}{542.00}{72.50}
\gL{c000000}{1.00}{324.00}{87.50}{542.00}{87.50}
\gL{c000000}{1.00}{324.00}{102.50}{542.00}{102.50}
\gL{c000000}{1.00}{324.00}{117.50}{542.00}{117.50}
\gL{c000000}{1.00}{324.00}{132.50}{542.00}{132.50}
\gL{c000000}{1.00}{324.00}{147.50}{542.00}{147.50}
\gL{c000000}{1.00}{324.00}{162.50}{542.00}{162.50}
\gL{c000000}{1.00}{324.00}{177.50}{542.00}{177.50}
\gL{c000000}{1.00}{324.00}{192.50}{542.00}{192.50}
\gL{c000000}{1.00}{324.00}{207.50}{542.00}{207.50}
\gL{c000000}{1.00}{324.00}{222.50}{542.00}{222.50}
\gL{c000000}{1.00}{324.00}{237.50}{542.00}{237.50}
\gL{c000000}{1.00}{324.00}{252.50}{542.00}{252.50}
\gL{c000000}{1.00}{324.00}{267.50}{542.00}{267.50}
\gL{c000000}{1.00}{324.00}{282.50}{542.00}{282.50}
\gL{c000000}{1.00}{324.00}{297.50}{542.00}{297.50}
\gL{c000000}{1.00}{324.00}{312.50}{542.00}{312.50}
\gL{c000000}{1.00}{324.00}{327.50}{542.00}{327.50}
\gL{c000000}{1.00}{324.00}{342.50}{542.00}{342.50}
\gL{c000000}{1.00}{324.00}{357.50}{542.00}{357.50}
\gL{c000000}{1.00}{324.00}{372.50}{542.00}{372.50}
\gL{c000000}{1.00}{324.00}{387.50}{542.00}{387.50}
\gL{c000000}{1.00}{324.00}{402.50}{542.00}{402.50}
\gL{c000000}{1.00}{324.00}{417.50}{542.00}{417.50}
\gL{c000000}{1.00}{324.00}{432.50}{542.00}{432.50}
\gL{c000000}{1.00}{324.00}{447.50}{542.00}{447.50}
\gL{c000000}{1.00}{324.00}{462.50}{542.00}{462.50}
\gL{c000000}{1.00}{324.00}{477.50}{542.00}{477.50}
\gL{c000000}{1.00}{324.00}{492.50}{542.00}{492.50}
\gL{c000000}{1.00}{324.00}{507.50}{542.00}{507.50}
\gL{c000000}{1.00}{324.00}{522.50}{542.00}{522.50}
\gL{c000000}{1.00}{324.00}{537.50}{542.00}{537.50}
\gL{c000000}{1.00}{324.00}{552.50}{542.00}{552.50}
\gL{c000000}{1.00}{324.00}{567.50}{542.00}{567.50}
\gL{c000000}{1.00}{324.00}{582.50}{542.00}{582.50}
\gL{c000000}{1.00}{324.00}{597.50}{542.00}{597.50}
\gL{c000000}{1.00}{324.00}{612.50}{542.00}{612.50}
\gL{c000000}{1.00}{324.00}{627.50}{542.00}{627.50}
\gL{c000000}{1.00}{324.00}{642.50}{542.00}{642.50}
\gL{c000000}{1.00}{324.00}{657.50}{542.00}{657.50}
\gL{c000000}{1.00}{324.00}{672.50}{542.00}{672.50}
\gL{c000000}{1.00}{324.00}{687.50}{542.00}{687.50}
\gL{c000000}{1.00}{324.00}{702.50}{542.00}{702.50}
\gL{c000000}{1.00}{324.00}{717.50}{542.00}{717.50}
\gT{532.01}{43.47}{\gS{R}{8}{c000000}{24 }}
\gT{78.00}{83.46}{\gS{R}{10}{c000000}{Astrophysics }}
\gT{188.25}{83.46}{\gS{R}{10}{c000000}{Astrophys. }}
\gT{78.00}{98.46}{\gS{R}{10}{c000000}{Atmosphere }}
\gT{188.25}{98.46}{\gS{R}{10}{c000000}{Atmos. }}
\gT{78.00}{113.46}{\gS{R}{10}{c000000}{Atomic, Atoms }}
\gT{188.25}{113.46}{\gS{R}{10}{c000000}{At. }}
\gT{78.00}{128.46}{\gS{R}{10}{c000000}{Australasian }}
\gT{188.25}{128.46}{\gS{R}{10}{c000000}{Australas. }}
\gT{78.00}{143.46}{\gS{R}{10}{c000000}{Australia }}
\gT{188.25}{143.46}{\gS{R}{10}{c000000}{Aust. }}
\gT{78.00}{158.46}{\gS{R}{10}{c000000}{Automatic }}
\gT{188.25}{158.46}{\gS{R}{10}{c000000}{Autom. }}
\gT{78.00}{173.46}{\gS{R}{10}{c000000}{Automation }}
\gT{188.25}{173.46}{\gS{R}{10}{c000000}{Automat. }}
\gT{78.00}{188.46}{\gS{R}{10}{c000000}{Automotive }}
\gT{188.25}{188.46}{\gS{R}{10}{c000000}{Automot. }}
\gT{78.00}{203.46}{\gS{R}{10}{c000000}{Autonomous }}
\gT{188.25}{203.46}{\gS{R}{10}{c000000}{Auton. }}
\gT{78.00}{233.46}{\gS{RB}{10}{c000000}{B }}
\gT{78.00}{248.46}{\gS{R}{10}{c000000}{Behavior(al) }}
\gT{188.25}{248.46}{\gS{R}{10}{c000000}{Behav. }}
\gT{78.00}{263.46}{\gS{R}{10}{c000000}{Belgian }}
\gT{188.25}{263.46}{\gS{R}{10}{c000000}{Belg. }}
\gT{78.00}{278.46}{\gS{R}{10}{c000000}{Biochemical }}
\gT{188.25}{278.46}{\gS{R}{10}{c000000}{Biochem. }}
\gT{78.00}{293.46}{\gS{R}{10}{c000000}{Bioinformatics }}
\gT{188.25}{293.46}{\gS{R}{10}{c000000}{Bioinf. }}
\gT{78.00}{308.46}{\gS{R}{10}{c000000}{Biology, Biological }}
\gT{188.25}{308.46}{\gS{R}{10}{c000000}{Biol. }}
\gT{78.00}{323.46}{\gS{R}{10}{c000000}{Biomedical }}
\gT{188.25}{323.46}{\gS{R}{10}{c000000}{Biomed. }}
\gT{78.00}{338.46}{\gS{R}{10}{c000000}{Biophysics }}
\gT{188.25}{338.46}{\gS{R}{10}{c000000}{Biophys. }}
\gT{78.00}{353.46}{\gS{R}{10}{c000000}{British }}
\gT{188.25}{353.46}{\gS{R}{10}{c000000}{Brit. }}
\gT{78.00}{368.46}{\gS{R}{10}{c000000}{Broadcasting }}
\gT{188.25}{368.46}{\gS{R}{10}{c000000}{Broadcast. }}
\gT{78.00}{383.46}{\gS{R}{10}{c000000}{Bulletin }}
\gT{188.25}{383.46}{\gS{R}{10}{c000000}{Bull. }}
\gT{78.00}{398.46}{\gS{R}{10}{c000000}{Bureau }}
\gT{188.25}{398.46}{\gS{R}{10}{c000000}{Bur. }}
\gT{78.00}{413.46}{\gS{R}{10}{c000000}{Business }}
\gT{188.25}{413.46}{\gS{R}{10}{c000000}{Bus. }}
\gT{78.00}{443.46}{\gS{RB}{10}{c000000}{C }}
\gT{78.00}{458.46}{\gS{R}{10}{c000000}{Canadian }}
\gT{188.25}{458.46}{\gS{R}{10}{c000000}{Can. }}
\gT{78.00}{473.46}{\gS{R}{10}{c000000}{Ceramic }}
\gT{188.25}{473.46}{\gS{R}{10}{c000000}{Ceram. }}
\gT{78.00}{488.46}{\gS{R}{10}{c000000}{Chemical }}
\gT{188.25}{488.46}{\gS{R}{10}{c000000}{Chem. }}
\gT{78.00}{503.46}{\gS{R}{10}{c000000}{Chinese }}
\gT{188.25}{503.46}{\gS{R}{10}{c000000}{Chin. }}
\gT{78.00}{518.46}{\gS{R}{10}{c000000}{Climatology }}
\gT{188.25}{518.46}{\gS{R}{10}{c000000}{Climatol. }}
\gT{78.00}{533.46}{\gS{R}{10}{c000000}{Clinical }}
\gT{188.25}{533.46}{\gS{R}{10}{c000000}{Clin. }}
\gT{78.00}{548.46}{\gS{R}{10}{c000000}{Cognitive }}
\gT{188.25}{548.46}{\gS{R}{10}{c000000}{Cogn. }}
\gT{78.00}{563.46}{\gS{R}{10}{c000000}{Colloquium }}
\gT{188.25}{563.46}{\gS{R}{10}{c000000}{Colloq. }}
\gT{78.00}{578.46}{\gS{R}{10}{c000000}{Communications }}
\gT{188.25}{578.46}{\gS{R}{10}{c000000}{Commun. }}
\gT{78.00}{593.46}{\gS{R}{10}{c000000}{Company }}
\gT{188.25}{593.46}{\gS{R}{10}{c000000}{Co. }}
\gT{78.00}{608.46}{\gS{R}{10}{c000000}{Compatibility }}
\gT{188.25}{608.46}{\gS{R}{10}{c000000}{Compat. }}
\gT{78.00}{623.46}{\gS{R}{10}{c000000}{Component(s) }}
\gT{188.25}{623.46}{\gS{R}{10}{c000000}{Compon. }}
\gT{78.00}{638.46}{\gS{R}{10}{c000000}{Computational }}
\gT{188.25}{638.46}{\gS{R}{10}{c000000}{Comput. }}
\gT{78.00}{653.46}{\gS{R}{10}{c000000}{Computer(s) }}
\gT{188.25}{653.46}{\gS{R}{10}{c000000}{Comput. }}
\gT{78.00}{668.46}{\gS{R}{10}{c000000}{Computing }}
\gT{188.25}{668.46}{\gS{R}{10}{c000000}{Comput. }}
\gT{78.00}{683.46}{\gS{R}{10}{c000000}{Condensed }}
\gT{188.25}{683.46}{\gS{R}{10}{c000000}{Condens. }}
\gT{78.00}{698.46}{\gS{R}{10}{c000000}{Conference }}
\gT{188.25}{698.46}{\gS{R}{10}{c000000}{Conf. }}
\gT{78.00}{713.46}{\gS{R}{10}{c000000}{Congress }}
\gT{188.25}{713.46}{\gS{R}{10}{c000000}{Congr. }}
\gT{330.00}{83.46}{\gS{R}{10}{c000000}{Consumer }}
\gT{440.25}{83.46}{\gS{R}{10}{c000000}{Consum. }}
\gT{330.00}{98.46}{\gS{R}{10}{c000000}{Conversion }}
\gT{440.25}{98.46}{\gS{R}{10}{c000000}{Convers. }}
\gT{330.00}{113.46}{\gS{R}{10}{c000000}{Convention }}
\gT{440.25}{113.46}{\gS{R}{10}{c000000}{Conv. }}
\gT{330.00}{128.46}{\gS{R}{10}{c000000}{Correspondence }}
\gT{440.25}{128.46}{\gS{R}{10}{c000000}{Corresp. }}
\gT{330.00}{143.46}{\gS{R}{10}{c000000}{Critical }}
\gT{440.25}{143.46}{\gS{R}{10}{c000000}{Crit. }}
\gT{330.00}{158.46}{\gS{R}{10}{c000000}{Crystal }}
\gT{440.25}{158.46}{\gS{R}{10}{c000000}{Cryst. }}
\gT{330.00}{173.46}{\gS{R}{10}{c000000}{Crystallography }}
\gT{440.25}{173.46}{\gS{R}{10}{c000000}{Crystallogr. }}
\gT{330.00}{188.46}{\gS{R}{10}{c000000}{Cybernetics }}
\gT{440.25}{188.46}{\gS{R}{10}{c000000}{Cybern. }}
\gT{330.00}{218.46}{\gS{RB}{10}{c000000}{D }}
\gT{330.00}{233.46}{\gS{R}{10}{c000000}{Decision }}
\gT{440.25}{233.46}{\gS{R}{10}{c000000}{Decis. }}
\gT{330.00}{248.46}{\gS{R}{10}{c000000}{Delivery }}
\gT{440.25}{248.46}{\gS{R}{10}{c000000}{Del. }}
\gT{330.00}{263.46}{\gS{R}{10}{c000000}{Department }}
\gT{440.25}{263.46}{\gS{R}{10}{c000000}{Dept. }}
\gT{330.00}{278.46}{\gS{R}{10}{c000000}{Design }}
\gT{440.25}{278.46}{\gS{R}{10}{c000000}{Des. }}
\gT{330.00}{293.46}{\gS{R}{10}{c000000}{Detector }}
\gT{440.25}{293.46}{\gS{R}{10}{c000000}{Detect. }}
\gT{330.00}{308.46}{\gS{R}{10}{c000000}{Development(al) }}
\gT{440.25}{308.46}{\gS{R}{10}{c000000}{Develop. }}
\gT{330.00}{323.46}{\gS{R}{10}{c000000}{Differential }}
\gT{440.25}{323.46}{\gS{R}{10}{c000000}{Differ. }}
\gT{330.00}{338.46}{\gS{R}{10}{c000000}{Digest }}
\gT{440.25}{338.46}{\gS{R}{10}{c000000}{Dig. }}
\gT{330.00}{353.46}{\gS{R}{10}{c000000}{Digital }}
\gT{440.25}{353.46}{\gS{R}{10}{c000000}{Digit. }}
\gT{330.00}{368.46}{\gS{R}{10}{c000000}{Disclosure }}
\gT{440.25}{368.46}{\gS{R}{10}{c000000}{Discl. }}
\gT{330.00}{383.46}{\gS{R}{10}{c000000}{Discussions }}
\gT{440.25}{383.46}{\gS{R}{10}{c000000}{Discuss. }}
\gT{330.00}{398.46}{\gS{R}{10}{c000000}{Dissertations }}
\gT{440.25}{398.46}{\gS{R}{10}{c000000}{Diss. }}
\gT{330.00}{413.46}{\gS{R}{10}{c000000}{Distributed }}
\gT{440.25}{413.46}{\gS{R}{10}{c000000}{Distrib. }}
\gT{330.00}{428.46}{\gS{R}{10}{c000000}{Dynamics }}
\gT{440.25}{428.46}{\gS{R}{10}{c000000}{Dyn. }}
\gT{330.00}{458.46}{\gS{RB}{10}{c000000}{E }}
\gT{330.00}{473.46}{\gS{R}{10}{c000000}{Earthquake }}
\gT{440.25}{473.46}{\gS{R}{10}{c000000}{Earthq. }}
\gT{330.00}{488.46}{\gS{R}{10}{c000000}{Economic(s) }}
\gT{440.25}{488.46}{\gS{R}{10}{c000000}{Econ. }}
\gT{330.00}{503.46}{\gS{R}{10}{c000000}{Edition }}
\gT{440.25}{503.46}{\gS{R}{10}{c000000}{Ed. }}
\gT{330.00}{518.46}{\gS{R}{10}{c000000}{Education }}
\gT{440.25}{518.46}{\gS{R}{10}{c000000}{Educ. }}
\gT{330.00}{533.46}{\gS{R}{10}{c000000}{Electrical }}
\gT{440.25}{533.46}{\gS{R}{10}{c000000}{Elect. }}
\gT{330.00}{548.46}{\gS{R}{10}{c000000}{Electrification }}
\gT{440.25}{548.46}{\gS{R}{10}{c000000}{Electrific. }}
\gT{330.00}{563.46}{\gS{R}{10}{c000000}{Electromagnetic }}
\gT{440.25}{563.46}{\gS{R}{10}{c000000}{Electromagn. }}
\gT{330.00}{578.46}{\gS{R}{10}{c000000}{Electroacoustic }}
\gT{440.25}{578.46}{\gS{R}{10}{c000000}{Electroacoust. }}
\gT{330.00}{593.46}{\gS{R}{10}{c000000}{Electronic }}
\gT{440.25}{593.46}{\gS{R}{10}{c000000}{Electron. }}
\gT{330.00}{608.46}{\gS{R}{10}{c000000}{Emerging }}
\gT{440.25}{608.46}{\gS{R}{10}{c000000}{Emerg. }}
\gT{330.00}{623.46}{\gS{R}{10}{c000000}{Engineering }}
\gT{440.25}{623.46}{\gS{R}{10}{c000000}{Eng. }}
\gT{330.00}{638.46}{\gS{R}{10}{c000000}{Environment }}
\gT{440.25}{638.46}{\gS{R}{10}{c000000}{Environ. }}
\gT{330.00}{653.46}{\gS{R}{10}{c000000}{Equations }}
\gT{440.25}{653.46}{\gS{R}{10}{c000000}{Equ. }}
\gT{330.00}{668.46}{\gS{R}{10}{c000000}{Equipment }}
\gT{440.25}{668.46}{\gS{R}{10}{c000000}{Equip. }}
\gT{330.00}{683.46}{\gS{R}{10}{c000000}{Ergonomics }}
\gT{440.25}{683.46}{\gS{R}{10}{c000000}{Ergonom. }}
\gT{330.00}{698.46}{\gS{R}{10}{c000000}{European }}
\gT{440.25}{698.46}{\gS{R}{10}{c000000}{Eur. }}
\gT{330.00}{713.46}{\gS{R}{10}{c000000}{Evaluation }}
\gT{440.25}{713.46}{\gS{R}{10}{c000000}{Eval. }}
\end{guidepage}
\begin{guidepage}% page 25
\gF{cEFEFEF}{72.00}{177.00}{289.50}{192.00}
\gF{cEFEFEF}{72.00}{312.00}{289.50}{327.00}
\gF{cEFEFEF}{72.00}{432.00}{289.50}{447.00}
\gF{cEFEFEF}{72.00}{552.00}{289.50}{567.00}
\gF{cEFEFEF}{324.00}{192.00}{541.50}{207.00}
\gF{cEFEFEF}{324.00}{252.00}{541.50}{267.00}
\gF{cEFEFEF}{324.00}{297.00}{541.50}{312.00}
\gF{cEFEFEF}{324.00}{432.00}{541.50}{447.00}
\gL{c000000}{1.00}{78.00}{189.50}{84.00}{189.50}
\gL{c000000}{1.00}{78.00}{324.50}{86.00}{324.50}
\gL{c000000}{1.00}{78.00}{444.50}{86.00}{444.50}
\gL{c000000}{1.00}{78.00}{564.50}{82.00}{564.50}
\gL{c000000}{1.00}{330.00}{204.50}{335.00}{204.50}
\gL{c000000}{1.00}{330.00}{264.50}{338.00}{264.50}
\gL{c000000}{1.00}{330.00}{309.50}{337.00}{309.50}
\gL{c000000}{1.00}{330.00}{444.50}{339.00}{444.50}
\gL{c000000}{1.00}{72.50}{72.00}{72.50}{717.00}
\gL{c000000}{1.00}{182.50}{72.00}{182.50}{177.00}
\gL{c000000}{1.00}{182.50}{192.00}{182.50}{312.00}
\gL{c000000}{1.00}{182.50}{327.00}{182.50}{432.00}
\gL{c000000}{1.00}{182.50}{447.00}{182.50}{552.00}
\gL{c000000}{1.00}{182.50}{567.00}{182.50}{717.00}
\gL{c000000}{1.00}{289.50}{72.00}{289.50}{717.00}
\gL{c000000}{1.00}{72.00}{72.50}{290.00}{72.50}
\gL{c000000}{1.00}{72.00}{87.50}{290.00}{87.50}
\gL{c000000}{1.00}{72.00}{102.50}{290.00}{102.50}
\gL{c000000}{1.00}{72.00}{117.50}{290.00}{117.50}
\gL{c000000}{1.00}{72.00}{132.50}{290.00}{132.50}
\gL{c000000}{1.00}{72.00}{147.50}{290.00}{147.50}
\gL{c000000}{1.00}{72.00}{162.50}{290.00}{162.50}
\gL{c000000}{1.00}{72.00}{177.50}{290.00}{177.50}
\gL{c000000}{1.00}{72.00}{192.50}{290.00}{192.50}
\gL{c000000}{1.00}{72.00}{207.50}{290.00}{207.50}
\gL{c000000}{1.00}{72.00}{222.50}{290.00}{222.50}
\gL{c000000}{1.00}{72.00}{237.50}{290.00}{237.50}
\gL{c000000}{1.00}{72.00}{252.50}{290.00}{252.50}
\gL{c000000}{1.00}{72.00}{267.50}{290.00}{267.50}
\gL{c000000}{1.00}{72.00}{282.50}{290.00}{282.50}
\gL{c000000}{1.00}{72.00}{297.50}{290.00}{297.50}
\gL{c000000}{1.00}{72.00}{312.50}{290.00}{312.50}
\gL{c000000}{1.00}{72.00}{327.50}{290.00}{327.50}
\gL{c000000}{1.00}{72.00}{342.50}{290.00}{342.50}
\gL{c000000}{1.00}{72.00}{357.50}{290.00}{357.50}
\gL{c000000}{1.00}{72.00}{372.50}{290.00}{372.50}
\gL{c000000}{1.00}{72.00}{387.50}{290.00}{387.50}
\gL{c000000}{1.00}{72.00}{402.50}{290.00}{402.50}
\gL{c000000}{1.00}{72.00}{417.50}{290.00}{417.50}
\gL{c000000}{1.00}{72.00}{432.50}{290.00}{432.50}
\gL{c000000}{1.00}{72.00}{447.50}{290.00}{447.50}
\gL{c000000}{1.00}{72.00}{462.50}{290.00}{462.50}
\gL{c000000}{1.00}{72.00}{477.50}{290.00}{477.50}
\gL{c000000}{1.00}{72.00}{492.50}{290.00}{492.50}
\gL{c000000}{1.00}{72.00}{507.50}{290.00}{507.50}
\gL{c000000}{1.00}{72.00}{522.50}{290.00}{522.50}
\gL{c000000}{1.00}{72.00}{537.50}{290.00}{537.50}
\gL{c000000}{1.00}{72.00}{552.50}{290.00}{552.50}
\gL{c000000}{1.00}{72.00}{567.50}{290.00}{567.50}
\gL{c000000}{1.00}{72.00}{582.50}{290.00}{582.50}
\gL{c000000}{1.00}{72.00}{597.50}{290.00}{597.50}
\gL{c000000}{1.00}{72.00}{612.50}{290.00}{612.50}
\gL{c000000}{1.00}{72.00}{627.50}{290.00}{627.50}
\gL{c000000}{1.00}{72.00}{642.50}{290.00}{642.50}
\gL{c000000}{1.00}{72.00}{657.50}{290.00}{657.50}
\gL{c000000}{1.00}{72.00}{672.50}{290.00}{672.50}
\gL{c000000}{1.00}{72.00}{687.50}{290.00}{687.50}
\gL{c000000}{1.00}{72.00}{702.50}{290.00}{702.50}
\gL{c000000}{1.00}{72.00}{717.50}{290.00}{717.50}
\gL{c000000}{1.00}{324.50}{72.00}{324.50}{717.00}
\gL{c000000}{1.00}{434.50}{72.00}{434.50}{192.00}
\gL{c000000}{1.00}{434.50}{207.00}{434.50}{252.00}
\gL{c000000}{1.00}{434.50}{267.00}{434.50}{297.00}
\gL{c000000}{1.00}{434.50}{312.00}{434.50}{432.00}
\gL{c000000}{1.00}{434.50}{447.00}{434.50}{717.00}
\gL{c000000}{1.00}{541.50}{72.00}{541.50}{717.00}
\gL{c000000}{1.00}{324.00}{72.50}{542.00}{72.50}
\gL{c000000}{1.00}{324.00}{87.50}{542.00}{87.50}
\gL{c000000}{1.00}{324.00}{102.50}{542.00}{102.50}
\gL{c000000}{1.00}{324.00}{117.50}{542.00}{117.50}
\gL{c000000}{1.00}{324.00}{132.50}{542.00}{132.50}
\gL{c000000}{1.00}{324.00}{147.50}{542.00}{147.50}
\gL{c000000}{1.00}{324.00}{162.50}{542.00}{162.50}
\gL{c000000}{1.00}{324.00}{177.50}{542.00}{177.50}
\gL{c000000}{1.00}{324.00}{192.50}{542.00}{192.50}
\gL{c000000}{1.00}{324.00}{207.50}{542.00}{207.50}
\gL{c000000}{1.00}{324.00}{222.50}{542.00}{222.50}
\gL{c000000}{1.00}{324.00}{237.50}{542.00}{237.50}
\gL{c000000}{1.00}{324.00}{252.50}{542.00}{252.50}
\gL{c000000}{1.00}{324.00}{267.50}{542.00}{267.50}
\gL{c000000}{1.00}{324.00}{282.50}{542.00}{282.50}
\gL{c000000}{1.00}{324.00}{297.50}{542.00}{297.50}
\gL{c000000}{1.00}{324.00}{312.50}{542.00}{312.50}
\gL{c000000}{1.00}{324.00}{327.50}{542.00}{327.50}
\gL{c000000}{1.00}{324.00}{342.50}{542.00}{342.50}
\gL{c000000}{1.00}{324.00}{357.50}{542.00}{357.50}
\gL{c000000}{1.00}{324.00}{372.50}{542.00}{372.50}
\gL{c000000}{1.00}{324.00}{387.50}{542.00}{387.50}
\gL{c000000}{1.00}{324.00}{402.50}{542.00}{402.50}
\gL{c000000}{1.00}{324.00}{417.50}{542.00}{417.50}
\gL{c000000}{1.00}{324.00}{432.50}{542.00}{432.50}
\gL{c000000}{1.00}{324.00}{447.50}{542.00}{447.50}
\gL{c000000}{1.00}{324.00}{462.50}{542.00}{462.50}
\gL{c000000}{1.00}{324.00}{477.50}{542.00}{477.50}
\gL{c000000}{1.00}{324.00}{492.50}{542.00}{492.50}
\gL{c000000}{1.00}{324.00}{507.50}{542.00}{507.50}
\gL{c000000}{1.00}{324.00}{522.50}{542.00}{522.50}
\gL{c000000}{1.00}{324.00}{537.50}{542.00}{537.50}
\gL{c000000}{1.00}{324.00}{552.50}{542.00}{552.50}
\gL{c000000}{1.00}{324.00}{567.50}{542.00}{567.50}
\gL{c000000}{1.00}{324.00}{582.50}{542.00}{582.50}
\gL{c000000}{1.00}{324.00}{597.50}{542.00}{597.50}
\gL{c000000}{1.00}{324.00}{612.50}{542.00}{612.50}
\gL{c000000}{1.00}{324.00}{627.50}{542.00}{627.50}
\gL{c000000}{1.00}{324.00}{642.50}{542.00}{642.50}
\gL{c000000}{1.00}{324.00}{657.50}{542.00}{657.50}
\gL{c000000}{1.00}{324.00}{672.50}{542.00}{672.50}
\gL{c000000}{1.00}{324.00}{687.50}{542.00}{687.50}
\gL{c000000}{1.00}{324.00}{702.50}{542.00}{702.50}
\gL{c000000}{1.00}{324.00}{717.50}{542.00}{717.50}
\gT{532.01}{43.47}{\gS{R}{8}{c000000}{25 }}
\gT{78.00}{83.46}{\gS{R}{10}{c000000}{Evolutionary }}
\gT{188.25}{83.46}{\gS{R}{10}{c000000}{Evol. }}
\gT{78.00}{98.46}{\gS{R}{10}{c000000}{Exhibition }}
\gT{188.25}{98.46}{\gS{R}{10}{c000000}{Exhib. }}
\gT{78.00}{113.46}{\gS{R}{10}{c000000}{Experimental }}
\gT{188.25}{113.46}{\gS{R}{10}{c000000}{Exp. }}
\gT{78.00}{128.46}{\gS{R}{10}{c000000}{Exploratory }}
\gT{188.25}{128.46}{\gS{R}{10}{c000000}{Explor. }}
\gT{78.00}{143.46}{\gS{R}{10}{c000000}{Exposition }}
\gT{188.25}{143.46}{\gS{R}{10}{c000000}{Expo. }}
\gT{78.00}{158.46}{\gS{R}{10}{c000000}{Express }}
\gT{188.25}{158.46}{\gS{R}{10}{c000000}{Exp. }}
\gT{78.00}{188.46}{\gS{RB}{10}{c000000}{F }}
\gT{78.00}{203.46}{\gS{R}{10}{c000000}{Fabrication }}
\gT{188.25}{203.46}{\gS{R}{10}{c000000}{Fabr. }}
\gT{78.00}{218.46}{\gS{R}{10}{c000000}{Faculty }}
\gT{188.25}{218.46}{\gS{R}{10}{c000000}{Fac. }}
\gT{78.00}{233.46}{\gS{R}{10}{c000000}{Ferroelectrics }}
\gT{188.25}{233.46}{\gS{R}{10}{c000000}{Ferroelect. }}
\gT{78.00}{248.46}{\gS{R}{10}{c000000}{Francais, French }}
\gT{188.25}{248.46}{\gS{R}{10}{c000000}{Fr. }}
\gT{78.00}{263.46}{\gS{R}{10}{c000000}{Frequency }}
\gT{188.25}{263.46}{\gS{R}{10}{c000000}{Freq. }}
\gT{78.00}{278.46}{\gS{R}{10}{c000000}{Foundation }}
\gT{188.25}{278.46}{\gS{R}{10}{c000000}{Found. }}
\gT{78.00}{293.46}{\gS{R}{10}{c000000}{Fundamental }}
\gT{188.25}{293.46}{\gS{R}{10}{c000000}{Fundam. }}
\gT{78.00}{323.46}{\gS{RB}{10}{c000000}{G }}
\gT{78.00}{338.46}{\gS{R}{10}{c000000}{Generation }}
\gT{188.25}{338.46}{\gS{R}{10}{c000000}{Gener. }}
\gT{78.00}{353.46}{\gS{R}{10}{c000000}{Geology }}
\gT{188.25}{353.46}{\gS{R}{10}{c000000}{Geol. }}
\gT{78.00}{368.46}{\gS{R}{10}{c000000}{Geophysics }}
\gT{188.25}{368.46}{\gS{R}{10}{c000000}{Geophys. }}
\gT{78.00}{383.46}{\gS{R}{10}{c000000}{Geoscience }}
\gT{188.25}{383.46}{\gS{R}{10}{c000000}{Geosci. }}
\gT{78.00}{398.46}{\gS{R}{10}{c000000}{Graphics }}
\gT{188.25}{398.46}{\gS{R}{10}{c000000}{Graph. }}
\gT{78.00}{413.46}{\gS{R}{10}{c000000}{Guidance }}
\gT{188.25}{413.46}{\gS{R}{10}{c000000}{Guid. }}
\gT{78.00}{443.46}{\gS{RB}{10}{c000000}{H }}
\gT{78.00}{458.46}{\gS{R}{10}{c000000}{Harmonic(s) }}
\gT{188.25}{458.46}{\gS{R}{10}{c000000}{Harmon. }}
\gT{78.00}{473.46}{\gS{R}{10}{c000000}{History }}
\gT{188.25}{473.46}{\gS{R}{10}{c000000}{Hist. }}
\gT{78.00}{488.46}{\gS{R}{10}{c000000}{Horizon }}
\gT{188.25}{488.46}{\gS{R}{10}{c000000}{Horiz. }}
\gT{78.00}{503.46}{\gS{R}{10}{c000000}{Hungary, Hungarian }}
\gT{188.25}{503.46}{\gS{R}{10}{c000000}{Hung. }}
\gT{78.00}{518.46}{\gS{R}{10}{c000000}{Hydraulics }}
\gT{188.25}{518.46}{\gS{R}{10}{c000000}{Hydraul. }}
\gT{78.00}{533.46}{\gS{R}{10}{c000000}{Hydrology }}
\gT{188.25}{533.46}{\gS{R}{10}{c000000}{Hydrol. }}
\gT{78.00}{563.46}{\gS{RB}{10}{c000000}{I }}
\gT{78.00}{578.46}{\gS{R}{10}{c000000}{Illuminating }}
\gT{188.25}{578.46}{\gS{R}{10}{c000000}{Illum. }}
\gT{78.00}{593.46}{\gS{R}{10}{c000000}{Imaging }}
\gT{188.25}{593.46}{\gS{R}{10}{c000000}{Imag. }}
\gT{78.00}{608.46}{\gS{R}{10}{c000000}{Industrial }}
\gT{188.25}{608.46}{\gS{R}{10}{c000000}{Ind. }}
\gT{78.00}{623.46}{\gS{R}{10}{c000000}{Information }}
\gT{188.25}{623.46}{\gS{R}{10}{c000000}{Inf. }}
\gT{78.00}{638.46}{\gS{R}{10}{c000000}{Informatics }}
\gT{188.25}{638.46}{\gS{R}{10}{c000000}{Inform. }}
\gT{78.00}{653.46}{\gS{R}{10}{c000000}{Innovation }}
\gT{188.25}{653.46}{\gS{R}{10}{c000000}{Innov. }}
\gT{78.00}{668.46}{\gS{R}{10}{c000000}{Institute }}
\gT{188.25}{668.46}{\gS{R}{10}{c000000}{Inst. }}
\gT{78.00}{683.46}{\gS{R}{10}{c000000}{Instrument }}
\gT{188.25}{683.46}{\gS{R}{10}{c000000}{Instrum. }}
\gT{78.00}{698.46}{\gS{R}{10}{c000000}{Instrumentation }}
\gT{188.25}{698.46}{\gS{R}{10}{c000000}{Instrum. }}
\gT{78.00}{713.46}{\gS{R}{10}{c000000}{Insulation }}
\gT{188.25}{713.46}{\gS{R}{10}{c000000}{Insul. }}
\gT{330.00}{83.46}{\gS{R}{10}{c000000}{Integrated }}
\gT{440.25}{83.46}{\gS{R}{10}{c000000}{Integr. }}
\gT{330.00}{98.46}{\gS{R}{10}{c000000}{Intelligence }}
\gT{440.25}{98.46}{\gS{R}{10}{c000000}{Intell. }}
\gT{330.00}{113.46}{\gS{R}{10}{c000000}{Intelligent }}
\gT{440.25}{113.46}{\gS{R}{10}{c000000}{Intell. }}
\gT{330.00}{128.46}{\gS{R}{10}{c000000}{Interactions }}
\gT{440.25}{128.46}{\gS{R}{10}{c000000}{Interact. }}
\gT{330.00}{143.46}{\gS{R}{10}{c000000}{International }}
\gT{440.25}{143.46}{\gS{R}{10}{c000000}{Int. }}
\gT{330.00}{158.46}{\gS{R}{10}{c000000}{Isotopes }}
\gT{440.25}{158.46}{\gS{R}{10}{c000000}{Isot. }}
\gT{330.00}{173.46}{\gS{R}{10}{c000000}{Israel }}
\gT{440.25}{173.46}{\gS{R}{10}{c000000}{Isr. }}
\gT{330.00}{203.46}{\gS{RB}{10}{c000000}{J }}
\gT{330.00}{218.46}{\gS{R}{10}{c000000}{Japan }}
\gT{440.25}{218.46}{\gS{R}{10}{c000000}{Jpn. }}
\gT{330.00}{233.46}{\gS{R}{10}{c000000}{Journal }}
\gT{440.25}{233.46}{\gS{R}{10}{c000000}{J. }}
\gT{330.00}{263.46}{\gS{RB}{10}{c000000}{K }}
\gT{330.00}{278.46}{\gS{R}{10}{c000000}{Knowledge }}
\gT{440.25}{278.46}{\gS{R}{10}{c000000}{Knowl. }}
\gT{330.00}{308.46}{\gS{RB}{10}{c000000}{L }}
\gT{330.00}{323.46}{\gS{R}{10}{c000000}{Laboratory(ies) }}
\gT{440.25}{323.46}{\gS{R}{10}{c000000}{Lab. }}
\gT{330.00}{338.46}{\gS{R}{10}{c000000}{Language }}
\gT{440.25}{338.46}{\gS{R}{10}{c000000}{Lang. }}
\gT{330.00}{353.46}{\gS{R}{10}{c000000}{Learning }}
\gT{440.25}{353.46}{\gS{R}{10}{c000000}{Learn. }}
\gT{330.00}{368.46}{\gS{R}{10}{c000000}{Letter(s) }}
\gT{440.25}{368.46}{\gS{R}{10}{c000000}{Lett. }}
\gT{330.00}{383.46}{\gS{R}{10}{c000000}{Lightwave }}
\gT{440.25}{383.46}{\gS{R}{10}{c000000}{Lightw. }}
\gT{330.00}{398.46}{\gS{R}{10}{c000000}{Logic, Logical }}
\gT{440.25}{398.46}{\gS{R}{10}{c000000}{Log. }}
\gT{330.00}{413.46}{\gS{R}{10}{c000000}{Luminescence }}
\gT{440.25}{413.46}{\gS{R}{10}{c000000}{Lumin. }}
\gT{330.00}{443.46}{\gS{RB}{10}{c000000}{M }}
\gT{330.00}{458.46}{\gS{R}{10}{c000000}{Machine }}
\gT{440.25}{458.46}{\gS{R}{10}{c000000}{Mach. }}
\gT{330.00}{473.46}{\gS{R}{10}{c000000}{Magazine }}
\gT{440.25}{473.46}{\gS{R}{10}{c000000}{Mag. }}
\gT{330.00}{488.46}{\gS{R}{10}{c000000}{Magnetics }}
\gT{440.25}{488.46}{\gS{R}{10}{c000000}{Magn. }}
\gT{330.00}{503.46}{\gS{R}{10}{c000000}{Management }}
\gT{440.25}{503.46}{\gS{R}{10}{c000000}{Manage. }}
\gT{330.00}{518.46}{\gS{R}{10}{c000000}{Managing }}
\gT{440.25}{518.46}{\gS{R}{10}{c000000}{Manag. }}
\gT{330.00}{533.46}{\gS{R}{10}{c000000}{Manufacturing }}
\gT{440.25}{533.46}{\gS{R}{10}{c000000}{Manuf. }}
\gT{330.00}{548.46}{\gS{R}{10}{c000000}{Marine }}
\gT{440.25}{548.46}{\gS{R}{10}{c000000}{Mar. }}
\gT{330.00}{563.46}{\gS{R}{10}{c000000}{Material }}
\gT{440.25}{563.46}{\gS{R}{10}{c000000}{Mater. }}
\gT{330.00}{578.46}{\gS{R}{10}{c000000}{Mathematical }}
\gT{440.25}{578.46}{\gS{R}{10}{c000000}{Math. }}
\gT{330.00}{593.46}{\gS{R}{10}{c000000}{Mathematics }}
\gT{440.25}{593.46}{\gS{R}{10}{c000000}{Math. }}
\gT{330.00}{608.46}{\gS{R}{10}{c000000}{Measurement }}
\gT{440.25}{608.46}{\gS{R}{10}{c000000}{Meas. }}
\gT{330.00}{623.46}{\gS{R}{10}{c000000}{Mechanical }}
\gT{440.25}{623.46}{\gS{R}{10}{c000000}{Mech. }}
\gT{330.00}{638.46}{\gS{R}{10}{c000000}{Medical, Medicine }}
\gT{440.25}{638.46}{\gS{R}{10}{c000000}{Med. }}
\gT{330.00}{653.46}{\gS{R}{10}{c000000}{Metals }}
\gT{440.25}{653.46}{\gS{R}{10}{c000000}{Met. }}
\gT{330.00}{668.46}{\gS{R}{10}{c000000}{Metallurgy }}
\gT{440.25}{668.46}{\gS{R}{10}{c000000}{Metall. }}
\gT{330.00}{683.46}{\gS{R}{10}{c000000}{Meteorology }}
\gT{440.25}{683.46}{\gS{R}{10}{c000000}{Meteorol. }}
\gT{330.00}{698.46}{\gS{R}{10}{c000000}{Metropolitan }}
\gT{440.25}{698.46}{\gS{R}{10}{c000000}{Metrop. }}
\gT{330.00}{713.46}{\gS{R}{10}{c000000}{Mexican, Mexico }}
\gT{440.25}{713.46}{\gS{R}{10}{c000000}{Mex. }}
\end{guidepage}
\begin{guidepage}% page 26
\gF{cEFEFEF}{72.00}{222.00}{289.50}{237.00}
\gF{cEFEFEF}{72.00}{402.00}{289.50}{417.00}
\gF{cEFEFEF}{72.00}{567.00}{289.50}{582.00}
\gF{cEFEFEF}{324.00}{357.00}{541.50}{372.00}
\gF{cEFEFEF}{324.00}{417.00}{541.50}{432.00}
\gF{cEFEFEF}{324.00}{672.00}{541.50}{687.00}
\gL{c000000}{1.00}{78.00}{234.50}{85.00}{234.50}
\gL{c000000}{1.00}{78.00}{414.50}{86.00}{414.50}
\gL{c000000}{1.00}{78.00}{579.50}{84.00}{579.50}
\gL{c000000}{1.00}{330.00}{369.50}{338.00}{369.50}
\gL{c000000}{1.00}{330.00}{429.50}{337.00}{429.50}
\gL{c000000}{1.00}{330.00}{684.50}{336.00}{684.50}
\gL{c000000}{1.00}{72.50}{72.00}{72.50}{717.00}
\gL{c000000}{1.00}{182.50}{72.00}{182.50}{222.00}
\gL{c000000}{1.00}{182.50}{237.00}{182.50}{402.00}
\gL{c000000}{1.00}{182.50}{417.00}{182.50}{567.00}
\gL{c000000}{1.00}{182.50}{582.00}{182.50}{717.00}
\gL{c000000}{1.00}{289.50}{72.00}{289.50}{717.00}
\gL{c000000}{1.00}{72.00}{72.50}{290.00}{72.50}
\gL{c000000}{1.00}{72.00}{87.50}{290.00}{87.50}
\gL{c000000}{1.00}{72.00}{102.50}{290.00}{102.50}
\gL{c000000}{1.00}{72.00}{117.50}{290.00}{117.50}
\gL{c000000}{1.00}{72.00}{132.50}{290.00}{132.50}
\gL{c000000}{1.00}{72.00}{147.50}{290.00}{147.50}
\gL{c000000}{1.00}{72.00}{162.50}{290.00}{162.50}
\gL{c000000}{1.00}{72.00}{177.50}{290.00}{177.50}
\gL{c000000}{1.00}{72.00}{192.50}{290.00}{192.50}
\gL{c000000}{1.00}{72.00}{207.50}{290.00}{207.50}
\gL{c000000}{1.00}{72.00}{222.50}{290.00}{222.50}
\gL{c000000}{1.00}{72.00}{237.50}{290.00}{237.50}
\gL{c000000}{1.00}{72.00}{252.50}{290.00}{252.50}
\gL{c000000}{1.00}{72.00}{267.50}{290.00}{267.50}
\gL{c000000}{1.00}{72.00}{282.50}{290.00}{282.50}
\gL{c000000}{1.00}{72.00}{297.50}{290.00}{297.50}
\gL{c000000}{1.00}{72.00}{312.50}{290.00}{312.50}
\gL{c000000}{1.00}{72.00}{327.50}{290.00}{327.50}
\gL{c000000}{1.00}{72.00}{342.50}{290.00}{342.50}
\gL{c000000}{1.00}{72.00}{357.50}{290.00}{357.50}
\gL{c000000}{1.00}{72.00}{372.50}{290.00}{372.50}
\gL{c000000}{1.00}{72.00}{387.50}{290.00}{387.50}
\gL{c000000}{1.00}{72.00}{402.50}{290.00}{402.50}
\gL{c000000}{1.00}{72.00}{417.50}{290.00}{417.50}
\gL{c000000}{1.00}{72.00}{432.50}{290.00}{432.50}
\gL{c000000}{1.00}{72.00}{447.50}{290.00}{447.50}
\gL{c000000}{1.00}{72.00}{462.50}{290.00}{462.50}
\gL{c000000}{1.00}{72.00}{477.50}{290.00}{477.50}
\gL{c000000}{1.00}{72.00}{492.50}{290.00}{492.50}
\gL{c000000}{1.00}{72.00}{507.50}{290.00}{507.50}
\gL{c000000}{1.00}{72.00}{522.50}{290.00}{522.50}
\gL{c000000}{1.00}{72.00}{537.50}{290.00}{537.50}
\gL{c000000}{1.00}{72.00}{552.50}{290.00}{552.50}
\gL{c000000}{1.00}{72.00}{567.50}{290.00}{567.50}
\gL{c000000}{1.00}{72.00}{582.50}{290.00}{582.50}
\gL{c000000}{1.00}{72.00}{597.50}{290.00}{597.50}
\gL{c000000}{1.00}{72.00}{612.50}{290.00}{612.50}
\gL{c000000}{1.00}{72.00}{627.50}{290.00}{627.50}
\gL{c000000}{1.00}{72.00}{642.50}{290.00}{642.50}
\gL{c000000}{1.00}{72.00}{657.50}{290.00}{657.50}
\gL{c000000}{1.00}{72.00}{672.50}{290.00}{672.50}
\gL{c000000}{1.00}{72.00}{687.50}{290.00}{687.50}
\gL{c000000}{1.00}{72.00}{702.50}{290.00}{702.50}
\gL{c000000}{1.00}{72.00}{717.50}{290.00}{717.50}
\gL{c000000}{1.00}{324.50}{72.00}{324.50}{717.00}
\gL{c000000}{1.00}{434.50}{72.00}{434.50}{357.00}
\gL{c000000}{1.00}{434.50}{372.00}{434.50}{417.00}
\gL{c000000}{1.00}{434.50}{432.00}{434.50}{672.00}
\gL{c000000}{1.00}{434.50}{687.00}{434.50}{717.00}
\gL{c000000}{1.00}{541.50}{72.00}{541.50}{717.00}
\gL{c000000}{1.00}{324.00}{72.50}{542.00}{72.50}
\gL{c000000}{1.00}{324.00}{87.50}{542.00}{87.50}
\gL{c000000}{1.00}{324.00}{102.50}{542.00}{102.50}
\gL{c000000}{1.00}{324.00}{117.50}{542.00}{117.50}
\gL{c000000}{1.00}{324.00}{132.50}{542.00}{132.50}
\gL{c000000}{1.00}{324.00}{147.50}{542.00}{147.50}
\gL{c000000}{1.00}{324.00}{162.50}{542.00}{162.50}
\gL{c000000}{1.00}{324.00}{177.50}{542.00}{177.50}
\gL{c000000}{1.00}{324.00}{192.50}{542.00}{192.50}
\gL{c000000}{1.00}{324.00}{207.50}{542.00}{207.50}
\gL{c000000}{1.00}{324.00}{222.50}{542.00}{222.50}
\gL{c000000}{1.00}{324.00}{237.50}{542.00}{237.50}
\gL{c000000}{1.00}{324.00}{252.50}{542.00}{252.50}
\gL{c000000}{1.00}{324.00}{267.50}{542.00}{267.50}
\gL{c000000}{1.00}{324.00}{282.50}{542.00}{282.50}
\gL{c000000}{1.00}{324.00}{297.50}{542.00}{297.50}
\gL{c000000}{1.00}{324.00}{312.50}{542.00}{312.50}
\gL{c000000}{1.00}{324.00}{327.50}{542.00}{327.50}
\gL{c000000}{1.00}{324.00}{342.50}{542.00}{342.50}
\gL{c000000}{1.00}{324.00}{357.50}{542.00}{357.50}
\gL{c000000}{1.00}{324.00}{372.50}{542.00}{372.50}
\gL{c000000}{1.00}{324.00}{387.50}{542.00}{387.50}
\gL{c000000}{1.00}{324.00}{402.50}{542.00}{402.50}
\gL{c000000}{1.00}{324.00}{417.50}{542.00}{417.50}
\gL{c000000}{1.00}{324.00}{432.50}{542.00}{432.50}
\gL{c000000}{1.00}{324.00}{447.50}{542.00}{447.50}
\gL{c000000}{1.00}{324.00}{462.50}{542.00}{462.50}
\gL{c000000}{1.00}{324.00}{477.50}{542.00}{477.50}
\gL{c000000}{1.00}{324.00}{492.50}{542.00}{492.50}
\gL{c000000}{1.00}{324.00}{507.50}{542.00}{507.50}
\gL{c000000}{1.00}{324.00}{522.50}{542.00}{522.50}
\gL{c000000}{1.00}{324.00}{537.50}{542.00}{537.50}
\gL{c000000}{1.00}{324.00}{552.50}{542.00}{552.50}
\gL{c000000}{1.00}{324.00}{567.50}{542.00}{567.50}
\gL{c000000}{1.00}{324.00}{582.50}{542.00}{582.50}
\gL{c000000}{1.00}{324.00}{597.50}{542.00}{597.50}
\gL{c000000}{1.00}{324.00}{612.50}{542.00}{612.50}
\gL{c000000}{1.00}{324.00}{627.50}{542.00}{627.50}
\gL{c000000}{1.00}{324.00}{642.50}{542.00}{642.50}
\gL{c000000}{1.00}{324.00}{657.50}{542.00}{657.50}
\gL{c000000}{1.00}{324.00}{672.50}{542.00}{672.50}
\gL{c000000}{1.00}{324.00}{687.50}{542.00}{687.50}
\gL{c000000}{1.00}{324.00}{702.50}{542.00}{702.50}
\gL{c000000}{1.00}{324.00}{717.50}{542.00}{717.50}
\gT{532.01}{43.47}{\gS{R}{8}{c000000}{26 }}
\gT{78.00}{83.46}{\gS{R}{10}{c000000}{Microelectromechanical }}
\gT{188.25}{83.46}{\gS{R}{10}{c000000}{Microelectromech. }}
\gT{78.00}{98.46}{\gS{R}{10}{c000000}{Microgravity }}
\gT{188.25}{98.46}{\gS{R}{10}{c000000}{Microgr. }}
\gT{78.00}{113.46}{\gS{R}{10}{c000000}{Microscopy }}
\gT{188.25}{113.46}{\gS{R}{10}{c000000}{Microsc. }}
\gT{78.00}{128.46}{\gS{R}{10}{c000000}{Microwave(s) }}
\gT{188.25}{128.46}{\gS{R}{10}{c000000}{Microw. }}
\gT{78.00}{143.46}{\gS{R}{10}{c000000}{Military }}
\gT{188.25}{143.46}{\gS{R}{10}{c000000}{Mil. }}
\gT{78.00}{158.46}{\gS{R}{10}{c000000}{Modeling }}
\gT{188.25}{158.46}{\gS{R}{10}{c000000}{Model. }}
\gT{78.00}{173.46}{\gS{R}{10}{c000000}{Molecular }}
\gT{188.25}{173.46}{\gS{R}{10}{c000000}{Mol. }}
\gT{78.00}{188.46}{\gS{R}{10}{c000000}{Monitoring }}
\gT{188.25}{188.46}{\gS{R}{10}{c000000}{Monit. }}
\gT{78.00}{203.46}{\gS{R}{10}{c000000}{Multiphysics }}
\gT{188.25}{203.46}{\gS{R}{10}{c000000}{Multiphys. }}
\gT{78.00}{233.46}{\gS{RB}{10}{c000000}{N }}
\gT{78.00}{248.46}{\gS{R}{10}{c000000}{Nanobioscience }}
\gT{188.25}{248.46}{\gS{R}{10}{c000000}{Nanobiosci. }}
\gT{78.00}{263.46}{\gS{R}{10}{c000000}{Nanotechnology }}
\gT{188.25}{263.46}{\gS{R}{10}{c000000}{Nanotechnol. }}
\gT{78.00}{278.46}{\gS{R}{10}{c000000}{National }}
\gT{188.25}{278.46}{\gS{R}{10}{c000000}{Nat. }}
\gT{78.00}{293.46}{\gS{R}{10}{c000000}{Naval }}
\gT{188.25}{293.46}{\gS{R}{10}{c000000}{Nav. }}
\gT{78.00}{308.46}{\gS{R}{10}{c000000}{Navigation }}
\gT{188.25}{308.46}{\gS{R}{10}{c000000}{Navig. }}
\gT{78.00}{323.46}{\gS{R}{10}{c000000}{Network, Networking }}
\gT{188.25}{323.46}{\gS{R}{10}{c000000}{Netw. }}
\gT{78.00}{338.46}{\gS{R}{10}{c000000}{Newsletter }}
\gT{188.25}{338.46}{\gS{R}{10}{c000000}{Newslett. }}
\gT{78.00}{353.46}{\gS{R}{10}{c000000}{Nondestructive }}
\gT{188.25}{353.46}{\gS{R}{10}{c000000}{Nondestruct. }}
\gT{78.00}{368.46}{\gS{R}{10}{c000000}{Nuclear }}
\gT{188.25}{368.46}{\gS{R}{10}{c000000}{Nucl. }}
\gT{78.00}{383.46}{\gS{R}{10}{c000000}{Numerical }}
\gT{188.25}{383.46}{\gS{R}{10}{c000000}{Numer. }}
\gT{78.00}{413.46}{\gS{RB}{10}{c000000}{O }}
\gT{78.00}{428.46}{\gS{R}{10}{c000000}{Observations }}
\gT{188.25}{428.46}{\gS{R}{10}{c000000}{Observ. }}
\gT{78.00}{443.46}{\gS{R}{10}{c000000}{Oceanic }}
\gT{188.25}{443.46}{\gS{R}{10}{c000000}{Ocean. }}
\gT{78.00}{458.46}{\gS{R}{10}{c000000}{Oceanography }}
\gT{188.25}{458.46}{\gS{R}{10}{c000000}{Oceanogr. }}
\gT{78.00}{473.46}{\gS{R}{10}{c000000}{Occupation }}
\gT{188.25}{473.46}{\gS{R}{10}{c000000}{Occupat. }}
\gT{78.00}{488.46}{\gS{R}{10}{c000000}{Operational }}
\gT{188.25}{488.46}{\gS{R}{10}{c000000}{Oper. }}
\gT{78.00}{503.46}{\gS{R}{10}{c000000}{Optical }}
\gT{188.25}{503.46}{\gS{R}{10}{c000000}{Opt. }}
\gT{78.00}{518.46}{\gS{R}{10}{c000000}{Optics }}
\gT{188.25}{518.46}{\gS{R}{10}{c000000}{Opt. }}
\gT{78.00}{533.46}{\gS{R}{10}{c000000}{Optimization }}
\gT{188.25}{533.46}{\gS{R}{10}{c000000}{Optim. }}
\gT{78.00}{548.46}{\gS{R}{10}{c000000}{Organization }}
\gT{188.25}{548.46}{\gS{R}{10}{c000000}{Org. }}
\gT{78.00}{578.46}{\gS{RB}{10}{c000000}{P }}
\gT{78.00}{593.46}{\gS{R}{10}{c000000}{Packaging }}
\gT{188.25}{593.46}{\gS{R}{10}{c000000}{Packag. }}
\gT{78.00}{608.46}{\gS{R}{10}{c000000}{Particle }}
\gT{188.25}{608.46}{\gS{R}{10}{c000000}{Part. }}
\gT{78.00}{623.46}{\gS{R}{10}{c000000}{Patent }}
\gT{188.25}{623.46}{\gS{R}{10}{c000000}{Pat. }}
\gT{78.00}{638.46}{\gS{R}{10}{c000000}{Performance }}
\gT{188.25}{638.46}{\gS{R}{10}{c000000}{Perform. }}
\gT{78.00}{653.46}{\gS{R}{10}{c000000}{Personal }}
\gT{188.25}{653.46}{\gS{R}{10}{c000000}{Pers. }}
\gT{78.00}{668.46}{\gS{R}{10}{c000000}{Philosophical }}
\gT{188.25}{668.46}{\gS{R}{10}{c000000}{Philos. }}
\gT{78.00}{683.46}{\gS{R}{10}{c000000}{Photonics }}
\gT{188.25}{683.46}{\gS{R}{10}{c000000}{Photon. }}
\gT{78.00}{698.46}{\gS{R}{10}{c000000}{Photovoltaics }}
\gT{188.25}{698.46}{\gS{R}{10}{c000000}{Photovolt. }}
\gT{78.00}{713.46}{\gS{R}{10}{c000000}{Physics }}
\gT{188.25}{713.46}{\gS{R}{10}{c000000}{Phys. }}
\gT{330.00}{83.46}{\gS{R}{10}{c000000}{Physiology }}
\gT{440.25}{83.46}{\gS{R}{10}{c000000}{Physiol. }}
\gT{330.00}{98.46}{\gS{R}{10}{c000000}{Planetary }}
\gT{440.25}{98.46}{\gS{R}{10}{c000000}{Planet. }}
\gT{330.00}{113.46}{\gS{R}{10}{c000000}{Pneumatics }}
\gT{440.25}{113.46}{\gS{R}{10}{c000000}{Pneum. }}
\gT{330.00}{128.46}{\gS{R}{10}{c000000}{Pollution }}
\gT{440.25}{128.46}{\gS{R}{10}{c000000}{Pollut. }}
\gT{330.00}{143.46}{\gS{R}{10}{c000000}{Polymer }}
\gT{440.25}{143.46}{\gS{R}{10}{c000000}{Polym. }}
\gT{330.00}{158.46}{\gS{R}{10}{c000000}{Polytechnic }}
\gT{440.25}{158.46}{\gS{R}{10}{c000000}{Polytech. }}
\gT{330.00}{173.46}{\gS{R}{10}{c000000}{Practice }}
\gT{440.25}{173.46}{\gS{R}{10}{c000000}{Pract. }}
\gT{330.00}{188.46}{\gS{R}{10}{c000000}{Precision }}
\gT{440.25}{188.46}{\gS{R}{10}{c000000}{Precis. }}
\gT{330.00}{203.46}{\gS{R}{10}{c000000}{Principles }}
\gT{440.25}{203.46}{\gS{R}{10}{c000000}{Princ. }}
\gT{330.00}{218.46}{\gS{R}{10}{c000000}{Proceedings }}
\gT{440.25}{218.46}{\gS{R}{10}{c000000}{Proc. }}
\gT{330.00}{233.46}{\gS{R}{10}{c000000}{Processing }}
\gT{440.25}{233.46}{\gS{R}{10}{c000000}{Process. }}
\gT{330.00}{248.46}{\gS{R}{10}{c000000}{Production }}
\gT{440.25}{248.46}{\gS{R}{10}{c000000}{Prod. }}
\gT{330.00}{263.46}{\gS{R}{10}{c000000}{Productivity }}
\gT{440.25}{263.46}{\gS{R}{10}{c000000}{Productiv. }}
\gT{330.00}{278.46}{\gS{R}{10}{c000000}{Programmable }}
\gT{440.25}{278.46}{\gS{R}{10}{c000000}{Program. }}
\gT{330.00}{293.46}{\gS{R}{10}{c000000}{Programming }}
\gT{440.25}{293.46}{\gS{R}{10}{c000000}{Program. }}
\gT{330.00}{308.46}{\gS{R}{10}{c000000}{Progress }}
\gT{440.25}{308.46}{\gS{R}{10}{c000000}{Prog. }}
\gT{330.00}{323.46}{\gS{R}{10}{c000000}{Propagation }}
\gT{440.25}{323.46}{\gS{R}{10}{c000000}{Propag. }}
\gT{330.00}{338.46}{\gS{R}{10}{c000000}{Psychology }}
\gT{440.25}{338.46}{\gS{R}{10}{c000000}{Psychol. }}
\gT{330.00}{368.46}{\gS{RB}{10}{c000000}{Q }}
\gT{330.00}{383.46}{\gS{R}{10}{c000000}{Quality }}
\gT{440.25}{383.46}{\gS{R}{10}{c000000}{Qual. }}
\gT{330.00}{398.46}{\gS{R}{10}{c000000}{Quarterly }}
\gT{440.25}{398.46}{\gS{R}{10}{c000000}{Quart. }}
\gT{330.00}{428.46}{\gS{RB}{10}{c000000}{R }}
\gT{330.00}{443.46}{\gS{R}{10}{c000000}{Radiation }}
\gT{440.25}{443.46}{\gS{R}{10}{c000000}{Radiat. }}
\gT{330.00}{458.46}{\gS{R}{10}{c000000}{Radiology }}
\gT{440.25}{458.46}{\gS{R}{10}{c000000}{Radiol. }}
\gT{330.00}{473.46}{\gS{R}{10}{c000000}{Reactor }}
\gT{440.25}{473.46}{\gS{R}{10}{c000000}{React. }}
\gT{330.00}{488.46}{\gS{R}{10}{c000000}{Receivers }}
\gT{440.25}{488.46}{\gS{R}{10}{c000000}{Receiv. }}
\gT{330.00}{503.46}{\gS{R}{10}{c000000}{Recognition }}
\gT{440.25}{503.46}{\gS{R}{10}{c000000}{Recognit. }}
\gT{330.00}{518.46}{\gS{R}{10}{c000000}{Record }}
\gT{440.25}{518.46}{\gS{R}{10}{c000000}{Rec. }}
\gT{330.00}{533.46}{\gS{R}{10}{c000000}{Rehabilitation }}
\gT{440.25}{533.46}{\gS{R}{10}{c000000}{Rehabil. }}
\gT{330.00}{548.46}{\gS{R}{10}{c000000}{Reliability }}
\gT{440.25}{548.46}{\gS{R}{10}{c000000}{Rel. }}
\gT{330.00}{563.46}{\gS{R}{10}{c000000}{Report }}
\gT{440.25}{563.46}{\gS{R}{10}{c000000}{Rep. }}
\gT{330.00}{578.46}{\gS{R}{10}{c000000}{Research }}
\gT{440.25}{578.46}{\gS{R}{10}{c000000}{Res. }}
\gT{330.00}{593.46}{\gS{R}{10}{c000000}{Resonance }}
\gT{440.25}{593.46}{\gS{R}{10}{c000000}{Reson. }}
\gT{330.00}{608.46}{\gS{R}{10}{c000000}{Resources }}
\gT{440.25}{608.46}{\gS{R}{10}{c000000}{Resour. }}
\gT{330.00}{623.46}{\gS{R}{10}{c000000}{Review }}
\gT{440.25}{623.46}{\gS{R}{10}{c000000}{Rev. }}
\gT{330.00}{638.46}{\gS{R}{10}{c000000}{Robotics }}
\gT{440.25}{638.46}{\gS{R}{10}{c000000}{Robot. }}
\gT{330.00}{653.46}{\gS{R}{10}{c000000}{Royal }}
\gT{440.25}{653.46}{\gS{R}{10}{c000000}{Roy. }}
\gT{330.00}{683.46}{\gS{RB}{10}{c000000}{S }}
\gT{330.00}{698.46}{\gS{R}{10}{c000000}{Safety }}
\gT{440.25}{698.46}{\gS{R}{10}{c000000}{Saf. }}
\gT{330.00}{713.46}{\gS{R}{10}{c000000}{Satellite }}
\gT{440.25}{713.46}{\gS{R}{10}{c000000}{Satell. }}
\end{guidepage}
\begin{guidepage}% page 27
\gF{cEFEFEF}{72.00}{537.00}{289.50}{552.00}
\gF{cEFEFEF}{324.00}{297.00}{541.50}{312.00}
\gF{cEFEFEF}{324.00}{357.00}{541.50}{372.00}
\gF{cEFEFEF}{324.00}{477.00}{541.50}{492.00}
\gL{c000000}{1.00}{78.00}{549.50}{85.00}{549.50}
\gL{c000000}{1.00}{330.00}{309.50}{337.00}{309.50}
\gL{c000000}{1.00}{330.00}{369.50}{337.00}{369.50}
\gL{c000000}{1.00}{330.00}{489.50}{340.00}{489.50}
\gL{c000000}{1.00}{72.50}{72.00}{72.50}{552.00}
\gL{c000000}{1.00}{182.50}{72.00}{182.50}{537.00}
\gL{c000000}{1.00}{289.50}{72.00}{289.50}{552.00}
\gL{c000000}{1.00}{72.00}{72.50}{290.00}{72.50}
\gL{c000000}{1.00}{72.00}{87.50}{290.00}{87.50}
\gL{c000000}{1.00}{72.00}{102.50}{290.00}{102.50}
\gL{c000000}{1.00}{72.00}{117.50}{290.00}{117.50}
\gL{c000000}{1.00}{72.00}{132.50}{290.00}{132.50}
\gL{c000000}{1.00}{72.00}{147.50}{290.00}{147.50}
\gL{c000000}{1.00}{72.00}{162.50}{290.00}{162.50}
\gL{c000000}{1.00}{72.00}{177.50}{290.00}{177.50}
\gL{c000000}{1.00}{72.00}{192.50}{290.00}{192.50}
\gL{c000000}{1.00}{72.00}{207.50}{290.00}{207.50}
\gL{c000000}{1.00}{72.00}{222.50}{290.00}{222.50}
\gL{c000000}{1.00}{72.00}{237.50}{290.00}{237.50}
\gL{c000000}{1.00}{72.00}{252.50}{290.00}{252.50}
\gL{c000000}{1.00}{72.00}{267.50}{290.00}{267.50}
\gL{c000000}{1.00}{72.00}{282.50}{290.00}{282.50}
\gL{c000000}{1.00}{72.00}{297.50}{290.00}{297.50}
\gL{c000000}{1.00}{72.00}{312.50}{290.00}{312.50}
\gL{c000000}{1.00}{72.00}{327.50}{290.00}{327.50}
\gL{c000000}{1.00}{72.00}{342.50}{290.00}{342.50}
\gL{c000000}{1.00}{72.00}{357.50}{290.00}{357.50}
\gL{c000000}{1.00}{72.00}{372.50}{290.00}{372.50}
\gL{c000000}{1.00}{72.00}{387.50}{290.00}{387.50}
\gL{c000000}{1.00}{72.00}{402.50}{290.00}{402.50}
\gL{c000000}{1.00}{72.00}{417.50}{290.00}{417.50}
\gL{c000000}{1.00}{72.00}{432.50}{290.00}{432.50}
\gL{c000000}{1.00}{72.00}{447.50}{290.00}{447.50}
\gL{c000000}{1.00}{72.00}{462.50}{290.00}{462.50}
\gL{c000000}{1.00}{72.00}{477.50}{290.00}{477.50}
\gL{c000000}{1.00}{72.00}{492.50}{290.00}{492.50}
\gL{c000000}{1.00}{72.00}{507.50}{290.00}{507.50}
\gL{c000000}{1.00}{72.00}{522.50}{290.00}{522.50}
\gL{c000000}{1.00}{72.00}{537.50}{290.00}{537.50}
\gL{c000000}{1.00}{72.00}{552.50}{290.00}{552.50}
\gL{c000000}{1.00}{324.50}{72.00}{324.50}{522.00}
\gL{c000000}{1.00}{434.50}{72.00}{434.50}{297.00}
\gL{c000000}{1.00}{434.50}{312.00}{434.50}{357.00}
\gL{c000000}{1.00}{434.50}{372.00}{434.50}{477.00}
\gL{c000000}{1.00}{434.50}{492.00}{434.50}{522.00}
\gL{c000000}{1.00}{541.50}{72.00}{541.50}{522.00}
\gL{c000000}{1.00}{324.00}{72.50}{542.00}{72.50}
\gL{c000000}{1.00}{324.00}{87.50}{542.00}{87.50}
\gL{c000000}{1.00}{324.00}{102.50}{542.00}{102.50}
\gL{c000000}{1.00}{324.00}{117.50}{542.00}{117.50}
\gL{c000000}{1.00}{324.00}{132.50}{542.00}{132.50}
\gL{c000000}{1.00}{324.00}{147.50}{542.00}{147.50}
\gL{c000000}{1.00}{324.00}{162.50}{542.00}{162.50}
\gL{c000000}{1.00}{324.00}{177.50}{542.00}{177.50}
\gL{c000000}{1.00}{324.00}{192.50}{542.00}{192.50}
\gL{c000000}{1.00}{324.00}{207.50}{542.00}{207.50}
\gL{c000000}{1.00}{324.00}{222.50}{542.00}{222.50}
\gL{c000000}{1.00}{324.00}{237.50}{542.00}{237.50}
\gL{c000000}{1.00}{324.00}{252.50}{542.00}{252.50}
\gL{c000000}{1.00}{324.00}{267.50}{542.00}{267.50}
\gL{c000000}{1.00}{324.00}{282.50}{542.00}{282.50}
\gL{c000000}{1.00}{324.00}{297.50}{542.00}{297.50}
\gL{c000000}{1.00}{324.00}{312.50}{542.00}{312.50}
\gL{c000000}{1.00}{324.00}{327.50}{542.00}{327.50}
\gL{c000000}{1.00}{324.00}{342.50}{542.00}{342.50}
\gL{c000000}{1.00}{324.00}{357.50}{542.00}{357.50}
\gL{c000000}{1.00}{324.00}{372.50}{542.00}{372.50}
\gL{c000000}{1.00}{324.00}{387.50}{542.00}{387.50}
\gL{c000000}{1.00}{324.00}{402.50}{542.00}{402.50}
\gL{c000000}{1.00}{324.00}{417.50}{542.00}{417.50}
\gL{c000000}{1.00}{324.00}{432.50}{542.00}{432.50}
\gL{c000000}{1.00}{324.00}{447.50}{542.00}{447.50}
\gL{c000000}{1.00}{324.00}{462.50}{542.00}{462.50}
\gL{c000000}{1.00}{324.00}{477.50}{542.00}{477.50}
\gL{c000000}{1.00}{324.00}{492.50}{542.00}{492.50}
\gL{c000000}{1.00}{324.00}{507.50}{542.00}{507.50}
\gL{c000000}{1.00}{324.00}{522.50}{542.00}{522.50}
\gT{532.01}{43.47}{\gS{R}{8}{c000000}{27 }}
\gT{78.00}{83.46}{\gS{R}{10}{c000000}{Scandinavian }}
\gT{188.25}{83.46}{\gS{R}{10}{c000000}{Scand. }}
\gT{78.00}{98.46}{\gS{R}{10}{c000000}{Science }}
\gT{188.25}{98.46}{\gS{R}{10}{c000000}{Sci. }}
\gT{78.00}{113.46}{\gS{R}{10}{c000000}{Section }}
\gT{188.25}{113.46}{\gS{R}{10}{c000000}{Sect. }}
\gT{78.00}{128.46}{\gS{R}{10}{c000000}{Security }}
\gT{188.25}{128.46}{\gS{R}{10}{c000000}{Secur. }}
\gT{78.00}{143.46}{\gS{R}{10}{c000000}{Seismology }}
\gT{188.25}{143.46}{\gS{R}{10}{c000000}{Seismol. }}
\gT{78.00}{158.46}{\gS{R}{10}{c000000}{Selected }}
\gT{188.25}{158.46}{\gS{R}{10}{c000000}{Sel. }}
\gT{78.00}{173.46}{\gS{R}{10}{c000000}{Semiconductor }}
\gT{188.25}{173.46}{\gS{R}{10}{c000000}{Semicond. }}
\gT{78.00}{188.46}{\gS{R}{10}{c000000}{Sensing }}
\gT{188.25}{188.46}{\gS{R}{10}{c000000}{Sens. }}
\gT{78.00}{203.46}{\gS{R}{10}{c000000}{Series }}
\gT{188.25}{203.46}{\gS{R}{10}{c000000}{Ser. }}
\gT{78.00}{218.46}{\gS{R}{10}{c000000}{Simulation }}
\gT{188.25}{218.46}{\gS{R}{10}{c000000}{Simul. }}
\gT{78.00}{233.46}{\gS{R}{10}{c000000}{Singapore }}
\gT{188.25}{233.46}{\gS{R}{10}{c000000}{Singap. }}
\gT{78.00}{248.46}{\gS{R}{10}{c000000}{Sistema }}
\gT{188.25}{248.46}{\gS{R}{10}{c000000}{Sist. }}
\gT{78.00}{263.46}{\gS{R}{10}{c000000}{Society }}
\gT{188.25}{263.46}{\gS{R}{10}{c000000}{Soc. }}
\gT{78.00}{278.46}{\gS{R}{10}{c000000}{Sociological }}
\gT{188.25}{278.46}{\gS{R}{10}{c000000}{Sociol. }}
\gT{78.00}{293.46}{\gS{R}{10}{c000000}{Software }}
\gT{188.25}{293.46}{\gS{R}{10}{c000000}{Softw. }}
\gT{78.00}{308.46}{\gS{R}{10}{c000000}{Solar }}
\gT{188.25}{308.46}{\gS{R}{10}{c000000}{Sol. }}
\gT{78.00}{323.46}{\gS{R}{10}{c000000}{Soviet }}
\gT{188.25}{323.46}{\gS{R}{10}{c000000}{Sov. }}
\gT{78.00}{338.46}{\gS{R}{10}{c000000}{Spectroscopy }}
\gT{188.25}{338.46}{\gS{R}{10}{c000000}{Spectrosc. }}
\gT{78.00}{353.46}{\gS{R}{10}{c000000}{Spectrum }}
\gT{188.25}{353.46}{\gS{R}{10}{c000000}{Spectr. }}
\gT{78.00}{368.46}{\gS{R}{10}{c000000}{Speculations }}
\gT{188.25}{368.46}{\gS{R}{10}{c000000}{Specul. }}
\gT{78.00}{383.46}{\gS{R}{10}{c000000}{Statistics }}
\gT{188.25}{383.46}{\gS{R}{10}{c000000}{Statist. }}
\gT{78.00}{398.46}{\gS{R}{10}{c000000}{Structure }}
\gT{188.25}{398.46}{\gS{R}{10}{c000000}{Struct. }}
\gT{78.00}{413.46}{\gS{R}{10}{c000000}{Studies }}
\gT{188.25}{413.46}{\gS{R}{10}{c000000}{Stud. }}
\gT{78.00}{428.46}{\gS{R}{10}{c000000}{Superconductivity }}
\gT{188.25}{428.46}{\gS{R}{10}{c000000}{Supercond. }}
\gT{78.00}{443.46}{\gS{R}{10}{c000000}{Supplement }}
\gT{188.25}{443.46}{\gS{R}{10}{c000000}{Suppl. }}
\gT{78.00}{458.46}{\gS{R}{10}{c000000}{Surface }}
\gT{188.25}{458.46}{\gS{R}{10}{c000000}{Surf. }}
\gT{78.00}{473.46}{\gS{R}{10}{c000000}{Survey }}
\gT{188.25}{473.46}{\gS{R}{10}{c000000}{Surv. }}
\gT{78.00}{488.46}{\gS{R}{10}{c000000}{Sustainable }}
\gT{188.25}{488.46}{\gS{R}{10}{c000000}{Sustain. }}
\gT{78.00}{503.46}{\gS{R}{10}{c000000}{Symposium }}
\gT{188.25}{503.46}{\gS{R}{10}{c000000}{Symp. }}
\gT{78.00}{518.46}{\gS{R}{10}{c000000}{Systems }}
\gT{188.25}{518.46}{\gS{R}{10}{c000000}{Syst. }}
\gT{78.00}{548.46}{\gS{RB}{10}{c000000}{T }}
\gT{330.00}{83.46}{\gS{R}{10}{c000000}{Technical }}
\gT{440.25}{83.46}{\gS{R}{10}{c000000}{Tech. }}
\gT{330.00}{98.46}{\gS{R}{10}{c000000}{Techniques }}
\gT{440.25}{98.46}{\gS{R}{10}{c000000}{Techn. }}
\gT{330.00}{113.46}{\gS{R}{10}{c000000}{Technology }}
\gT{440.25}{113.46}{\gS{R}{10}{c000000}{Technol. }}
\gT{330.00}{128.46}{\gS{R}{10}{c000000}{Telecommunications }}
\gT{440.25}{128.46}{\gS{R}{10}{c000000}{Telecommun. }}
\gT{330.00}{143.46}{\gS{R}{10}{c000000}{Television }}
\gT{440.25}{143.46}{\gS{R}{10}{c000000}{Telev. }}
\gT{330.00}{158.46}{\gS{R}{10}{c000000}{Temperature }}
\gT{440.25}{158.46}{\gS{R}{10}{c000000}{Temp. }}
\gT{330.00}{173.46}{\gS{R}{10}{c000000}{Terrestrial }}
\gT{440.25}{173.46}{\gS{R}{10}{c000000}{Terr. }}
\gT{330.00}{188.46}{\gS{R}{10}{c000000}{Theorem }}
\gT{440.25}{188.46}{\gS{R}{10}{c000000}{Thm. }}
\gT{330.00}{203.46}{\gS{R}{10}{c000000}{Theoretical }}
\gT{440.25}{203.46}{\gS{R}{10}{c000000}{Theor. }}
\gT{330.00}{218.46}{\gS{R}{10}{c000000}{Transactions }}
\gT{440.25}{218.46}{\gS{R}{10}{c000000}{Trans. }}
\gT{330.00}{233.46}{\gS{R}{10}{c000000}{Translation }}
\gT{440.25}{233.46}{\gS{R}{10}{c000000}{Transl. }}
\gT{330.00}{248.46}{\gS{R}{10}{c000000}{Transmission }}
\gT{440.25}{248.46}{\gS{R}{10}{c000000}{Transmiss. }}
\gT{330.00}{263.46}{\gS{R}{10}{c000000}{Transportation }}
\gT{440.25}{263.46}{\gS{R}{10}{c000000}{Transp. }}
\gT{330.00}{278.46}{\gS{R}{10}{c000000}{Tutorials }}
\gT{440.25}{278.46}{\gS{R}{10}{c000000}{Tut. }}
\gT{330.00}{308.46}{\gS{RB}{10}{c000000}{U }}
\gT{330.00}{323.46}{\gS{R}{10}{c000000}{Ultrasonic }}
\gT{440.25}{323.46}{\gS{R}{10}{c000000}{Ultrason. }}
\gT{330.00}{338.46}{\gS{R}{10}{c000000}{University }}
\gT{440.25}{338.46}{\gS{R}{10}{c000000}{Univ. }}
\gT{330.00}{368.46}{\gS{RB}{10}{c000000}{V }}
\gT{330.00}{383.46}{\gS{R}{10}{c000000}{Vacuum }}
\gT{440.25}{383.46}{\gS{R}{10}{c000000}{Vac. }}
\gT{330.00}{398.46}{\gS{R}{10}{c000000}{Vehicular }}
\gT{440.25}{398.46}{\gS{R}{10}{c000000}{Veh. }}
\gT{330.00}{413.46}{\gS{R}{10}{c000000}{Vibration }}
\gT{440.25}{413.46}{\gS{R}{10}{c000000}{Vib. }}
\gT{330.00}{428.46}{\gS{R}{10}{c000000}{Vision }}
\gT{440.25}{428.46}{\gS{R}{10}{c000000}{Vis. }}
\gT{330.00}{443.46}{\gS{R}{10}{c000000}{Visual }}
\gT{440.25}{443.46}{\gS{R}{10}{c000000}{Vis. }}
\gT{330.00}{458.46}{\gS{R}{10}{c000000}{Volume }}
\gT{440.25}{458.46}{\gS{R}{10}{c000000}{Vol. }}
\gT{330.00}{488.46}{\gS{RB}{10}{c000000}{W }}
\gT{330.00}{503.46}{\gS{R}{10}{c000000}{Welding }}
\gT{440.25}{503.46}{\gS{R}{10}{c000000}{Weld. }}
\gT{330.00}{518.46}{\gS{R}{10}{c000000}{Working }}
\gT{440.25}{518.46}{\gS{R}{10}{c000000}{Work. }}
\end{guidepage}
\begin{guidepage}% page 28
\gF{cEFEFEF}{72.00}{118.80}{501.75}{136.80}
\gL{c000000}{1.00}{78.00}{132.50}{85.00}{132.50}
\gL{c000000}{1.00}{72.50}{118.00}{72.50}{713.00}
\gL{c000000}{1.00}{501.50}{118.00}{501.50}{713.00}
\gL{c000000}{1.00}{72.00}{118.50}{502.00}{118.50}
\gL{c000000}{1.00}{72.00}{136.50}{502.00}{136.50}
\gL{c000000}{1.00}{72.00}{154.50}{502.00}{154.50}
\gL{c000000}{1.00}{72.00}{172.50}{502.00}{172.50}
\gL{c000000}{1.00}{72.00}{190.50}{502.00}{190.50}
\gL{c000000}{1.00}{72.00}{208.50}{502.00}{208.50}
\gL{c000000}{1.00}{72.00}{226.50}{502.00}{226.50}
\gL{c000000}{1.00}{72.00}{244.50}{502.00}{244.50}
\gL{c000000}{1.00}{72.00}{262.50}{502.00}{262.50}
\gL{c000000}{1.00}{72.00}{280.50}{502.00}{280.50}
\gL{c000000}{1.00}{72.00}{298.50}{502.00}{298.50}
\gL{c000000}{1.00}{72.00}{316.50}{502.00}{316.50}
\gL{c000000}{1.00}{72.00}{334.50}{502.00}{334.50}
\gL{c000000}{1.00}{72.00}{352.50}{502.00}{352.50}
\gL{c000000}{1.00}{72.00}{370.50}{502.00}{370.50}
\gL{c000000}{1.00}{72.00}{388.50}{502.00}{388.50}
\gL{c000000}{1.00}{72.00}{406.50}{502.00}{406.50}
\gL{c000000}{1.00}{72.00}{424.50}{502.00}{424.50}
\gL{c000000}{1.00}{72.00}{442.50}{502.00}{442.50}
\gL{c000000}{1.00}{72.00}{460.50}{502.00}{460.50}
\gL{c000000}{1.00}{72.00}{478.50}{502.00}{478.50}
\gL{c000000}{1.00}{72.00}{496.50}{502.00}{496.50}
\gL{c000000}{1.00}{72.00}{514.50}{502.00}{514.50}
\gL{c000000}{1.00}{72.00}{532.50}{502.00}{532.50}
\gL{c000000}{1.00}{72.00}{550.50}{502.00}{550.50}
\gL{c000000}{1.00}{72.00}{568.50}{502.00}{568.50}
\gL{c000000}{1.00}{72.00}{586.50}{502.00}{586.50}
\gL{c000000}{1.00}{72.00}{604.50}{502.00}{604.50}
\gL{c000000}{1.00}{72.00}{622.50}{502.00}{622.50}
\gL{c000000}{1.00}{72.00}{640.50}{502.00}{640.50}
\gL{c000000}{1.00}{72.00}{658.50}{502.00}{658.50}
\gL{c000000}{1.00}{72.00}{676.50}{502.00}{676.50}
\gL{c000000}{1.00}{72.00}{694.50}{502.00}{694.50}
\gL{c000000}{1.00}{72.00}{712.50}{502.00}{712.50}
\gT{532.01}{43.47}{\gS{R}{8}{c000000}{28 }}
\gT{85.50}{84.36}{\gS{RB}{14}{c000000}{B.}\hspace*{5.17bp}\gS{R}{14}{c000000}{List of Publishers }}
\gT{72.00}{103.63}{\gS{RB}{10}{c000000}{NOTE:}\gS{R}{10}{c000000}{ Each publisher’s complete title is shown, but do not use what is in parentheses. This is just to show the }}
\gT{86.25}{115.63}{\gS{R}{10}{c000000}{entire title. }}
\gT{78.00}{131.01}{\gS{RB}{10}{c000000}{A}}
\gT{78.00}{149.01}{\gS{R}{10}{c000000}{Abelard-Schuman (Ltd.), New York, NY, USA }}
\gT{78.00}{167.01}{\gS{R}{10}{c000000}{Abingdon (Press), Nashville, TN, USA }}
\gT{78.00}{185.01}{\gS{R}{10}{c000000}{Ablex (Publishing Corp.), Norwood, NJ, USA }}
\gT{78.00}{203.01}{\gS{R}{10}{c000000}{Harry N. Abrams (Inc.), New York, NY, USA }}
\gT{78.00}{221.01}{\gS{R}{10}{c000000}{Academic (Press Inc.), New York, NY, USA; Orlando, FL, USA; San Diego, CA, USA }}
\gT{78.00}{239.01}{\gS{R}{10}{c000000}{Ace Books (Inc.), New York, NY, USA }}
\gT{78.00}{257.01}{\gS{R}{10}{c000000}{Adam Hilger, New York, NY, USA }}
\gT{78.00}{275.01}{\gS{R}{10}{c000000}{Addison-Wesley, Reading, MA, USA }}
\gT{78.00}{293.01}{\gS{R}{10}{c000000}{Aegean Park (Press), Laguna Hills, CA, USA }}
\gT{78.00}{311.01}{\gS{R}{10}{c000000}{Aiglon (Press, Ltd.), London, U.K. }}
\gT{78.00}{329.01}{\gS{R}{10}{c000000}{AIP (American Inst. Phys.), New York, NY, USA }}
\gT{78.00}{347.01}{\gS{R}{10}{c000000}{Alden (\& Co., Ltd.), Oxford, U.K. }}
\gT{78.00}{365.01}{\gS{R}{10}{c000000}{Aldine (Publishing Co.), London, U.K. }}
\gT{78.00}{383.01}{\gS{R}{10}{c000000}{Allen \& Sons (Ltd.), London, U.K. }}
\gT{78.00}{401.01}{\gS{R}{10}{c000000}{G. (George) Allen (\& Unwin Ltd.), London, U.K. }}
\gT{78.00}{419.01}{\gS{R}{10}{c000000}{H. (Howard) Allen (Inc.), Cleveland, OH, USA }}
\gT{78.00}{437.01}{\gS{R}{10}{c000000}{Allenson (\& Co., Ltd.), London, U.K. }}
\gT{78.00}{455.01}{\gS{R}{10}{c000000}{Allyn \& Bacon (Inc.), Boston, MA, USA }}
\gT{78.00}{473.01}{\gS{R}{10}{c000000}{American Book (Co.), New York, NY, USA }}
\gT{78.00}{491.01}{\gS{R}{10}{c000000}{American Council on Public Affairs, Washington, DC, USA }}
\gT{78.00}{509.01}{\gS{R}{10}{c000000}{American Elsevier, New York, NY, USA }}
\gT{78.00}{527.01}{\gS{R}{10}{c000000}{American Geographical Soc., New York, NY, USA }}
\gT{78.00}{545.01}{\gS{R}{10}{c000000}{American Library Assoc., Chicago, IL, USA }}
\gT{78.00}{563.01}{\gS{R}{10}{c000000}{American Management Assoc. (Inc.), New York, NY, USA }}
\gT{78.00}{581.01}{\gS{R}{10}{c000000}{American Philosophical Soc., Philadelphia, PA, USA }}
\gT{78.00}{599.01}{\gS{R}{10}{c000000}{American Psychological Assoc. (Inc.), Washington, DC, USA }}
\gT{78.00}{617.01}{\gS{R}{10}{c000000}{American Technical Soc., Chicago, IL, USA }}
\gT{78.00}{635.01}{\gS{R}{10}{c000000}{Amphoto Books (American Photographic Book Publishing Co. Inc.), New York, NY, USA }}
\gT{78.00}{653.01}{\gS{R}{10}{c000000}{Antioch (Press), Yellow Springs, OH, USA }}
\gT{78.00}{671.01}{\gS{R}{10}{c000000}{Appleton (-Century-Crofts), New York, NY, USA }}
\gT{78.00}{689.01}{\gS{R}{10}{c000000}{Architectural Book (Publishing Co., Inc.), New York, NY, USA }}
\gT{78.00}{707.01}{\gS{R}{10}{c000000}{Arco (Publishing Co., Inc.), New York, NY, USA }}
\end{guidepage}
\begin{guidepage}% page 29
\gF{cEFEFEF}{72.00}{396.00}{501.75}{414.00}
\gL{c000000}{1.00}{78.00}{409.50}{85.00}{409.50}
\gL{c000000}{1.00}{72.50}{72.00}{72.50}{720.00}
\gL{c000000}{1.00}{501.50}{72.00}{501.50}{720.00}
\gL{c000000}{1.00}{72.00}{72.50}{502.00}{72.50}
\gL{c000000}{1.00}{72.00}{90.50}{502.00}{90.50}
\gL{c000000}{1.00}{72.00}{108.50}{502.00}{108.50}
\gL{c000000}{1.00}{72.00}{126.50}{502.00}{126.50}
\gL{c000000}{1.00}{72.00}{144.50}{502.00}{144.50}
\gL{c000000}{1.00}{72.00}{162.50}{502.00}{162.50}
\gL{c000000}{1.00}{72.00}{180.50}{502.00}{180.50}
\gL{c000000}{1.00}{72.00}{198.50}{502.00}{198.50}
\gL{c000000}{1.00}{72.00}{216.50}{502.00}{216.50}
\gL{c000000}{1.00}{72.00}{234.50}{502.00}{234.50}
\gL{c000000}{1.00}{72.00}{252.50}{502.00}{252.50}
\gL{c000000}{1.00}{72.00}{270.50}{502.00}{270.50}
\gL{c000000}{1.00}{72.00}{288.50}{502.00}{288.50}
\gL{c000000}{1.00}{72.00}{306.50}{502.00}{306.50}
\gL{c000000}{1.00}{72.00}{324.50}{502.00}{324.50}
\gL{c000000}{1.00}{72.00}{342.50}{502.00}{342.50}
\gL{c000000}{1.00}{72.00}{360.50}{502.00}{360.50}
\gL{c000000}{1.00}{72.00}{378.50}{502.00}{378.50}
\gL{c000000}{1.00}{72.00}{396.50}{502.00}{396.50}
\gL{c000000}{1.00}{72.00}{414.50}{502.00}{414.50}
\gL{c000000}{1.00}{72.00}{432.50}{502.00}{432.50}
\gL{c000000}{1.00}{72.00}{450.50}{502.00}{450.50}
\gL{c000000}{1.00}{72.00}{468.50}{502.00}{468.50}
\gL{c000000}{1.00}{72.00}{486.50}{502.00}{486.50}
\gL{c000000}{1.00}{72.00}{504.50}{502.00}{504.50}
\gL{c000000}{1.00}{72.00}{522.50}{502.00}{522.50}
\gL{c000000}{1.00}{72.00}{540.50}{502.00}{540.50}
\gL{c000000}{1.00}{72.00}{558.50}{502.00}{558.50}
\gL{c000000}{1.00}{72.00}{576.50}{502.00}{576.50}
\gL{c000000}{1.00}{72.00}{594.50}{502.00}{594.50}
\gL{c000000}{1.00}{72.00}{612.50}{502.00}{612.50}
\gL{c000000}{1.00}{72.00}{630.50}{502.00}{630.50}
\gL{c000000}{1.00}{72.00}{648.50}{502.00}{648.50}
\gL{c000000}{1.00}{72.00}{666.50}{502.00}{666.50}
\gL{c000000}{1.00}{72.00}{684.50}{502.00}{684.50}
\gL{c000000}{1.00}{72.00}{702.50}{502.00}{702.50}
\gL{c000000}{1.00}{72.00}{720.50}{502.00}{720.50}
\gT{532.01}{43.47}{\gS{R}{8}{c000000}{29 }}
\gT{78.00}{84.21}{\gS{R}{10}{c000000}{Argus (Books, Inc.), Mohegan Lake, NY, USA }}
\gT{78.00}{102.21}{\gS{R}{10}{c000000}{E. (Edward) Arnold (\& Co.), London, U.K. }}
\gT{78.00}{120.21}{\gS{R}{10}{c000000}{E. J. Arnold (\& Sons, Ltd.), London, U.K. }}
\gT{78.00}{138.21}{\gS{R}{10}{c000000}{Artech House, Norwood, MA, USA }}
\gT{78.00}{156.21}{\gS{R}{10}{c000000}{ASME (American Soc. Mechanical Eng.), New York, NY, USA }}
\gT{78.00}{174.21}{\gS{R}{10}{c000000}{Aspen (Publishing, Inc.), Gaithersburg, MD, USA }}
\gT{78.00}{192.21}{\gS{R}{10}{c000000}{Associated Press, New York, NY, USA }}
\gT{78.00}{210.21}{\gS{R}{10}{c000000}{Athena Scientific, Belmont, MA, USA }}
\gT{78.00}{228.21}{\gS{R}{10}{c000000}{Atheneum Press, Boston, MA, USA }}
\gT{78.00}{246.21}{\gS{R}{10}{c000000}{Atheneum Press (Ltd.), London, U.K. }}
\gT{78.00}{264.20}{\gS{R}{10}{c000000}{Atheneum (Publishers), New York, NY, USA }}
\gT{78.00}{282.20}{\gS{R}{10}{c000000}{Atherton (Press), New York, NY, USA }}
\gT{78.00}{300.20}{\gS{R}{10}{c000000}{Atlantic Monthly (Press), Boston, MA, USA }}
\gT{78.00}{318.20}{\gS{R}{10}{c000000}{Auerbach (Publishing), New York, NY, USA }}
\gT{78.00}{336.20}{\gS{R}{10}{c000000}{Augsburg (Publishing House), Minneapolis, MN, USA }}
\gT{78.00}{354.20}{\gS{R}{10}{c000000}{(J. J.) Augustin (Inc.), Locust Valley, NY, USA }}
\gT{78.00}{372.20}{\gS{R}{10}{c000000}{(N.W.) Ayer (\& Sons, Inc.), Philadelphia, PA, USA }}
\gT{78.00}{408.20}{\gS{RB}{10}{c000000}{B}}
\gT{78.00}{426.20}{\gS{R}{10}{c000000}{Babson (Statistical Organization, Inc.), Wellesley Hills, MA, USA }}
\gT{78.00}{444.20}{\gS{R}{10}{c000000}{(Karl) Baedeker, Leipzig, Germany }}
\gT{78.00}{462.20}{\gS{R}{10}{c000000}{Bailliere (Tindall \& Cox, Ltd.), London, U.K. }}
\gT{78.00}{480.20}{\gS{R}{10}{c000000}{W. H. Baker (Co.), Boston, MA, USA }}
\gT{78.00}{498.20}{\gS{R}{10}{c000000}{Ballentine (Books, Inc.), New York, NY, USA }}
\gT{78.00}{516.21}{\gS{R}{10}{c000000}{Bankers (Publishing Co.), Boston, MA, USA }}
\gT{78.00}{534.21}{\gS{R}{10}{c000000}{Banta (George Banta Publishing Co. - The Collegiate Press), Menasha, WI, USA }}
\gT{78.00}{552.21}{\gS{R}{10}{c000000}{A. S. Barnes (\& Co., Inc.), Cranbury, NJ, USA }}
\gT{78.00}{570.21}{\gS{R}{10}{c000000}{Barnes \& Noble (Inc.), New York, NY, USA }}
\gT{78.00}{588.21}{\gS{R}{10}{c000000}{Barren (Educational Series, Inc.), Woodbury, NY, USA }}
\gT{78.00}{606.21}{\gS{R}{10}{c000000}{(M.) Barrows (\& Co., Inc.), New York, NY, USA }}
\gT{78.00}{624.21}{\gS{R}{10}{c000000}{(Johann Ambrosius) Barth, Munich, Germany }}
\gT{78.00}{642.21}{\gS{R}{10}{c000000}{Basic Books (Inc.), New York, NY, USA }}
\gT{78.00}{660.21}{\gS{R}{10}{c000000}{Battelle (Press), Columbus, OH, USA }}
\gT{78.00}{678.21}{\gS{R}{10}{c000000}{Baylor Univ. Press, Waco, TX, USA }}
\gT{78.00}{696.21}{\gS{R}{10}{c000000}{Beacon (Press), Boston, MA, USA }}
\gT{78.00}{714.21}{\gS{R}{10}{c000000}{G. Bell (\& Sons, Ltd.), London, U.K. }}
\end{guidepage}
\begin{guidepage}% page 30
\gL{c000000}{1.00}{72.50}{72.00}{72.50}{720.00}
\gL{c000000}{1.00}{501.50}{72.00}{501.50}{720.00}
\gL{c000000}{1.00}{72.00}{72.50}{502.00}{72.50}
\gL{c000000}{1.00}{72.00}{90.50}{502.00}{90.50}
\gL{c000000}{1.00}{72.00}{108.50}{502.00}{108.50}
\gL{c000000}{1.00}{72.00}{126.50}{502.00}{126.50}
\gL{c000000}{1.00}{72.00}{144.50}{502.00}{144.50}
\gL{c000000}{1.00}{72.00}{162.50}{502.00}{162.50}
\gL{c000000}{1.00}{72.00}{180.50}{502.00}{180.50}
\gL{c000000}{1.00}{72.00}{198.50}{502.00}{198.50}
\gL{c000000}{1.00}{72.00}{216.50}{502.00}{216.50}
\gL{c000000}{1.00}{72.00}{234.50}{502.00}{234.50}
\gL{c000000}{1.00}{72.00}{252.50}{502.00}{252.50}
\gL{c000000}{1.00}{72.00}{270.50}{502.00}{270.50}
\gL{c000000}{1.00}{72.00}{288.50}{502.00}{288.50}
\gL{c000000}{1.00}{72.00}{306.50}{502.00}{306.50}
\gL{c000000}{1.00}{72.00}{324.50}{502.00}{324.50}
\gL{c000000}{1.00}{72.00}{342.50}{502.00}{342.50}
\gL{c000000}{1.00}{72.00}{360.50}{502.00}{360.50}
\gL{c000000}{1.00}{72.00}{378.50}{502.00}{378.50}
\gL{c000000}{1.00}{72.00}{396.50}{502.00}{396.50}
\gL{c000000}{1.00}{72.00}{414.50}{502.00}{414.50}
\gL{c000000}{1.00}{72.00}{432.50}{502.00}{432.50}
\gL{c000000}{1.00}{72.00}{450.50}{502.00}{450.50}
\gL{c000000}{1.00}{72.00}{468.50}{502.00}{468.50}
\gL{c000000}{1.00}{72.00}{486.50}{502.00}{486.50}
\gL{c000000}{1.00}{72.00}{504.50}{502.00}{504.50}
\gL{c000000}{1.00}{72.00}{522.50}{502.00}{522.50}
\gL{c000000}{1.00}{72.00}{540.50}{502.00}{540.50}
\gL{c000000}{1.00}{72.00}{558.50}{502.00}{558.50}
\gL{c000000}{1.00}{72.00}{576.50}{502.00}{576.50}
\gL{c000000}{1.00}{72.00}{594.50}{502.00}{594.50}
\gL{c000000}{1.00}{72.00}{612.50}{502.00}{612.50}
\gL{c000000}{1.00}{72.00}{630.50}{502.00}{630.50}
\gL{c000000}{1.00}{72.00}{648.50}{502.00}{648.50}
\gL{c000000}{1.00}{72.00}{666.50}{502.00}{666.50}
\gL{c000000}{1.00}{72.00}{684.50}{502.00}{684.50}
\gL{c000000}{1.00}{72.00}{702.50}{502.00}{702.50}
\gL{c000000}{1.00}{72.00}{720.50}{502.00}{720.50}
\gT{532.01}{43.47}{\gS{R}{8}{c000000}{30 }}
\gT{78.00}{84.21}{\gS{R}{10}{c000000}{J. P. Bell (Co.), Lynchburg, VA, USA }}
\gT{78.00}{102.21}{\gS{R}{10}{c000000}{(Matthew) Bender (\& Co., Inc.), New York, NY, USA }}
\gT{78.00}{120.21}{\gS{R}{10}{c000000}{W. A. Benjamin (Inc.), New York, NY, USA }}
\gT{78.00}{138.21}{\gS{R}{10}{c000000}{Benjamin Cummings (Publishing Co.), Redwood City, CA, USA }}
\gT{78.00}{156.21}{\gS{R}{10}{c000000}{(Ernest) Benn (Ltd.), London, U.K. }}
\gT{78.00}{174.21}{\gS{R}{10}{c000000}{(Charles A.) Bennett (Co.), Peoria, IL, USA }}
\gT{78.00}{192.21}{\gS{R}{10}{c000000}{Benziger (Brothers, Inc.), New York, NY, USA }}
\gT{78.00}{210.21}{\gS{R}{10}{c000000}{(Librairie Polytechnique Charles) Beranger, Paris, France }}
\gT{78.00}{228.21}{\gS{R}{10}{c000000}{Berlinische Verlagsanstalt, Berlin, Germany }}
\gT{78.00}{246.21}{\gS{R}{10}{c000000}{Beuth Vertrieb, G.m.b.H., Berlin, Germany }}
\gT{78.00}{264.20}{\gS{R}{10}{c000000}{Birkhaüser, Cambridge, MA, USA }}
\gT{78.00}{282.20}{\gS{R}{10}{c000000}{A. \& C. Black (Ltd.), London, U.K. }}
\gT{78.00}{300.20}{\gS{R}{10}{c000000}{(Walter J.) Black (Inc.), Roslyn, NY, USA }}
\gT{78.00}{318.20}{\gS{R}{10}{c000000}{Blackie (\& Son, Ltd.), Glasgow, U.K. }}
\gT{78.00}{336.20}{\gS{R}{10}{c000000}{(Basil) Blackwell (\& Mott, Ltd.), Oxford, U.K. }}
\gT{78.00}{354.20}{\gS{R}{10}{c000000}{(William) Blackwood (\& Sons, Ltd.), Edinburgh, U.K. }}
\gT{78.00}{372.20}{\gS{R}{10}{c000000}{Blaisdell, London, U.K. }}
\gT{78.00}{390.20}{\gS{R}{10}{c000000}{Blandford (Press. Ltd.), London, U.K. }}
\gT{78.00}{408.20}{\gS{R}{10}{c000000}{Blue Ribbon (Books, Inc.) New York, NY, USA }}
\gT{78.00}{426.20}{\gS{R}{10}{c000000}{(Clark) Boardman (Co., Ltd.) New York, NY, USA }}
\gT{78.00}{444.20}{\gS{R}{10}{c000000}{Bobbs-Merrill (Co., Inc.), Indianapolis, IN, USA }}
\gT{78.00}{462.20}{\gS{R}{10}{c000000}{Bonniers (Inc.), New York, NY, USA }}
\gT{78.00}{480.20}{\gS{R}{10}{c000000}{A. (Albert) Bonniers Forlag, Stockholm, Sweden }}
\gT{78.00}{498.20}{\gS{R}{10}{c000000}{(Thomas) Bouregy (\& Co.), New York, NY, USA }}
\gT{78.00}{516.21}{\gS{R}{10}{c000000}{(R. R.) Bowker, New Providence, NJ, USA }}
\gT{78.00}{534.21}{\gS{R}{10}{c000000}{Boyd \& Fraser (Publishing Co.), Boston, MA, USA }}
\gT{78.00}{552.21}{\gS{R}{10}{c000000}{(Charles T.) Branford (Co.), Newton Center, MA, USA }}
\gT{78.00}{570.21}{\gS{R}{10}{c000000}{(George) Braziller (Inc.), New York, NY, USA }}
\gT{78.00}{588.21}{\gS{R}{10}{c000000}{(E. J.) Brill, Leiden, The Netherlands }}
\gT{78.00}{606.21}{\gS{R}{10}{c000000}{British Book Centre (Inc.), New York, NY, USA }}
\gT{78.00}{624.21}{\gS{R}{10}{c000000}{(F. A.) Brockhaus, Wiesbaden, Germany }}
\gT{78.00}{642.21}{\gS{R}{10}{c000000}{Brookings Institution, Washington, DC, USA }}
\gT{78.00}{660.21}{\gS{R}{10}{c000000}{Brooks/Cole (Publishing Co.), Pacific Grove, CA, USA }}
\gT{78.00}{678.21}{\gS{R}{10}{c000000}{William C. Brown (Co.), Dubuque, IA, USA }}
\gT{78.00}{696.21}{\gS{R}{10}{c000000}{Bruce (Publishing Co.), Milwaukee, WI, USA }}
\gT{78.00}{714.21}{\gS{R}{10}{c000000}{Burns (Oates \& Washbourne, Ltd.), London, U.K. }}
\end{guidepage}
\begin{guidepage}% page 31
\gF{cEFEFEF}{72.00}{108.00}{501.75}{126.00}
\gL{c000000}{1.00}{78.00}{121.50}{85.00}{121.50}
\gL{c000000}{1.00}{72.50}{72.00}{72.50}{720.00}
\gL{c000000}{1.00}{501.50}{72.00}{501.50}{720.00}
\gL{c000000}{1.00}{72.00}{72.50}{502.00}{72.50}
\gL{c000000}{1.00}{72.00}{90.50}{502.00}{90.50}
\gL{c000000}{1.00}{72.00}{108.50}{502.00}{108.50}
\gL{c000000}{1.00}{72.00}{126.50}{502.00}{126.50}
\gL{c000000}{1.00}{72.00}{144.50}{502.00}{144.50}
\gL{c000000}{1.00}{72.00}{162.50}{502.00}{162.50}
\gL{c000000}{1.00}{72.00}{180.50}{502.00}{180.50}
\gL{c000000}{1.00}{72.00}{198.50}{502.00}{198.50}
\gL{c000000}{1.00}{72.00}{216.50}{502.00}{216.50}
\gL{c000000}{1.00}{72.00}{234.50}{502.00}{234.50}
\gL{c000000}{1.00}{72.00}{252.50}{502.00}{252.50}
\gL{c000000}{1.00}{72.00}{270.50}{502.00}{270.50}
\gL{c000000}{1.00}{72.00}{288.50}{502.00}{288.50}
\gL{c000000}{1.00}{72.00}{306.50}{502.00}{306.50}
\gL{c000000}{1.00}{72.00}{324.50}{502.00}{324.50}
\gL{c000000}{1.00}{72.00}{342.50}{502.00}{342.50}
\gL{c000000}{1.00}{72.00}{360.50}{502.00}{360.50}
\gL{c000000}{1.00}{72.00}{378.50}{502.00}{378.50}
\gL{c000000}{1.00}{72.00}{396.50}{502.00}{396.50}
\gL{c000000}{1.00}{72.00}{414.50}{502.00}{414.50}
\gL{c000000}{1.00}{72.00}{432.50}{502.00}{432.50}
\gL{c000000}{1.00}{72.00}{450.50}{502.00}{450.50}
\gL{c000000}{1.00}{72.00}{468.50}{502.00}{468.50}
\gL{c000000}{1.00}{72.00}{486.50}{502.00}{486.50}
\gL{c000000}{1.00}{72.00}{504.50}{502.00}{504.50}
\gL{c000000}{1.00}{72.00}{522.50}{502.00}{522.50}
\gL{c000000}{1.00}{72.00}{540.50}{502.00}{540.50}
\gL{c000000}{1.00}{72.00}{558.50}{502.00}{558.50}
\gL{c000000}{1.00}{72.00}{576.50}{502.00}{576.50}
\gL{c000000}{1.00}{72.00}{594.50}{502.00}{594.50}
\gL{c000000}{1.00}{72.00}{612.50}{502.00}{612.50}
\gL{c000000}{1.00}{72.00}{630.50}{502.00}{630.50}
\gL{c000000}{1.00}{72.00}{648.50}{502.00}{648.50}
\gL{c000000}{1.00}{72.00}{666.50}{502.00}{666.50}
\gL{c000000}{1.00}{72.00}{684.50}{502.00}{684.50}
\gL{c000000}{1.00}{72.00}{702.50}{502.00}{702.50}
\gL{c000000}{1.00}{72.00}{720.50}{502.00}{720.50}
\gT{532.01}{43.47}{\gS{R}{8}{c000000}{31 }}
\gT{78.00}{84.21}{\gS{R}{10}{c000000}{Butterworth (\& Co., Ltd.), London, U.K. }}
\gT{78.00}{120.21}{\gS{RB}{10}{c000000}{C}}
\gT{78.00}{138.21}{\gS{R}{10}{c000000}{Cambridge Univ. Press, Cambridge, U.K. }}
\gT{78.00}{156.21}{\gS{R}{10}{c000000}{(Jonathon) Cape (Ltd.), London, U.K. }}
\gT{78.00}{174.21}{\gS{R}{10}{c000000}{Carnegie Endowment (for International Peace), New York, NY, USA }}
\gT{78.00}{192.21}{\gS{R}{10}{c000000}{Carnegie Foundation (for the Advancement of Teaching), New York, NY, USA }}
\gT{78.00}{210.21}{\gS{R}{10}{c000000}{Carnegie Press, Pittsburgh, PA, USA }}
\gT{78.00}{228.21}{\gS{R}{10}{c000000}{Cassell (\& Co., Ltd.), London, U.K. }}
\gT{78.00}{246.21}{\gS{R}{10}{c000000}{Catholic Univ. of America Press, Washington, DC, USA }}
\gT{78.00}{264.20}{\gS{R}{10}{c000000}{Caxton (Printers, Ltd.), Caldwell, ID, USA }}
\gT{78.00}{282.20}{\gS{R}{10}{c000000}{Century House, Watkins Glen, NY, USA }}
\gT{78.00}{300.20}{\gS{R}{10}{c000000}{(W. \& R.) Chambers (Ltd.), Edinburgh, U.K. }}
\gT{78.00}{318.20}{\gS{R}{10}{c000000}{Chandler (Publishing Co.), San Francisco, CA, USA }}
\gT{78.00}{336.20}{\gS{R}{10}{c000000}{Channel Press, New York, NY, USA }}
\gT{78.00}{354.20}{\gS{R}{10}{c000000}{Chapman \& Grimes (Inc.), Boston, MA, USA }}
\gT{78.00}{372.20}{\gS{R}{10}{c000000}{Chapman \& Hall (Ltd.), London, U.K. }}
\gT{78.00}{390.20}{\gS{R}{10}{c000000}{Chatto \& Windus (Ltd.), London, U.K. }}
\gT{78.00}{408.20}{\gS{R}{10}{c000000}{Chelsea (Publishing Co.), New York, NY, USA }}
\gT{78.00}{426.20}{\gS{R}{10}{c000000}{Chemical Publishing (Co.), New York, NY, USA }}
\gT{78.00}{444.20}{\gS{R}{10}{c000000}{Chilton (Co.), Philadelphia, PA, USA }}
\gT{78.00}{462.20}{\gS{R}{10}{c000000}{Christian Science Publishing (Society), Boston, MA, USA }}
\gT{78.00}{480.20}{\gS{R}{10}{c000000}{Citadel (Press), New York, NY, USA }}
\gT{78.00}{498.20}{\gS{R}{10}{c000000}{A. H. Clark (Co.), Glendale, CA, USA }}
\gT{78.00}{516.21}{\gS{R}{10}{c000000}{Clarendon (Press), Oxford, U.K. }}
\gT{78.00}{534.21}{\gS{R}{10}{c000000}{Clark Univ. Press, Worcester, MA, USA }}
\gT{78.00}{552.21}{\gS{R}{10}{c000000}{Collegiate (Press), Ames, IA, USA }}
\gT{78.00}{570.21}{\gS{R}{10}{c000000}{(P. F.) Collier (\& Son Corp.), New York, NY, USA }}
\gT{78.00}{588.21}{\gS{R}{10}{c000000}{Collier-Macmillian (Ltd.), London, U.K. }}
\gT{78.00}{606.21}{\gS{R}{10}{c000000}{Columbia Univ. Press, New York, NY, USA }}
\gT{78.00}{624.21}{\gS{R}{10}{c000000}{Combat (Forces Press), Washington, DC, USA }}
\gT{78.00}{642.21}{\gS{R}{10}{c000000}{Computer Science (Press), Rockville, MD, USA }}
\gT{78.00}{660.21}{\gS{R}{10}{c000000}{Comstock (Publishing Associates), Ithaca, NY, USA }}
\gT{78.00}{678.21}{\gS{R}{10}{c000000}{Concordia (Publishing House), St. Louis, MO, USA }}
\gT{78.00}{696.21}{\gS{R}{10}{c000000}{Concrete (Publications, Ltd.), London, U.K. }}
\gT{78.00}{714.21}{\gS{R}{10}{c000000}{Cornell Univ. Press, Ithaca, NY, USA }}
\end{guidepage}
\begin{guidepage}% page 32
\gF{cEFEFEF}{72.00}{198.00}{501.75}{216.00}
\gF{cEFEFEF}{72.00}{630.00}{501.75}{648.00}
\gL{c000000}{1.00}{78.00}{211.50}{85.00}{211.50}
\gL{c000000}{1.00}{78.00}{643.50}{85.00}{643.50}
\gL{c000000}{1.00}{72.50}{72.00}{72.50}{720.00}
\gL{c000000}{1.00}{501.50}{72.00}{501.50}{720.00}
\gL{c000000}{1.00}{72.00}{72.50}{502.00}{72.50}
\gL{c000000}{1.00}{72.00}{90.50}{502.00}{90.50}
\gL{c000000}{1.00}{72.00}{108.50}{502.00}{108.50}
\gL{c000000}{1.00}{72.00}{126.50}{502.00}{126.50}
\gL{c000000}{1.00}{72.00}{144.50}{502.00}{144.50}
\gL{c000000}{1.00}{72.00}{162.50}{502.00}{162.50}
\gL{c000000}{1.00}{72.00}{180.50}{502.00}{180.50}
\gL{c000000}{1.00}{72.00}{198.50}{502.00}{198.50}
\gL{c000000}{1.00}{72.00}{216.50}{502.00}{216.50}
\gL{c000000}{1.00}{72.00}{234.50}{502.00}{234.50}
\gL{c000000}{1.00}{72.00}{252.50}{502.00}{252.50}
\gL{c000000}{1.00}{72.00}{270.50}{502.00}{270.50}
\gL{c000000}{1.00}{72.00}{288.50}{502.00}{288.50}
\gL{c000000}{1.00}{72.00}{306.50}{502.00}{306.50}
\gL{c000000}{1.00}{72.00}{324.50}{502.00}{324.50}
\gL{c000000}{1.00}{72.00}{342.50}{502.00}{342.50}
\gL{c000000}{1.00}{72.00}{360.50}{502.00}{360.50}
\gL{c000000}{1.00}{72.00}{378.50}{502.00}{378.50}
\gL{c000000}{1.00}{72.00}{396.50}{502.00}{396.50}
\gL{c000000}{1.00}{72.00}{414.50}{502.00}{414.50}
\gL{c000000}{1.00}{72.00}{432.50}{502.00}{432.50}
\gL{c000000}{1.00}{72.00}{450.50}{502.00}{450.50}
\gL{c000000}{1.00}{72.00}{468.50}{502.00}{468.50}
\gL{c000000}{1.00}{72.00}{486.50}{502.00}{486.50}
\gL{c000000}{1.00}{72.00}{504.50}{502.00}{504.50}
\gL{c000000}{1.00}{72.00}{522.50}{502.00}{522.50}
\gL{c000000}{1.00}{72.00}{540.50}{502.00}{540.50}
\gL{c000000}{1.00}{72.00}{558.50}{502.00}{558.50}
\gL{c000000}{1.00}{72.00}{576.50}{502.00}{576.50}
\gL{c000000}{1.00}{72.00}{594.50}{502.00}{594.50}
\gL{c000000}{1.00}{72.00}{612.50}{502.00}{612.50}
\gL{c000000}{1.00}{72.00}{630.50}{502.00}{630.50}
\gL{c000000}{1.00}{72.00}{648.50}{502.00}{648.50}
\gL{c000000}{1.00}{72.00}{666.50}{502.00}{666.50}
\gL{c000000}{1.00}{72.00}{684.50}{502.00}{684.50}
\gL{c000000}{1.00}{72.00}{702.50}{502.00}{702.50}
\gL{c000000}{1.00}{72.00}{720.50}{502.00}{720.50}
\gT{532.01}{43.47}{\gS{R}{8}{c000000}{32 }}
\gT{78.00}{84.21}{\gS{R}{10}{c000000}{Coward-McCann (Inc.), New York, NY, USA }}
\gT{78.00}{102.21}{\gS{R}{10}{c000000}{CRC Press (Inc.), Boca Raton, FL, USA; (also Cleveland, OH, USA) }}
\gT{78.00}{120.21}{\gS{R}{10}{c000000}{Criterion (Press, Inc.), Torrance, CA, USA }}
\gT{78.00}{138.21}{\gS{R}{10}{c000000}{(Thomas Y.) Crowell (Co.), New York, NY, USA }}
\gT{78.00}{156.21}{\gS{R}{10}{c000000}{Crower-Collier-Macmillian (Inc.), New York, NY, USA }}
\gT{78.00}{174.21}{\gS{R}{10}{c000000}{Crown (Publishers, Inc.), New York, NY, USA }}
\gT{78.00}{210.21}{\gS{RB}{10}{c000000}{D}}
\gT{78.00}{228.21}{\gS{R}{10}{c000000}{(Peter) Davies (Ltd.), London, U.K. }}
\gT{78.00}{246.21}{\gS{R}{10}{c000000}{(F. A.) Davis (Co.), Philadelphia, PA, USA }}
\gT{78.00}{264.20}{\gS{R}{10}{c000000}{John Day (Co., Inc.), New York, NY, USA }}
\gT{78.00}{282.20}{\gS{R}{10}{c000000}{(John) De Graff (Inc.), Tuckahoe, NY, USA }}
\gT{78.00}{300.20}{\gS{R}{10}{c000000}{Delmar (Publishers, Inc.), Albany, NY, USA }}
\gT{78.00}{318.20}{\gS{R}{10}{c000000}{(T. S.) Denison (\& Co., Inc.), Minneapolis, MN, USA }}
\gT{78.00}{336.20}{\gS{R}{10}{c000000}{(J. M.) Dent (\& Sons, Ltd.), London, U.K. }}
\gT{78.00}{354.20}{\gS{R}{10}{c000000}{Devin-Adair (Co.), New York, NY, USA }}
\gT{78.00}{372.20}{\gS{R}{10}{c000000}{Dial (Press, Inc.), New York, NY, USA }}
\gT{78.00}{390.20}{\gS{R}{10}{c000000}{Digital (Press, Inc.), Bedford, MA, USA }}
\gT{78.00}{408.20}{\gS{R}{10}{c000000}{Dodd, Mead (\& Co., Inc.), New York, NY, USA }}
\gT{78.00}{426.20}{\gS{R}{10}{c000000}{Doubleday (\& Co., Inc.), New York, NY, USA }}
\gT{78.00}{444.20}{\gS{R}{10}{c000000}{Dover (Publications, Inc.), New York, NY, USA }}
\gT{78.00}{462.20}{\gS{R}{10}{c000000}{Dow Jones-Irwin (Inc.), Homewood, IL, USA }}
\gT{78.00}{480.20}{\gS{R}{10}{c000000}{(Frederick J.) Drake (\& Co., Inc.), Chicago, IL, USA }}
\gT{78.00}{498.20}{\gS{R}{10}{c000000}{Dryad (Press), Leicester, U.K. }}
\gT{78.00}{516.21}{\gS{R}{10}{c000000}{(Gerald) Duckworth (\& Co., Ltd.), London, U.K. }}
\gT{78.00}{534.21}{\gS{R}{10}{c000000}{Duke Univ. Press, Durham, NC, USA }}
\gT{78.00}{552.21}{\gS{R}{10}{c000000}{Duncker \& Humblet, Berlin, Germany }}
\gT{78.00}{570.21}{\gS{R}{10}{c000000}{Dunod, Paris, France }}
\gT{78.00}{588.21}{\gS{R}{10}{c000000}{P. Dupont, Paris, France }}
\gT{78.00}{606.21}{\gS{R}{10}{c000000}{(E. P.) Dutton (\& Co., Inc.), New York, NY, USA }}
\gT{78.00}{642.21}{\gS{RB}{10}{c000000}{E}}
\gT{78.00}{660.21}{\gS{R}{10}{c000000}{Edinburgh House (Press), London, U.K. }}
\gT{78.00}{678.21}{\gS{R}{10}{c000000}{Educational Technology (Publications), Englewood Cliffs, NJ, USA }}
\gT{78.00}{696.21}{\gS{R}{10}{c000000}{(J. W.) Edwards (Publisher, Inc.), Ann Arbor, MI, USA }}
\gT{78.00}{714.21}{\gS{R}{10}{c000000}{Eldon (Press, Ltd.), London, U.K. }}
\end{guidepage}
\begin{guidepage}% page 33
\gF{cEFEFEF}{72.00}{180.00}{501.75}{198.00}
\gF{cEFEFEF}{72.00}{486.00}{501.75}{504.00}
\gL{c000000}{1.00}{78.00}{193.50}{84.00}{193.50}
\gL{c000000}{1.00}{78.00}{499.50}{86.00}{499.50}
\gL{c000000}{1.00}{72.50}{72.00}{72.50}{720.00}
\gL{c000000}{1.00}{501.50}{72.00}{501.50}{720.00}
\gL{c000000}{1.00}{72.00}{72.50}{502.00}{72.50}
\gL{c000000}{1.00}{72.00}{90.50}{502.00}{90.50}
\gL{c000000}{1.00}{72.00}{108.50}{502.00}{108.50}
\gL{c000000}{1.00}{72.00}{126.50}{502.00}{126.50}
\gL{c000000}{1.00}{72.00}{144.50}{502.00}{144.50}
\gL{c000000}{1.00}{72.00}{162.50}{502.00}{162.50}
\gL{c000000}{1.00}{72.00}{180.50}{502.00}{180.50}
\gL{c000000}{1.00}{72.00}{198.50}{502.00}{198.50}
\gL{c000000}{1.00}{72.00}{216.50}{502.00}{216.50}
\gL{c000000}{1.00}{72.00}{234.50}{502.00}{234.50}
\gL{c000000}{1.00}{72.00}{252.50}{502.00}{252.50}
\gL{c000000}{1.00}{72.00}{270.50}{502.00}{270.50}
\gL{c000000}{1.00}{72.00}{288.50}{502.00}{288.50}
\gL{c000000}{1.00}{72.00}{306.50}{502.00}{306.50}
\gL{c000000}{1.00}{72.00}{324.50}{502.00}{324.50}
\gL{c000000}{1.00}{72.00}{342.50}{502.00}{342.50}
\gL{c000000}{1.00}{72.00}{360.50}{502.00}{360.50}
\gL{c000000}{1.00}{72.00}{378.50}{502.00}{378.50}
\gL{c000000}{1.00}{72.00}{396.50}{502.00}{396.50}
\gL{c000000}{1.00}{72.00}{414.50}{502.00}{414.50}
\gL{c000000}{1.00}{72.00}{432.50}{502.00}{432.50}
\gL{c000000}{1.00}{72.00}{450.50}{502.00}{450.50}
\gL{c000000}{1.00}{72.00}{468.50}{502.00}{468.50}
\gL{c000000}{1.00}{72.00}{486.50}{502.00}{486.50}
\gL{c000000}{1.00}{72.00}{504.50}{502.00}{504.50}
\gL{c000000}{1.00}{72.00}{522.50}{502.00}{522.50}
\gL{c000000}{1.00}{72.00}{540.50}{502.00}{540.50}
\gL{c000000}{1.00}{72.00}{558.50}{502.00}{558.50}
\gL{c000000}{1.00}{72.00}{576.50}{502.00}{576.50}
\gL{c000000}{1.00}{72.00}{594.50}{502.00}{594.50}
\gL{c000000}{1.00}{72.00}{612.50}{502.00}{612.50}
\gL{c000000}{1.00}{72.00}{630.50}{502.00}{630.50}
\gL{c000000}{1.00}{72.00}{648.50}{502.00}{648.50}
\gL{c000000}{1.00}{72.00}{666.50}{502.00}{666.50}
\gL{c000000}{1.00}{72.00}{684.50}{502.00}{684.50}
\gL{c000000}{1.00}{72.00}{702.50}{502.00}{702.50}
\gL{c000000}{1.00}{72.00}{720.50}{502.00}{720.50}
\gT{532.01}{43.47}{\gS{R}{8}{c000000}{33 }}
\gT{78.00}{84.21}{\gS{R}{10}{c000000}{Elsevier (Science Publishing Co.), Amsterdam, The Netherlands; New York, NY, USA }}
\gT{78.00}{102.21}{\gS{R}{10}{c000000}{Emerson (Books, Inc.), New York, NY, USA }}
\gT{78.00}{120.21}{\gS{R}{10}{c000000}{Encyclopedia Britannica (Inc.), Chicago, IL, USA }}
\gT{78.00}{138.21}{\gS{R}{10}{c000000}{Engineering Pub. (Inc.), Elizabeth, NJ, USA }}
\gT{78.00}{156.21}{\gS{R}{10}{c000000}{(W.) Englemann, Leipzig, Germany }}
\gT{78.00}{192.21}{\gS{RB}{10}{c000000}{F}}
\gT{78.00}{210.21}{\gS{R}{10}{c000000}{Fairchild (Publications, Inc.), New York, NY, USA }}
\gT{78.00}{228.21}{\gS{R}{10}{c000000}{Farrar, Straus (\& Giroux, Inc.), New York, NY, USA }}
\gT{78.00}{246.21}{\gS{R}{10}{c000000}{(F. W.) Faxon (Inc.), Boston, MA, USA }}
\gT{78.00}{264.20}{\gS{R}{10}{c000000}{(Carl) Fischer (Inc.), New York, NY, USA }}
\gT{78.00}{282.20}{\gS{R}{10}{c000000}{Fleet (Publishing Corp.), New York, NY, USA }}
\gT{78.00}{300.20}{\gS{R}{10}{c000000}{Follett (Publishing Co.), Chicago, IL, USA }}
\gT{78.00}{318.20}{\gS{R}{10}{c000000}{Foundation (Press, Inc.), New York, NY, USA }}
\gT{78.00}{336.20}{\gS{R}{10}{c000000}{Franey (\& Co., Inc.), London, U.K. }}
\gT{78.00}{354.20}{\gS{R}{10}{c000000}{(W.H.) Freeman (and Co.), San Francisco, CA, USA }}
\gT{78.00}{372.20}{\gS{R}{10}{c000000}{Free Press, New York, NY, USA }}
\gT{78.00}{390.20}{\gS{R}{10}{c000000}{(Samuel) French (Inc.), New York, NY, USA }}
\gT{78.00}{408.20}{\gS{R}{10}{c000000}{Friedlander \& Sohn, Berlin, Germany }}
\gT{78.00}{426.20}{\gS{R}{10}{c000000}{Friendship Press, New York, NY, USA }}
\gT{78.00}{444.20}{\gS{R}{10}{c000000}{Funk \& Wagnalls (Co.), New York, NY, USA }}
\gT{78.00}{462.20}{\gS{R}{10}{c000000}{Futura (Publishing Co., Inc.), Mt. Kisco, NY, USA }}
\gT{78.00}{498.20}{\gS{RB}{10}{c000000}{G}}
\gT{78.00}{516.21}{\gS{R}{10}{c000000}{(W. J.) Gage (Ltd.), London, U.K. }}
\gT{78.00}{534.21}{\gS{R}{10}{c000000}{Gale \& Polden (Ltd.), London, U.K. }}
\gT{78.00}{552.21}{\gS{R}{10}{c000000}{Garden City Books, New York, NY, USA }}
\gT{78.00}{570.21}{\gS{R}{10}{c000000}{Garland (Publishing, Inc.), New York, NY, USA }}
\gT{78.00}{588.21}{\gS{R}{10}{c000000}{Garrett (\& Massie, Inc.), Richmond, VA, USA }}
\gT{78.00}{606.21}{\gS{R}{10}{c000000}{Gautheir-Villars, Paris, France }}
\gT{78.00}{624.21}{\gS{R}{10}{c000000}{(Bernard) Geis (Associates), New York, NY, USA }}
\gT{78.00}{642.21}{\gS{R}{10}{c000000}{George Mason Univ. Press, Fairfax, VA, USA }}
\gT{78.00}{660.21}{\gS{R}{10}{c000000}{Globe Book (Co., Inc.), New York, NY, USA }}
\gT{78.00}{678.21}{\gS{R}{10}{c000000}{Golden Cockerel (Press, Inc.), London, U.K. }}
\gT{78.00}{696.21}{\gS{R}{10}{c000000}{Golden (Press, Inc.), New York, NY, USA }}
\gT{78.00}{714.21}{\gS{R}{10}{c000000}{(Victor) Gollancz (Ltd.), London, U.K. }}
\end{guidepage}
\begin{guidepage}% page 34
\gF{cEFEFEF}{72.00}{342.00}{501.75}{360.00}
\gL{c000000}{1.00}{78.00}{355.50}{86.00}{355.50}
\gL{c000000}{1.00}{72.50}{72.00}{72.50}{720.00}
\gL{c000000}{1.00}{501.50}{72.00}{501.50}{720.00}
\gL{c000000}{1.00}{72.00}{72.50}{502.00}{72.50}
\gL{c000000}{1.00}{72.00}{90.50}{502.00}{90.50}
\gL{c000000}{1.00}{72.00}{108.50}{502.00}{108.50}
\gL{c000000}{1.00}{72.00}{126.50}{502.00}{126.50}
\gL{c000000}{1.00}{72.00}{144.50}{502.00}{144.50}
\gL{c000000}{1.00}{72.00}{162.50}{502.00}{162.50}
\gL{c000000}{1.00}{72.00}{180.50}{502.00}{180.50}
\gL{c000000}{1.00}{72.00}{198.50}{502.00}{198.50}
\gL{c000000}{1.00}{72.00}{216.50}{502.00}{216.50}
\gL{c000000}{1.00}{72.00}{234.50}{502.00}{234.50}
\gL{c000000}{1.00}{72.00}{252.50}{502.00}{252.50}
\gL{c000000}{1.00}{72.00}{270.50}{502.00}{270.50}
\gL{c000000}{1.00}{72.00}{288.50}{502.00}{288.50}
\gL{c000000}{1.00}{72.00}{306.50}{502.00}{306.50}
\gL{c000000}{1.00}{72.00}{324.50}{502.00}{324.50}
\gL{c000000}{1.00}{72.00}{342.50}{502.00}{342.50}
\gL{c000000}{1.00}{72.00}{360.50}{502.00}{360.50}
\gL{c000000}{1.00}{72.00}{378.50}{502.00}{378.50}
\gL{c000000}{1.00}{72.00}{396.50}{502.00}{396.50}
\gL{c000000}{1.00}{72.00}{414.50}{502.00}{414.50}
\gL{c000000}{1.00}{72.00}{432.50}{502.00}{432.50}
\gL{c000000}{1.00}{72.00}{450.50}{502.00}{450.50}
\gL{c000000}{1.00}{72.00}{468.50}{502.00}{468.50}
\gL{c000000}{1.00}{72.00}{486.50}{502.00}{486.50}
\gL{c000000}{1.00}{72.00}{504.50}{502.00}{504.50}
\gL{c000000}{1.00}{72.00}{522.50}{502.00}{522.50}
\gL{c000000}{1.00}{72.00}{540.50}{502.00}{540.50}
\gL{c000000}{1.00}{72.00}{558.50}{502.00}{558.50}
\gL{c000000}{1.00}{72.00}{576.50}{502.00}{576.50}
\gL{c000000}{1.00}{72.00}{594.50}{502.00}{594.50}
\gL{c000000}{1.00}{72.00}{612.50}{502.00}{612.50}
\gL{c000000}{1.00}{72.00}{630.50}{502.00}{630.50}
\gL{c000000}{1.00}{72.00}{648.50}{502.00}{648.50}
\gL{c000000}{1.00}{72.00}{666.50}{502.00}{666.50}
\gL{c000000}{1.00}{72.00}{684.50}{502.00}{684.50}
\gL{c000000}{1.00}{72.00}{702.50}{502.00}{702.50}
\gL{c000000}{1.00}{72.00}{720.50}{502.00}{720.50}
\gT{532.01}{43.47}{\gS{R}{8}{c000000}{34 }}
\gT{78.00}{84.21}{\gS{R}{10}{c000000}{Goodheart-Willcox (Co.), Homewood, IL, USA }}
\gT{78.00}{102.21}{\gS{R}{10}{c000000}{Gordon \& Breach (Publishers, Inc.), New York, NY, USA }}
\gT{78.00}{120.21}{\gS{R}{10}{c000000}{Gower (Medical Pub.), New York, NY, USA }}
\gT{78.00}{138.21}{\gS{R}{10}{c000000}{GPO (Government Printing Office), Washington, DC, USA }}
\gT{78.00}{156.21}{\gS{R}{10}{c000000}{Greenberg (Publishers, Inc.), New York, NY, USA }}
\gT{78.00}{174.21}{\gS{R}{10}{c000000}{Greystone (Press), New York, NY, USA }}
\gT{78.00}{192.21}{\gS{R}{10}{c000000}{(Charles) Griffin (\& Co., Inc.), London, U.K. }}
\gT{78.00}{210.21}{\gS{R}{10}{c000000}{Grolier Club, New York, NY, USA }}
\gT{78.00}{228.21}{\gS{R}{10}{c000000}{Grolier (Inc.), New York, NY, USA }}
\gT{78.00}{246.21}{\gS{R}{10}{c000000}{Grosset \& Dunlap (Inc.), New York, NY, USA }}
\gT{78.00}{264.20}{\gS{R}{10}{c000000}{Grove (Press), New York, NY, USA }}
\gT{78.00}{282.20}{\gS{R}{10}{c000000}{Grune \& Stratton (Inc.), New York, NY, USA }}
\gT{78.00}{300.20}{\gS{R}{10}{c000000}{Gulf (Publishing Co.), Houston, TX, USA }}
\gT{78.00}{318.20}{\gS{R}{10}{c000000}{Gurney (\& Jackson), London, U.K. }}
\gT{78.00}{354.20}{\gS{RB}{10}{c000000}{H}}
\gT{78.00}{372.20}{\gS{R}{10}{c000000}{Hafner (Publishing Co.), New York, NY, USA }}
\gT{78.00}{390.20}{\gS{R}{10}{c000000}{E. M. Hale (\& Co.), Eau Claire, WI, USA }}
\gT{78.00}{408.20}{\gS{R}{10}{c000000}{Halsted (Press), New York, NY, USA }}
\gT{78.00}{426.20}{\gS{R}{10}{c000000}{Hammond (Inc.), Maplewood, NJ, USA }}
\gT{78.00}{444.20}{\gS{R}{10}{c000000}{Harcourt Brace Jovanovich, Orlando, FL, USA }}
\gT{78.00}{462.20}{\gS{R}{10}{c000000}{(A.) Harleben, Vienna, Austria }}
\gT{78.00}{480.20}{\gS{R}{10}{c000000}{Harper \& Row, New York, NY, USA }}
\gT{78.00}{498.20}{\gS{R}{10}{c000000}{Harper Collins, New York, NY, USA }}
\gT{78.00}{516.21}{\gS{R}{10}{c000000}{W. E. Harrison (\& Sons), Ipswich, U.K. }}
\gT{78.00}{534.21}{\gS{R}{10}{c000000}{Harrison (\& Sons, Ltd.), London, U.K. }}
\gT{78.00}{552.21}{\gS{R}{10}{c000000}{Harvard Univ. Press, Cambridge, MA, USA }}
\gT{78.00}{570.21}{\gS{R}{10}{c000000}{Hastings (House, Publishers), New York, NY, USA }}
\gT{78.00}{588.21}{\gS{R}{10}{c000000}{Hawthorn (Books, Inc.), New York, NY, USA }}
\gT{78.00}{606.21}{\gS{R}{10}{c000000}{Hayden (Book Co.), New York, NY, USA }}
\gT{78.00}{624.21}{\gS{R}{10}{c000000}{(D. C.) Heath (and Co.), Boston, MA, USA }}
\gT{78.00}{642.21}{\gS{R}{10}{c000000}{Heath, Cranton (Ltd.), London, U.K. }}
\gT{78.00}{660.21}{\gS{R}{10}{c000000}{(W.) Heffer (\& Sons, Ltd.), Cambridge, U.K. }}
\gT{78.00}{678.21}{\gS{R}{10}{c000000}{(William) Heinemann (Ltd.), London, U.K. }}
\gT{78.00}{696.21}{\gS{R}{10}{c000000}{Helicon (Press, Inc.), Baltimore, MD, USA }}
\gT{78.00}{714.21}{\gS{R}{10}{c000000}{Hemisphere (Publishing Corp.), Bristol, PA, USA }}
\end{guidepage}
\begin{guidepage}% page 35
\gF{cEFEFEF}{72.00}{414.00}{501.75}{432.00}
\gL{c000000}{1.00}{78.00}{427.50}{82.00}{427.50}
\gL{c000000}{1.00}{72.50}{72.00}{72.50}{720.00}
\gL{c000000}{1.00}{501.50}{72.00}{501.50}{720.00}
\gL{c000000}{1.00}{72.00}{72.50}{502.00}{72.50}
\gL{c000000}{1.00}{72.00}{90.50}{502.00}{90.50}
\gL{c000000}{1.00}{72.00}{108.50}{502.00}{108.50}
\gL{c000000}{1.00}{72.00}{126.50}{502.00}{126.50}
\gL{c000000}{1.00}{72.00}{144.50}{502.00}{144.50}
\gL{c000000}{1.00}{72.00}{162.50}{502.00}{162.50}
\gL{c000000}{1.00}{72.00}{180.50}{502.00}{180.50}
\gL{c000000}{1.00}{72.00}{198.50}{502.00}{198.50}
\gL{c000000}{1.00}{72.00}{216.50}{502.00}{216.50}
\gL{c000000}{1.00}{72.00}{234.50}{502.00}{234.50}
\gL{c000000}{1.00}{72.00}{252.50}{502.00}{252.50}
\gL{c000000}{1.00}{72.00}{270.50}{502.00}{270.50}
\gL{c000000}{1.00}{72.00}{288.50}{502.00}{288.50}
\gL{c000000}{1.00}{72.00}{306.50}{502.00}{306.50}
\gL{c000000}{1.00}{72.00}{324.50}{502.00}{324.50}
\gL{c000000}{1.00}{72.00}{342.50}{502.00}{342.50}
\gL{c000000}{1.00}{72.00}{360.50}{502.00}{360.50}
\gL{c000000}{1.00}{72.00}{378.50}{502.00}{378.50}
\gL{c000000}{1.00}{72.00}{396.50}{502.00}{396.50}
\gL{c000000}{1.00}{72.00}{414.50}{502.00}{414.50}
\gL{c000000}{1.00}{72.00}{432.50}{502.00}{432.50}
\gL{c000000}{1.00}{72.00}{450.50}{502.00}{450.50}
\gL{c000000}{1.00}{72.00}{468.50}{502.00}{468.50}
\gL{c000000}{1.00}{72.00}{486.50}{502.00}{486.50}
\gL{c000000}{1.00}{72.00}{504.50}{502.00}{504.50}
\gL{c000000}{1.00}{72.00}{522.50}{502.00}{522.50}
\gL{c000000}{1.00}{72.00}{540.50}{502.00}{540.50}
\gL{c000000}{1.00}{72.00}{558.50}{502.00}{558.50}
\gL{c000000}{1.00}{72.00}{576.50}{502.00}{576.50}
\gL{c000000}{1.00}{72.00}{594.50}{502.00}{594.50}
\gL{c000000}{1.00}{72.00}{612.50}{502.00}{612.50}
\gL{c000000}{1.00}{72.00}{630.50}{502.00}{630.50}
\gL{c000000}{1.00}{72.00}{648.50}{502.00}{648.50}
\gL{c000000}{1.00}{72.00}{666.50}{502.00}{666.50}
\gL{c000000}{1.00}{72.00}{684.50}{502.00}{684.50}
\gL{c000000}{1.00}{72.00}{702.50}{502.00}{702.50}
\gL{c000000}{1.00}{72.00}{720.50}{502.00}{720.50}
\gT{532.01}{43.47}{\gS{R}{8}{c000000}{35 }}
\gT{78.00}{84.21}{\gS{R}{10}{c000000}{(B.) Herder (Book Co.), St. Louis, MO, USA }}
\gT{78.00}{102.21}{\gS{R}{10}{c000000}{Heritage (Press), New York, NY, USA }}
\gT{78.00}{120.21}{\gS{R}{10}{c000000}{Hermann \& Cie, Paris, France }}
\gT{78.00}{138.21}{\gS{R}{10}{c000000}{Hermitage (House, Inc.), New York, NY, USA }}
\gT{78.00}{156.21}{\gS{R}{10}{c000000}{Hill \& Wang (Inc.), New York, NY, USA }}
\gT{78.00}{174.21}{\gS{R}{10}{c000000}{Hindawi, Cairo, Egypt; New York, NY, USA }}
\gT{78.00}{192.21}{\gS{R}{10}{c000000}{H. M. Stationary Office, London, U.K. }}
\gT{78.00}{210.21}{\gS{R}{10}{c000000}{Hodder (\& Stroughton, Ltd.), London, U.K. }}
\gT{78.00}{228.21}{\gS{R}{10}{c000000}{(Paul B.) Hoeber, New York, NY, USA }}
\gT{78.00}{246.21}{\gS{R}{10}{c000000}{Hogarth (Press, Ltd.), London, U.K. }}
\gT{78.00}{264.20}{\gS{R}{10}{c000000}{Holiday (House, Inc.), New York, NY, USA }}
\gT{78.00}{282.20}{\gS{R}{10}{c000000}{Holt, Rinehart and Winston (Inc.), New York, NY, USA }}
\gT{78.00}{300.20}{\gS{R}{10}{c000000}{Horizon (Press, Inc.), New York, NY, USA }}
\gT{78.00}{318.20}{\gS{R}{10}{c000000}{Houghton Mifflin (Co.), Boston, MA, USA }}
\gT{78.00}{336.20}{\gS{R}{10}{c000000}{Humanities Press (Inc.), New York, NY, USA }}
\gT{78.00}{354.20}{\gS{R}{10}{c000000}{(Bruce) Humphries (Inc.), New York, NY, USA }}
\gT{78.00}{372.20}{\gS{R}{10}{c000000}{Hutchinson (\& Co., Publisher, Ltd.), London, U.K. }}
\gT{78.00}{390.20}{\gS{R}{10}{c000000}{Hyperion (Press, Inc.), New York, NY, USA }}
\gT{78.00}{426.20}{\gS{RB}{10}{c000000}{I}}
\gT{78.00}{444.20}{\gS{R}{10}{c000000}{IEEE Comput. Soc. Press, Los Alamitos, CA, USA }}
\gT{78.00}{462.20}{\gS{R}{10}{c000000}{IEEE Press, Piscataway, NJ, USA }}
\gT{78.00}{480.20}{\gS{R}{10}{c000000}{Indiana Univ. Press, Bloomington, IN, USA }}
\gT{78.00}{498.20}{\gS{R}{10}{c000000}{Industrial Press, New York, NY, USA }}
\gT{78.00}{516.21}{\gS{R}{10}{c000000}{Industrial (Publications, Inc.), Chicago, IL, USA }}
\gT{78.00}{534.21}{\gS{R}{10}{c000000}{Inst. for Advanced Studies, Princeton, NJ, USA }}
\gT{78.00}{552.21}{\gS{R}{10}{c000000}{Instrum. Soc. Amer., Research Triangle Park, NC, USA }}
\gT{78.00}{570.21}{\gS{R}{10}{c000000}{International Labor Office, Washington, DC, USA }}
\gT{78.00}{588.21}{\gS{R}{10}{c000000}{International Publishers (Co., Inc.), New York, NY, USA }}
\gT{78.00}{606.21}{\gS{R}{10}{c000000}{International Textbook (Co.), Scranton, PA, USA }}
\gT{78.00}{624.21}{\gS{R}{10}{c000000}{International Universities (Press, Inc.), New York, NY, USA }}
\gT{78.00}{642.21}{\gS{R}{10}{c000000}{Interscience (Publishers), New York, NY, USA }}
\gT{78.00}{660.21}{\gS{R}{10}{c000000}{IOS Press, Amsterdam, The Netherlands }}
\gT{78.00}{678.21}{\gS{R}{10}{c000000}{Iowa State Univ. Press, Ames, IA, USA }}
\gT{78.00}{696.21}{\gS{R}{10}{c000000}{Iroquois (Publishing Co., Inc.), Columbus, OH, USA }}
\gT{78.00}{714.21}{\gS{R}{10}{c000000}{(Richard D.) Irwin (Inc.), Homewood, IL, USA }}
\end{guidepage}
\begin{guidepage}% page 36
\gF{cEFEFEF}{72.00}{90.00}{501.75}{108.00}
\gF{cEFEFEF}{72.00}{252.00}{501.75}{270.00}
\gF{cEFEFEF}{72.00}{504.00}{501.75}{522.00}
\gF{cEFEFEF}{72.00}{576.00}{501.75}{594.00}
\gL{c000000}{1.00}{78.00}{103.50}{83.00}{103.50}
\gL{c000000}{1.00}{78.00}{265.50}{86.00}{265.50}
\gL{c000000}{1.00}{78.00}{517.50}{85.00}{517.50}
\gL{c000000}{1.00}{78.00}{589.50}{87.00}{589.50}
\gL{c000000}{1.00}{72.50}{72.00}{72.50}{720.00}
\gL{c000000}{1.00}{501.50}{72.00}{501.50}{720.00}
\gL{c000000}{1.00}{72.00}{72.50}{502.00}{72.50}
\gL{c000000}{1.00}{72.00}{90.50}{502.00}{90.50}
\gL{c000000}{1.00}{72.00}{108.50}{502.00}{108.50}
\gL{c000000}{1.00}{72.00}{126.50}{502.00}{126.50}
\gL{c000000}{1.00}{72.00}{144.50}{502.00}{144.50}
\gL{c000000}{1.00}{72.00}{162.50}{502.00}{162.50}
\gL{c000000}{1.00}{72.00}{180.50}{502.00}{180.50}
\gL{c000000}{1.00}{72.00}{198.50}{502.00}{198.50}
\gL{c000000}{1.00}{72.00}{216.50}{502.00}{216.50}
\gL{c000000}{1.00}{72.00}{234.50}{502.00}{234.50}
\gL{c000000}{1.00}{72.00}{252.50}{502.00}{252.50}
\gL{c000000}{1.00}{72.00}{270.50}{502.00}{270.50}
\gL{c000000}{1.00}{72.00}{288.50}{502.00}{288.50}
\gL{c000000}{1.00}{72.00}{306.50}{502.00}{306.50}
\gL{c000000}{1.00}{72.00}{324.50}{502.00}{324.50}
\gL{c000000}{1.00}{72.00}{342.50}{502.00}{342.50}
\gL{c000000}{1.00}{72.00}{360.50}{502.00}{360.50}
\gL{c000000}{1.00}{72.00}{378.50}{502.00}{378.50}
\gL{c000000}{1.00}{72.00}{396.50}{502.00}{396.50}
\gL{c000000}{1.00}{72.00}{414.50}{502.00}{414.50}
\gL{c000000}{1.00}{72.00}{432.50}{502.00}{432.50}
\gL{c000000}{1.00}{72.00}{450.50}{502.00}{450.50}
\gL{c000000}{1.00}{72.00}{468.50}{502.00}{468.50}
\gL{c000000}{1.00}{72.00}{486.50}{502.00}{486.50}
\gL{c000000}{1.00}{72.00}{504.50}{502.00}{504.50}
\gL{c000000}{1.00}{72.00}{522.50}{502.00}{522.50}
\gL{c000000}{1.00}{72.00}{540.50}{502.00}{540.50}
\gL{c000000}{1.00}{72.00}{558.50}{502.00}{558.50}
\gL{c000000}{1.00}{72.00}{576.50}{502.00}{576.50}
\gL{c000000}{1.00}{72.00}{594.50}{502.00}{594.50}
\gL{c000000}{1.00}{72.00}{612.50}{502.00}{612.50}
\gL{c000000}{1.00}{72.00}{630.50}{502.00}{630.50}
\gL{c000000}{1.00}{72.00}{648.50}{502.00}{648.50}
\gL{c000000}{1.00}{72.00}{666.50}{502.00}{666.50}
\gL{c000000}{1.00}{72.00}{684.50}{502.00}{684.50}
\gL{c000000}{1.00}{72.00}{702.50}{502.00}{702.50}
\gL{c000000}{1.00}{72.00}{720.50}{502.00}{720.50}
\gT{532.01}{43.47}{\gS{R}{8}{c000000}{36 }}
\gT{78.00}{102.21}{\gS{RB}{10}{c000000}{J}}
\gT{78.00}{120.21}{\gS{R}{10}{c000000}{JAI Press, Greenwich, CT, USA }}
\gT{78.00}{138.21}{\gS{R}{10}{c000000}{Jarrolds (Publishers, Ltd.), London, U.K. }}
\gT{78.00}{156.21}{\gS{R}{10}{c000000}{The Johns Hopkins Univ. Press, Baltimore, MD, USA }}
\gT{78.00}{174.21}{\gS{R}{10}{c000000}{Jones \& Bartlett, Boston, MA, USA }}
\gT{78.00}{192.21}{\gS{R}{10}{c000000}{Jossey-Bass (Inc.), San Francisco, CA, USA }}
\gT{78.00}{210.21}{\gS{R}{10}{c000000}{Julian (Press, Inc.), New York, NY, USA }}
\gT{78.00}{228.21}{\gS{R}{10}{c000000}{Juta (\& Co., Ltd.), Cape Town, South Africa }}
\gT{78.00}{264.20}{\gS{RB}{10}{c000000}{K}}
\gT{78.00}{282.20}{\gS{R}{10}{c000000}{(Willibald) Keller, Leipzig, Germany }}
\gT{78.00}{300.20}{\gS{R}{10}{c000000}{(P. J.) Kenedy (\& Sons), New York, NY, USA }}
\gT{78.00}{318.20}{\gS{R}{10}{c000000}{King’s Crown (Press), New York, NY, USA }}
\gT{78.00}{336.20}{\gS{R}{10}{c000000}{Kluwer (Academic), Norwell, MA, USA (also, Boston, MA, USA) }}
\gT{78.00}{354.20}{\gS{R}{10}{c000000}{Kluwer-Nijhoff (Publishing), Norwell, MA, USA }}
\gT{78.00}{372.20}{\gS{R}{10}{c000000}{(W.) Knapp Verlag, Dusseldorf, Germany }}
\gT{78.00}{390.20}{\gS{R}{10}{c000000}{Knickerbocker Press, New York, NY, USA }}
\gT{78.00}{408.20}{\gS{R}{10}{c000000}{Knight (Publishers, Inc.), New York, NY, USA }}
\gT{78.00}{426.20}{\gS{R}{10}{c000000}{(Alfred A.) Knopf, New York, NY, USA }}
\gT{78.00}{444.20}{\gS{R}{10}{c000000}{Krieger (Publishing Co.), Melbourne, FL, USA }}
\gT{78.00}{462.20}{\gS{R}{10}{c000000}{(Alfred) Kroener, Stuttgart, Germany }}
\gT{78.00}{480.20}{\gS{R}{10}{c000000}{(Edward) Kummer, Leipzig, Germany }}
\gT{78.00}{516.21}{\gS{RB}{10}{c000000}{L}}
\gT{78.00}{534.21}{\gS{R}{10}{c000000}{Laidlaw (Brothers, Inc.), River Forest, IL, USA }}
\gT{78.00}{552.21}{\gS{R}{10}{c000000}{Longman (Publishing Group), White Plains, NY, USA }}
\gT{78.00}{588.21}{\gS{RB}{10}{c000000}{M}}
\gT{78.00}{606.21}{\gS{R}{10}{c000000}{Macmillan (Inc.), New York, NY, USA }}
\gT{78.00}{624.21}{\gS{R}{10}{c000000}{Marcel Dekker (Inc.), New York, NY, USA }}
\gT{78.00}{642.21}{\gS{R}{10}{c000000}{Materials Research Soc., Pittsburgh, PA, USA }}
\gT{78.00}{660.21}{\gS{R}{10}{c000000}{Matrix, Beaverton, OR, USA }}
\gT{78.00}{678.21}{\gS{R}{10}{c000000}{McGraw-Hill (Inc.), New York, NY, USA }}
\gT{78.00}{696.21}{\gS{R}{10}{c000000}{MIT Press, Cambridge, MA, USA }}
\gT{78.00}{714.21}{\gS{R}{10}{c000000}{Morgan \& Claypool (Publishers), San Rafael, CA, USA }}
\end{guidepage}
\begin{guidepage}% page 37
\gF{cEFEFEF}{72.00}{198.00}{501.75}{216.00}
\gF{cEFEFEF}{72.00}{630.00}{501.75}{648.00}
\gL{c000000}{1.00}{78.00}{211.50}{85.00}{211.50}
\gL{c000000}{1.00}{78.00}{643.50}{86.00}{643.50}
\gL{c000000}{1.00}{72.50}{72.00}{72.50}{720.00}
\gL{c000000}{1.00}{501.50}{72.00}{501.50}{720.00}
\gL{c000000}{1.00}{72.00}{72.50}{502.00}{72.50}
\gL{c000000}{1.00}{72.00}{90.50}{502.00}{90.50}
\gL{c000000}{1.00}{72.00}{108.50}{502.00}{108.50}
\gL{c000000}{1.00}{72.00}{126.50}{502.00}{126.50}
\gL{c000000}{1.00}{72.00}{144.50}{502.00}{144.50}
\gL{c000000}{1.00}{72.00}{162.50}{502.00}{162.50}
\gL{c000000}{1.00}{72.00}{180.50}{502.00}{180.50}
\gL{c000000}{1.00}{72.00}{198.50}{502.00}{198.50}
\gL{c000000}{1.00}{72.00}{216.50}{502.00}{216.50}
\gL{c000000}{1.00}{72.00}{234.50}{502.00}{234.50}
\gL{c000000}{1.00}{72.00}{252.50}{502.00}{252.50}
\gL{c000000}{1.00}{72.00}{270.50}{502.00}{270.50}
\gL{c000000}{1.00}{72.00}{288.50}{502.00}{288.50}
\gL{c000000}{1.00}{72.00}{306.50}{502.00}{306.50}
\gL{c000000}{1.00}{72.00}{324.50}{502.00}{324.50}
\gL{c000000}{1.00}{72.00}{342.50}{502.00}{342.50}
\gL{c000000}{1.00}{72.00}{360.50}{502.00}{360.50}
\gL{c000000}{1.00}{72.00}{378.50}{502.00}{378.50}
\gL{c000000}{1.00}{72.00}{396.50}{502.00}{396.50}
\gL{c000000}{1.00}{72.00}{414.50}{502.00}{414.50}
\gL{c000000}{1.00}{72.00}{432.50}{502.00}{432.50}
\gL{c000000}{1.00}{72.00}{450.50}{502.00}{450.50}
\gL{c000000}{1.00}{72.00}{468.50}{502.00}{468.50}
\gL{c000000}{1.00}{72.00}{486.50}{502.00}{486.50}
\gL{c000000}{1.00}{72.00}{504.50}{502.00}{504.50}
\gL{c000000}{1.00}{72.00}{522.50}{502.00}{522.50}
\gL{c000000}{1.00}{72.00}{540.50}{502.00}{540.50}
\gL{c000000}{1.00}{72.00}{558.50}{502.00}{558.50}
\gL{c000000}{1.00}{72.00}{576.50}{502.00}{576.50}
\gL{c000000}{1.00}{72.00}{594.50}{502.00}{594.50}
\gL{c000000}{1.00}{72.00}{612.50}{502.00}{612.50}
\gL{c000000}{1.00}{72.00}{630.50}{502.00}{630.50}
\gL{c000000}{1.00}{72.00}{648.50}{502.00}{648.50}
\gL{c000000}{1.00}{72.00}{666.50}{502.00}{666.50}
\gL{c000000}{1.00}{72.00}{684.50}{502.00}{684.50}
\gL{c000000}{1.00}{72.00}{702.50}{502.00}{702.50}
\gL{c000000}{1.00}{72.00}{720.50}{502.00}{720.50}
\gT{532.01}{43.47}{\gS{R}{8}{c000000}{37 }}
\gT{78.00}{84.21}{\gS{R}{10}{c000000}{Morgan Kaufmann (Publishers, Inc.), San Mateo, CA, USA }}
\gT{78.00}{102.21}{\gS{R}{10}{c000000}{(The C. V.) Mosby (Co.), St. Louis, MO, USA }}
\gT{78.00}{120.21}{\gS{R}{10}{c000000}{Muhlenbert Press, Philadelphia, PA, USA }}
\gT{78.00}{138.21}{\gS{R}{10}{c000000}{J. Murray, London, U.K. }}
\gT{78.00}{156.21}{\gS{R}{10}{c000000}{Murray Hill Books, New York, NY, USA }}
\gT{78.00}{174.21}{\gS{R}{10}{c000000}{Museum of Modern Art, New York, NY, USA }}
\gT{78.00}{210.21}{\gS{RB}{10}{c000000}{N}}
\gT{78.00}{228.21}{\gS{R}{10}{c000000}{National Academy of Sciences–Nat. Res. Council, Washington, DC, USA }}
\gT{78.00}{246.21}{\gS{R}{10}{c000000}{National Academy Press, Washington, DC, USA }}
\gT{78.00}{264.20}{\gS{R}{10}{c000000}{National Education Association, Washington, DC, USA }}
\gT{78.00}{282.20}{\gS{R}{10}{c000000}{National Geographic Society, Washington, DC, USA }}
\gT{78.00}{300.20}{\gS{R}{10}{c000000}{National Industrial Conference Board, Inc., New York, NY, USA }}
\gT{78.00}{318.20}{\gS{R}{10}{c000000}{Naval Inst. Press, Annapolis, MD, USA }}
\gT{78.00}{336.20}{\gS{R}{10}{c000000}{Naylor (Co.), San Antonio, TX, USA }}
\gT{78.00}{354.20}{\gS{R}{10}{c000000}{C. C. Nelson (Publishing Co.), Appleton, WI, USA }}
\gT{78.00}{372.20}{\gS{R}{10}{c000000}{(Thomas) Nelson (\& Sons), Camden, NJ, USA }}
\gT{78.00}{390.20}{\gS{R}{10}{c000000}{New American Library of World Literature (Inc.), New York, NY, USA }}
\gT{78.00}{408.20}{\gS{R}{10}{c000000}{New Directions, New York, NY, USA }}
\gT{78.00}{426.20}{\gS{R}{10}{c000000}{New York Graphic (Society), Greenwich, CT, USA }}
\gT{78.00}{444.20}{\gS{R}{10}{c000000}{New York Univ. Press, New York, NY, USA }}
\gT{78.00}{462.20}{\gS{R}{10}{c000000}{Noble (\& Noble, Publishers, Inc.), New York, NY, USA }}
\gT{78.00}{480.20}{\gS{R}{10}{c000000}{Nonesuch (Library, Ltd.), London, U.K. }}
\gT{78.00}{498.20}{\gS{R}{10}{c000000}{Noonday (Press), New York, NY, USA }}
\gT{78.00}{516.21}{\gS{R}{10}{c000000}{(P.) Noordhoff (N. V.), Gronigen, The Netherlands }}
\gT{78.00}{534.21}{\gS{R}{10}{c000000}{North Holland, Amsterdam, The Netherlands }}
\gT{78.00}{552.21}{\gS{R}{10}{c000000}{(W. W.) Norton (\& Co.), New York, NY, USA }}
\gT{78.00}{570.21}{\gS{R}{10}{c000000}{Nova (Science Publishers), Commack, NY, USA }}
\gT{78.00}{588.21}{\gS{R}{10}{c000000}{Now (Publishers), Boston, MA, USA }}
\gT{78.00}{606.21}{\gS{R}{10}{c000000}{Noyes (Publications), Park Ridge, NJ, USA }}
\gT{78.00}{642.21}{\gS{RB}{10}{c000000}{O}}
\gT{78.00}{660.21}{\gS{R}{10}{c000000}{(Ivan) Obolensky (Inc.), New York, NY, USA }}
\gT{78.00}{678.21}{\gS{R}{10}{c000000}{Oceana (Publications, Inc.), Dobbs Ferry, NY, USA }}
\gT{78.00}{696.21}{\gS{R}{10}{c000000}{Odyssey (Press, Inc.), New York, NY, USA }}
\gT{78.00}{714.21}{\gS{R}{10}{c000000}{Ohio State Univ. Press, Columbus, OH, USA }}
\end{guidepage}
\begin{guidepage}% page 38
\gF{cEFEFEF}{72.00}{270.00}{501.75}{288.00}
\gL{c000000}{1.00}{78.00}{283.50}{84.00}{283.50}
\gL{c000000}{1.00}{72.50}{72.00}{72.50}{719.00}
\gL{c000000}{1.00}{501.50}{72.00}{501.50}{719.00}
\gL{c000000}{1.00}{72.00}{72.50}{502.00}{72.50}
\gL{c000000}{1.00}{72.00}{90.50}{502.00}{90.50}
\gL{c000000}{1.00}{72.00}{108.50}{502.00}{108.50}
\gL{c000000}{1.00}{72.00}{126.50}{502.00}{126.50}
\gL{c000000}{1.00}{72.00}{144.50}{502.00}{144.50}
\gL{c000000}{1.00}{72.00}{162.50}{502.00}{162.50}
\gL{c000000}{1.00}{72.00}{180.50}{502.00}{180.50}
\gL{c000000}{1.00}{72.00}{198.50}{502.00}{198.50}
\gL{c000000}{1.00}{72.00}{216.50}{502.00}{216.50}
\gL{c000000}{1.00}{72.00}{234.50}{502.00}{234.50}
\gL{c000000}{1.00}{72.00}{252.50}{502.00}{252.50}
\gL{c000000}{1.00}{72.00}{270.50}{502.00}{270.50}
\gL{c000000}{1.00}{72.00}{288.50}{502.00}{288.50}
\gL{c000000}{1.00}{72.00}{306.50}{502.00}{306.50}
\gL{c000000}{1.00}{72.00}{324.50}{502.00}{324.50}
\gL{c000000}{1.00}{72.00}{342.50}{502.00}{342.50}
\gL{c000000}{1.00}{72.00}{360.50}{502.00}{360.50}
\gL{c000000}{1.00}{72.00}{378.50}{502.00}{378.50}
\gL{c000000}{1.00}{72.00}{396.50}{502.00}{396.50}
\gL{c000000}{1.00}{72.00}{414.50}{502.00}{414.50}
\gL{c000000}{1.00}{72.00}{432.50}{502.00}{432.50}
\gL{c000000}{1.00}{72.00}{450.50}{502.00}{450.50}
\gL{c000000}{1.00}{72.00}{468.50}{502.00}{468.50}
\gL{c000000}{1.00}{72.00}{486.50}{502.00}{486.50}
\gL{c000000}{1.00}{72.00}{504.50}{502.00}{504.50}
\gL{c000000}{1.00}{72.00}{522.50}{502.00}{522.50}
\gL{c000000}{1.00}{72.00}{540.50}{502.00}{540.50}
\gL{c000000}{1.00}{72.00}{558.50}{502.00}{558.50}
\gL{c000000}{1.00}{72.00}{576.50}{502.00}{576.50}
\gL{c000000}{1.00}{72.00}{594.50}{502.00}{594.50}
\gL{c000000}{1.00}{72.00}{612.50}{502.00}{612.50}
\gL{c000000}{1.00}{72.00}{630.50}{502.00}{630.50}
\gL{c000000}{1.00}{72.00}{648.50}{502.00}{648.50}
\gL{c000000}{1.00}{72.00}{666.50}{502.00}{666.50}
\gL{c000000}{1.00}{72.00}{684.50}{502.00}{684.50}
\gL{c000000}{1.00}{72.00}{702.50}{502.00}{702.50}
\gL{c000000}{1.00}{72.00}{718.50}{502.00}{718.50}
\gT{532.01}{43.47}{\gS{R}{8}{c000000}{38 }}
\gT{78.00}{84.21}{\gS{R}{10}{c000000}{R. Oldenbourg-Verlag, Munich, Germany }}
\gT{78.00}{102.21}{\gS{R}{10}{c000000}{Oliver \& Boyd (Ltd.), Edinburgh, U.K.; London, U.K. }}
\gT{78.00}{120.21}{\gS{R}{10}{c000000}{Open Court (Publishing Co.), La Salle, IL, USA }}
\gT{78.00}{138.21}{\gS{R}{10}{c000000}{Orange Judd (Publishing Co., Inc.), New York, NY, USA }}
\gT{78.00}{156.21}{\gS{R}{10}{c000000}{Orell Fussli-Verlag, Zurich, Switzerland }}
\gT{78.00}{174.21}{\gS{R}{10}{c000000}{Orion (Press, Inc.), New York, NY, USA }}
\gT{78.00}{192.21}{\gS{R}{10}{c000000}{Oxford, New York, NY, USA }}
\gT{78.00}{210.21}{\gS{R}{10}{c000000}{Oxford Book (Co., Inc.), New York, NY, USA }}
\gT{78.00}{228.21}{\gS{R}{10}{c000000}{Oxford Science (Publications), Oxford, U.K. }}
\gT{78.00}{246.21}{\gS{R}{10}{c000000}{Oxford Univ. Press, London, U.K. (also, New York, NY, USA) }}
\gT{78.00}{282.20}{\gS{RB}{10}{c000000}{P}}
\gT{78.00}{300.20}{\gS{R}{10}{c000000}{(L. C.) Page (\& Co.), Boston, MA, USA }}
\gT{78.00}{318.20}{\gS{R}{10}{c000000}{Pageant Books (Inc.), New York, NY, USA }}
\gT{78.00}{336.20}{\gS{R}{10}{c000000}{Pantheon Books (Inc.), New York, NY, USA }}
\gT{78.00}{354.20}{\gS{R}{10}{c000000}{Parents’ Inst. (Inc.), Baltimore, MD, USA }}
\gT{78.00}{372.20}{\gS{R}{10}{c000000}{Parthenon (Publ. Group), Park Ridge, NJ, USA }}
\gT{78.00}{390.20}{\gS{R}{10}{c000000}{Paul Parey Verlag, Berlin, Germany }}
\gT{78.00}{408.20}{\gS{R}{10}{c000000}{Pellegrini \& Cudahy, New York, NY, USA }}
\gT{78.00}{426.20}{\gS{R}{10}{c000000}{Penguin (Books), Baltimore, MD, USA }}
\gT{78.00}{444.20}{\gS{R}{10}{c000000}{Peninsula (Publishing), Los Altos, CA, USA }}
\gT{78.00}{462.20}{\gS{R}{10}{c000000}{(The William) Penn (Publishing Corp.), New York, NY, USA }}
\gT{78.00}{480.20}{\gS{R}{10}{c000000}{Penn. State Univ. Press, University Park, PA, USA }}
\gT{78.00}{498.20}{\gS{R}{10}{c000000}{Pentech, Mountain View, CA, USA }}
\gT{78.00}{516.21}{\gS{R}{10}{c000000}{Penton (Publishing Co.), Cleveland, OH, USA }}
\gT{78.00}{534.21}{\gS{R}{10}{c000000}{(Peter) Peregrinus, Stevenage, U.K. }}
\gT{78.00}{552.21}{\gS{R}{10}{c000000}{Pergamon (Press, Inc.), New York, NY, USA }}
\gT{78.00}{570.21}{\gS{R}{10}{c000000}{Phaidon (Press, Ltd.), London, U.K. }}
\gT{78.00}{588.21}{\gS{R}{10}{c000000}{Pharos Editions, London, U.K. }}
\gT{78.00}{606.21}{\gS{R}{10}{c000000}{(George) Philip (\& Sons, Ltd.), London, U.K. }}
\gT{78.00}{624.21}{\gS{R}{10}{c000000}{Philosophical Library (Inc.), New York, NY, USA }}
\gT{78.00}{642.21}{\gS{R}{10}{c000000}{Pitman (Publishing Corp.), New York, NY, USA }}
\gT{78.00}{660.21}{\gS{R}{10}{c000000}{Sir Isaac Pitman (\& Sons, Ltd.), London, U.K. }}
\gT{78.00}{678.21}{\gS{R}{10}{c000000}{Platt \& Monk (Co., Inc.), New York, NY, USA }}
\gT{78.00}{696.21}{\gS{R}{10}{c000000}{Plays (Inc.), Boston, MA, USA }}
\gT{78.00}{713.46}{\gS{R}{10}{c000000}{Pleiad Press, New York, NY, USA }}
\end{guidepage}
\begin{guidepage}% page 39
\gF{cEFEFEF}{72.00}{270.00}{501.75}{288.00}
\gF{cEFEFEF}{72.00}{360.00}{501.75}{378.00}
\gL{c000000}{1.00}{78.00}{283.50}{86.00}{283.50}
\gL{c000000}{1.00}{78.00}{373.50}{85.00}{373.50}
\gL{c000000}{1.00}{72.50}{72.00}{72.50}{720.00}
\gL{c000000}{1.00}{501.50}{72.00}{501.50}{720.00}
\gL{c000000}{1.00}{72.00}{72.50}{502.00}{72.50}
\gL{c000000}{1.00}{72.00}{90.50}{502.00}{90.50}
\gL{c000000}{1.00}{72.00}{108.50}{502.00}{108.50}
\gL{c000000}{1.00}{72.00}{126.50}{502.00}{126.50}
\gL{c000000}{1.00}{72.00}{144.50}{502.00}{144.50}
\gL{c000000}{1.00}{72.00}{162.50}{502.00}{162.50}
\gL{c000000}{1.00}{72.00}{180.50}{502.00}{180.50}
\gL{c000000}{1.00}{72.00}{198.50}{502.00}{198.50}
\gL{c000000}{1.00}{72.00}{216.50}{502.00}{216.50}
\gL{c000000}{1.00}{72.00}{234.50}{502.00}{234.50}
\gL{c000000}{1.00}{72.00}{252.50}{502.00}{252.50}
\gL{c000000}{1.00}{72.00}{270.50}{502.00}{270.50}
\gL{c000000}{1.00}{72.00}{288.50}{502.00}{288.50}
\gL{c000000}{1.00}{72.00}{306.50}{502.00}{306.50}
\gL{c000000}{1.00}{72.00}{324.50}{502.00}{324.50}
\gL{c000000}{1.00}{72.00}{342.50}{502.00}{342.50}
\gL{c000000}{1.00}{72.00}{360.50}{502.00}{360.50}
\gL{c000000}{1.00}{72.00}{378.50}{502.00}{378.50}
\gL{c000000}{1.00}{72.00}{396.50}{502.00}{396.50}
\gL{c000000}{1.00}{72.00}{414.50}{502.00}{414.50}
\gL{c000000}{1.00}{72.00}{432.50}{502.00}{432.50}
\gL{c000000}{1.00}{72.00}{450.50}{502.00}{450.50}
\gL{c000000}{1.00}{72.00}{468.50}{502.00}{468.50}
\gL{c000000}{1.00}{72.00}{486.50}{502.00}{486.50}
\gL{c000000}{1.00}{72.00}{504.50}{502.00}{504.50}
\gL{c000000}{1.00}{72.00}{522.50}{502.00}{522.50}
\gL{c000000}{1.00}{72.00}{540.50}{502.00}{540.50}
\gL{c000000}{1.00}{72.00}{558.50}{502.00}{558.50}
\gL{c000000}{1.00}{72.00}{576.50}{502.00}{576.50}
\gL{c000000}{1.00}{72.00}{594.50}{502.00}{594.50}
\gL{c000000}{1.00}{72.00}{612.50}{502.00}{612.50}
\gL{c000000}{1.00}{72.00}{630.50}{502.00}{630.50}
\gL{c000000}{1.00}{72.00}{648.50}{502.00}{648.50}
\gL{c000000}{1.00}{72.00}{666.50}{502.00}{666.50}
\gL{c000000}{1.00}{72.00}{684.50}{502.00}{684.50}
\gL{c000000}{1.00}{72.00}{702.50}{502.00}{702.50}
\gL{c000000}{1.00}{72.00}{720.50}{502.00}{720.50}
\gT{532.01}{43.47}{\gS{R}{8}{c000000}{39 }}
\gT{78.00}{84.21}{\gS{R}{10}{c000000}{Plenum (Press, Inc.), New York, NY, USA }}
\gT{78.00}{102.21}{\gS{R}{10}{c000000}{Popular Mechanics (Press), New York, NY, USA }}
\gT{78.00}{120.21}{\gS{R}{10}{c000000}{(Frederick A.) Praeger (Inc.), New York, NY, USA }}
\gT{78.00}{138.21}{\gS{R}{10}{c000000}{Prentice-Hall (Inc.), Englewood Cliffs, NJ, USA (also, Upper Saddle River, NJ, USA) }}
\gT{78.00}{156.21}{\gS{R}{10}{c000000}{Princeton Univ. Press, Princeton, NJ, USA }}
\gT{78.00}{174.21}{\gS{R}{10}{c000000}{Public Affairs Press, Washington, DC, USA }}
\gT{78.00}{192.21}{\gS{R}{10}{c000000}{Public School (Publishing Co.), Bloomington, IL, USA }}
\gT{78.00}{210.21}{\gS{R}{10}{c000000}{(G. P.) Putnam (’s Sons), New York, NY, USA }}
\gT{78.00}{228.21}{\gS{R}{10}{c000000}{Putnam (\& Co., Ltd.), London, U.K. }}
\gT{78.00}{246.21}{\gS{R}{10}{c000000}{PWS-Kent (Publishing Co.), Boston, MA, USA }}
\gT{78.00}{282.20}{\gS{RB}{10}{c000000}{Q}}
\gT{78.00}{300.20}{\gS{R}{10}{c000000}{Quality (Press), Milwaukee, WI, USA }}
\gT{78.00}{318.20}{\gS{R}{10}{c000000}{Quandt, Leipzig, Germany }}
\gT{78.00}{336.20}{\gS{R}{10}{c000000}{Quelle und Meyer, G.m.b.H., Heidelberg, Germany }}
\gT{78.00}{372.20}{\gS{RB}{10}{c000000}{R}}
\gT{78.00}{390.20}{\gS{R}{10}{c000000}{Rand McNally (\& Co.), Skokie, IL, USA }}
\gT{78.00}{408.20}{\gS{R}{10}{c000000}{Random House (Inc.), New York, NY, USA }}
\gT{78.00}{426.20}{\gS{R}{10}{c000000}{Raven (Press), New York, NY, USA }}
\gT{78.00}{444.20}{\gS{R}{10}{c000000}{(Henry) Regnery (Co.), Chicago, IL, USA }}
\gT{78.00}{462.20}{\gS{R}{10}{c000000}{(D.) Reidel, Amsterdam, The Netherlands }}
\gT{78.00}{480.20}{\gS{R}{10}{c000000}{(The Peter) P. Reilly (Co.), Philadelphia, PA, USA }}
\gT{78.00}{498.20}{\gS{R}{10}{c000000}{Reilly \& Lee (Co.), Chicago, IL, USA }}
\gT{78.00}{516.21}{\gS{R}{10}{c000000}{Reinhold (Publishing Corp.), New York, NY, USA }}
\gT{78.00}{534.21}{\gS{R}{10}{c000000}{Research \& Education (Assoc.), Piscataway, NJ, USA }}
\gT{78.00}{552.21}{\gS{R}{10}{c000000}{(Fleming H.) Revell (Co.), Westwood, NJ, USA }}
\gT{78.00}{570.21}{\gS{R}{10}{c000000}{Reynal (\& Co.), New York, NY, USA }}
\gT{78.00}{588.21}{\gS{R}{10}{c000000}{Reynal \& Hitchcock (Inc.), New York, NY, USA }}
\gT{78.00}{606.21}{\gS{R}{10}{c000000}{Rice Univ. Press, Houston, TX, USA }}
\gT{78.00}{624.21}{\gS{R}{10}{c000000}{Richards Press (Ltd.), London, U.K. }}
\gT{78.00}{642.21}{\gS{R}{10}{c000000}{(John Francis) Rider (Publisher, Inc.), New York, NY, USA }}
\gT{78.00}{660.21}{\gS{R}{10}{c000000}{Rinehart (\& Co.), New York, NY, USA }}
\gT{78.00}{678.21}{\gS{R}{10}{c000000}{Ronald (Press Co.), New York, NY, USA }}
\gT{78.00}{696.21}{\gS{R}{10}{c000000}{Routledge (Peterson \& Co.), Evanston, IL, USA }}
\gT{78.00}{714.21}{\gS{R}{10}{c000000}{Roy (Publishers), New York, NY, USA }}
\end{guidepage}
\begin{guidepage}% page 40
\gF{cEFEFEF}{72.00}{144.00}{501.75}{162.00}
\gL{c000000}{1.00}{78.00}{157.50}{84.00}{157.50}
\gL{c000000}{1.00}{72.50}{72.00}{72.50}{720.00}
\gL{c000000}{1.00}{501.50}{72.00}{501.50}{720.00}
\gL{c000000}{1.00}{72.00}{72.50}{502.00}{72.50}
\gL{c000000}{1.00}{72.00}{90.50}{502.00}{90.50}
\gL{c000000}{1.00}{72.00}{108.50}{502.00}{108.50}
\gL{c000000}{1.00}{72.00}{126.50}{502.00}{126.50}
\gL{c000000}{1.00}{72.00}{144.50}{502.00}{144.50}
\gL{c000000}{1.00}{72.00}{162.50}{502.00}{162.50}
\gL{c000000}{1.00}{72.00}{180.50}{502.00}{180.50}
\gL{c000000}{1.00}{72.00}{198.50}{502.00}{198.50}
\gL{c000000}{1.00}{72.00}{216.50}{502.00}{216.50}
\gL{c000000}{1.00}{72.00}{234.50}{502.00}{234.50}
\gL{c000000}{1.00}{72.00}{252.50}{502.00}{252.50}
\gL{c000000}{1.00}{72.00}{270.50}{502.00}{270.50}
\gL{c000000}{1.00}{72.00}{288.50}{502.00}{288.50}
\gL{c000000}{1.00}{72.00}{306.50}{502.00}{306.50}
\gL{c000000}{1.00}{72.00}{324.50}{502.00}{324.50}
\gL{c000000}{1.00}{72.00}{342.50}{502.00}{342.50}
\gL{c000000}{1.00}{72.00}{360.50}{502.00}{360.50}
\gL{c000000}{1.00}{72.00}{378.50}{502.00}{378.50}
\gL{c000000}{1.00}{72.00}{396.50}{502.00}{396.50}
\gL{c000000}{1.00}{72.00}{414.50}{502.00}{414.50}
\gL{c000000}{1.00}{72.00}{432.50}{502.00}{432.50}
\gL{c000000}{1.00}{72.00}{450.50}{502.00}{450.50}
\gL{c000000}{1.00}{72.00}{468.50}{502.00}{468.50}
\gL{c000000}{1.00}{72.00}{486.50}{502.00}{486.50}
\gL{c000000}{1.00}{72.00}{504.50}{502.00}{504.50}
\gL{c000000}{1.00}{72.00}{522.50}{502.00}{522.50}
\gL{c000000}{1.00}{72.00}{540.50}{502.00}{540.50}
\gL{c000000}{1.00}{72.00}{558.50}{502.00}{558.50}
\gL{c000000}{1.00}{72.00}{576.50}{502.00}{576.50}
\gL{c000000}{1.00}{72.00}{594.50}{502.00}{594.50}
\gL{c000000}{1.00}{72.00}{612.50}{502.00}{612.50}
\gL{c000000}{1.00}{72.00}{630.50}{502.00}{630.50}
\gL{c000000}{1.00}{72.00}{648.50}{502.00}{648.50}
\gL{c000000}{1.00}{72.00}{666.50}{502.00}{666.50}
\gL{c000000}{1.00}{72.00}{684.50}{502.00}{684.50}
\gL{c000000}{1.00}{72.00}{702.50}{502.00}{702.50}
\gL{c000000}{1.00}{72.00}{720.50}{502.00}{720.50}
\gT{532.01}{43.47}{\gS{R}{8}{c000000}{40 }}
\gT{78.00}{84.21}{\gS{R}{10}{c000000}{Russell Sage (Foundation), New York, NY, USA }}
\gT{78.00}{102.21}{\gS{R}{10}{c000000}{Rutgers Univ. Press, New Brunswick, NJ, USA }}
\gT{78.00}{120.21}{\gS{R}{10}{c000000}{Ryerson Press, Toronto, ON, Canada }}
\gT{78.00}{156.21}{\gS{RB}{10}{c000000}{S}}
\gT{78.00}{174.21}{\gS{R}{10}{c000000}{SAE (Soc. Automotive Eng.), Warrendale, PA }}
\gT{78.00}{192.21}{\gS{R}{10}{c000000}{Sage, Newbury Park, CA, USA }}
\gT{78.00}{210.21}{\gS{R}{10}{c000000}{St. Martin’s (Press, Inc.), New York, NY, USA }}
\gT{78.00}{228.21}{\gS{R}{10}{c000000}{(Howard W.) Sams (\& Co.), Indianapolis, IN, USA }}
\gT{78.00}{246.21}{\gS{R}{10}{c000000}{(Benjamin H.) Sanborn (\& Co.), Chicago, IL, USA }}
\gT{78.00}{264.20}{\gS{R}{10}{c000000}{San Francisco Press, San Francisco, CA, USA }}
\gT{78.00}{282.20}{\gS{R}{10}{c000000}{(W. B.) Saunders (Co.), Philadelphia, PA, USA }}
\gT{78.00}{300.20}{\gS{R}{10}{c000000}{Scarecrow (Press, Inc.), New York, NY, USA }}
\gT{78.00}{318.20}{\gS{R}{10}{c000000}{(G.) Schirmer (Inc.), New York, NY, USA }}
\gT{78.00}{336.20}{\gS{R}{10}{c000000}{(Henry) Schuman (Inc., Publishers), New York, NY, USA }}
\gT{78.00}{354.20}{\gS{R}{10}{c000000}{Science Research (Associates, Inc.), Chicago, IL, USA }}
\gT{78.00}{372.20}{\gS{R}{10}{c000000}{Scientific (Press, Inc.), San Francisco, CA, USA }}
\gT{78.00}{390.20}{\gS{R}{10}{c000000}{Scientific Book (Corp.), New York, NY, USA }}
\gT{78.00}{408.20}{\gS{R}{10}{c000000}{SciTech, Rijeka, Croatia }}
\gT{78.00}{426.20}{\gS{R}{10}{c000000}{W. R. Scott (Inc.), New York, NY, USA }}
\gT{78.00}{444.20}{\gS{R}{10}{c000000}{Scott, Foresman (\& Co.), Glenview, IL, USA }}
\gT{78.00}{462.20}{\gS{R}{10}{c000000}{Scott (Publications, Inc.), New York, NY, USA }}
\gT{78.00}{480.20}{\gS{R}{10}{c000000}{(Charles) Scribner, New York, NY, USA }}
\gT{78.00}{498.20}{\gS{R}{10}{c000000}{Seabury (Press, Inc.), New York, NY, USA }}
\gT{78.00}{516.21}{\gS{R}{10}{c000000}{(Martin) Secker \& Warburg (Ltd.), London, U.K. }}
\gT{78.00}{534.21}{\gS{R}{10}{c000000}{Sergel \& Hempel, Brunswick, Germany }}
\gT{78.00}{552.21}{\gS{R}{10}{c000000}{A. W. Shaw (Co.), Chicago, IL, USA }}
\gT{78.00}{570.21}{\gS{R}{10}{c000000}{Shaker, Aachen, Germany }}
\gT{78.00}{588.21}{\gS{R}{10}{c000000}{Sheed \& Ward (Ltd.), London, U.K. }}
\gT{78.00}{606.21}{\gS{R}{10}{c000000}{Sheed \& Ward (Inc.), New York, NY, USA }}
\gT{78.00}{624.21}{\gS{R}{10}{c000000}{Shenandoah (Publishing House, Inc.), Strasburg, VA, USA }}
\gT{78.00}{642.21}{\gS{R}{10}{c000000}{Sheridan (House, Inc.), Dobbs Ferry, NY, USA }}
\gT{78.00}{660.21}{\gS{R}{10}{c000000}{Shoe Lane (Publishing Co.), London, U.K. }}
\gT{78.00}{678.21}{\gS{R}{10}{c000000}{Shoe String Press (Inc.), Hamden, CT, USA }}
\gT{78.00}{696.21}{\gS{R}{10}{c000000}{SIAM (Soc. Ind. Appl. Math.), Philadelphia, PA, USA }}
\gT{78.00}{714.21}{\gS{R}{10}{c000000}{Sidgewick \& Jackson (Ltd.), London, U.K. }}
\end{guidepage}
\begin{guidepage}% page 41
\gF{cEFEFEF}{72.00}{612.00}{501.75}{630.00}
\gL{c000000}{1.00}{78.00}{625.50}{85.00}{625.50}
\gL{c000000}{1.00}{72.50}{72.00}{72.50}{720.00}
\gL{c000000}{1.00}{501.50}{72.00}{501.50}{720.00}
\gL{c000000}{1.00}{72.00}{72.50}{502.00}{72.50}
\gL{c000000}{1.00}{72.00}{90.50}{502.00}{90.50}
\gL{c000000}{1.00}{72.00}{108.50}{502.00}{108.50}
\gL{c000000}{1.00}{72.00}{126.50}{502.00}{126.50}
\gL{c000000}{1.00}{72.00}{144.50}{502.00}{144.50}
\gL{c000000}{1.00}{72.00}{162.50}{502.00}{162.50}
\gL{c000000}{1.00}{72.00}{180.50}{502.00}{180.50}
\gL{c000000}{1.00}{72.00}{198.50}{502.00}{198.50}
\gL{c000000}{1.00}{72.00}{216.50}{502.00}{216.50}
\gL{c000000}{1.00}{72.00}{234.50}{502.00}{234.50}
\gL{c000000}{1.00}{72.00}{252.50}{502.00}{252.50}
\gL{c000000}{1.00}{72.00}{270.50}{502.00}{270.50}
\gL{c000000}{1.00}{72.00}{288.50}{502.00}{288.50}
\gL{c000000}{1.00}{72.00}{306.50}{502.00}{306.50}
\gL{c000000}{1.00}{72.00}{324.50}{502.00}{324.50}
\gL{c000000}{1.00}{72.00}{342.50}{502.00}{342.50}
\gL{c000000}{1.00}{72.00}{360.50}{502.00}{360.50}
\gL{c000000}{1.00}{72.00}{378.50}{502.00}{378.50}
\gL{c000000}{1.00}{72.00}{396.50}{502.00}{396.50}
\gL{c000000}{1.00}{72.00}{414.50}{502.00}{414.50}
\gL{c000000}{1.00}{72.00}{432.50}{502.00}{432.50}
\gL{c000000}{1.00}{72.00}{450.50}{502.00}{450.50}
\gL{c000000}{1.00}{72.00}{468.50}{502.00}{468.50}
\gL{c000000}{1.00}{72.00}{486.50}{502.00}{486.50}
\gL{c000000}{1.00}{72.00}{504.50}{502.00}{504.50}
\gL{c000000}{1.00}{72.00}{522.50}{502.00}{522.50}
\gL{c000000}{1.00}{72.00}{540.50}{502.00}{540.50}
\gL{c000000}{1.00}{72.00}{558.50}{502.00}{558.50}
\gL{c000000}{1.00}{72.00}{576.50}{502.00}{576.50}
\gL{c000000}{1.00}{72.00}{594.50}{502.00}{594.50}
\gL{c000000}{1.00}{72.00}{612.50}{502.00}{612.50}
\gL{c000000}{1.00}{72.00}{630.50}{502.00}{630.50}
\gL{c000000}{1.00}{72.00}{648.50}{502.00}{648.50}
\gL{c000000}{1.00}{72.00}{666.50}{502.00}{666.50}
\gL{c000000}{1.00}{72.00}{684.50}{502.00}{684.50}
\gL{c000000}{1.00}{72.00}{702.50}{502.00}{702.50}
\gL{c000000}{1.00}{72.00}{720.50}{502.00}{720.50}
\gT{532.01}{43.47}{\gS{R}{8}{c000000}{41 }}
\gT{78.00}{84.21}{\gS{R}{10}{c000000}{Silver Burdett (Co.), Morristown, NJ, USA }}
\gT{78.00}{102.21}{\gS{R}{10}{c000000}{Simmons-Boardman (Publishing Corp.), New York, NY, USA }}
\gT{78.00}{120.21}{\gS{R}{10}{c000000}{Simon and Schuster (Inc.), New York, NY, USA }}
\gT{78.00}{138.21}{\gS{R}{10}{c000000}{(L. W.) Singer (Co., Inc.), Syracuse, NY, USA }}
\gT{78.00}{156.21}{\gS{R}{10}{c000000}{(Thomas) Skinner (of Canada, Ltd.), Montreal, PQ, Canada }}
\gT{78.00}{174.21}{\gS{R}{10}{c000000}{(Thomas) Skinner (\& Co.), New York, NY, USA }}
\gT{78.00}{192.21}{\gS{R}{10}{c000000}{Skira (International Corp.), New York, NY, USA }}
\gT{78.00}{210.21}{\gS{R}{10}{c000000}{(William) Sloane (Associates), New York, NY, USA }}
\gT{78.00}{228.21}{\gS{R}{10}{c000000}{SME (Soc. Manufact. Eng.), Dearborn, MI, USA }}
\gT{78.00}{246.21}{\gS{R}{10}{c000000}{Richard R. Smith (Co. Inc.), Peterborough, NH, USA }}
\gT{78.00}{264.20}{\gS{R}{10}{c000000}{Social Sci. Res. Council, New York, NY, USA }}
\gT{78.00}{282.20}{\gS{R}{10}{c000000}{Soncino (Press, Ltd.), London, U.K. }}
\gT{78.00}{300.20}{\gS{R}{10}{c000000}{Southern Methodist Univ. Press, Dallas, TX, USA }}
\gT{78.00}{318.20}{\gS{R}{10}{c000000}{South-Western (Publishing Co.), Cincinnati, OH, USA }}
\gT{78.00}{336.20}{\gS{R}{10}{c000000}{Spartan (Books), Washington, DC, USA }}
\gT{78.00}{354.20}{\gS{R}{10}{c000000}{SPIE (Int. Soc. Opt. Eng.), Bellingham, WA, USA }}
\gT{78.00}{372.20}{\gS{R}{10}{c000000}{Springer-Verlag, New York; Berlin, Germany; Vienna, Austria }}
\gT{78.00}{390.20}{\gS{R}{10}{c000000}{Stackpole (Co.), Harrisburg, PA, USA }}
\gT{78.00}{408.20}{\gS{R}{10}{c000000}{Standard \& Poor’s (Corp.), New York, NY, USA }}
\gT{78.00}{426.20}{\gS{R}{10}{c000000}{Standard (Publishing Co.), Cincinnati, OH, USA }}
\gT{78.00}{444.20}{\gS{R}{10}{c000000}{Stanford Univ. Press, Stanford, CA, USA }}
\gT{78.00}{462.20}{\gS{R}{10}{c000000}{Staples (Press, Ltd.), London, U.K. }}
\gT{78.00}{480.20}{\gS{R}{10}{c000000}{Stechert-Hafner (Inc.), New York, NY, USA }}
\gT{78.00}{498.20}{\gS{R}{10}{c000000}{Steck-Vaughn (Co.), Austin, TX, USA }}
\gT{78.00}{516.21}{\gS{R}{10}{c000000}{Sterling (Publishing Co., Inc.), New York, NY, USA }}
\gT{78.00}{534.21}{\gS{R}{10}{c000000}{Stockton (Press), New York, NY, USA }}
\gT{78.00}{552.21}{\gS{R}{10}{c000000}{(Frederick A.) Stokes (Co.), Philadelphia, PA, USA }}
\gT{78.00}{570.21}{\gS{R}{10}{c000000}{Studio (Publications, Inc.), New York, NY, USA }}
\gT{78.00}{588.21}{\gS{R}{10}{c000000}{Alan Swallow, Denver, CO, USA }}
\gT{78.00}{624.21}{\gS{RB}{10}{c000000}{T}}
\gT{78.00}{642.21}{\gS{R}{10}{c000000}{TAB (Books), New York, NY, USA }}
\gT{78.00}{660.21}{\gS{R}{10}{c000000}{Talbot (Press, Ltd.), Dublin, Ireland }}
\gT{78.00}{678.21}{\gS{R}{10}{c000000}{Taplinger (Publishing Co., Inc.), New York, NY, USA }}
\gT{78.00}{696.21}{\gS{R}{10}{c000000}{Taylor \& Francis, New York, NY, USA }}
\gT{78.00}{714.21}{\gS{R}{10}{c000000}{Technical Press (Ltd.), Surrey, U.K. }}
\end{guidepage}
\begin{guidepage}% page 42
\gF{cEFEFEF}{72.00}{342.00}{501.75}{360.00}
\gL{c000000}{1.00}{78.00}{355.50}{85.00}{355.50}
\gL{c000000}{1.00}{72.50}{72.00}{72.50}{720.00}
\gL{c000000}{1.00}{501.50}{72.00}{501.50}{720.00}
\gL{c000000}{1.00}{72.00}{72.50}{502.00}{72.50}
\gL{c000000}{1.00}{72.00}{90.50}{502.00}{90.50}
\gL{c000000}{1.00}{72.00}{108.50}{502.00}{108.50}
\gL{c000000}{1.00}{72.00}{126.50}{502.00}{126.50}
\gL{c000000}{1.00}{72.00}{144.50}{502.00}{144.50}
\gL{c000000}{1.00}{72.00}{162.50}{502.00}{162.50}
\gL{c000000}{1.00}{72.00}{180.50}{502.00}{180.50}
\gL{c000000}{1.00}{72.00}{198.50}{502.00}{198.50}
\gL{c000000}{1.00}{72.00}{216.50}{502.00}{216.50}
\gL{c000000}{1.00}{72.00}{234.50}{502.00}{234.50}
\gL{c000000}{1.00}{72.00}{252.50}{502.00}{252.50}
\gL{c000000}{1.00}{72.00}{270.50}{502.00}{270.50}
\gL{c000000}{1.00}{72.00}{288.50}{502.00}{288.50}
\gL{c000000}{1.00}{72.00}{306.50}{502.00}{306.50}
\gL{c000000}{1.00}{72.00}{324.50}{502.00}{324.50}
\gL{c000000}{1.00}{72.00}{342.50}{502.00}{342.50}
\gL{c000000}{1.00}{72.00}{360.50}{502.00}{360.50}
\gL{c000000}{1.00}{72.00}{378.50}{502.00}{378.50}
\gL{c000000}{1.00}{72.00}{396.50}{502.00}{396.50}
\gL{c000000}{1.00}{72.00}{414.50}{502.00}{414.50}
\gL{c000000}{1.00}{72.00}{432.50}{502.00}{432.50}
\gL{c000000}{1.00}{72.00}{450.50}{502.00}{450.50}
\gL{c000000}{1.00}{72.00}{468.50}{502.00}{468.50}
\gL{c000000}{1.00}{72.00}{486.50}{502.00}{486.50}
\gL{c000000}{1.00}{72.00}{504.50}{502.00}{504.50}
\gL{c000000}{1.00}{72.00}{522.50}{502.00}{522.50}
\gL{c000000}{1.00}{72.00}{540.50}{502.00}{540.50}
\gL{c000000}{1.00}{72.00}{558.50}{502.00}{558.50}
\gL{c000000}{1.00}{72.00}{576.50}{502.00}{576.50}
\gL{c000000}{1.00}{72.00}{594.50}{502.00}{594.50}
\gL{c000000}{1.00}{72.00}{612.50}{502.00}{612.50}
\gL{c000000}{1.00}{72.00}{630.50}{502.00}{630.50}
\gL{c000000}{1.00}{72.00}{648.50}{502.00}{648.50}
\gL{c000000}{1.00}{72.00}{666.50}{502.00}{666.50}
\gL{c000000}{1.00}{72.00}{684.50}{502.00}{684.50}
\gL{c000000}{1.00}{72.00}{702.50}{502.00}{702.50}
\gL{c000000}{1.00}{72.00}{720.50}{502.00}{720.50}
\gT{532.01}{43.47}{\gS{R}{8}{c000000}{42 }}
\gT{78.00}{84.21}{\gS{R}{10}{c000000}{Textbooks (Ltd.), London, U.K. }}
\gT{78.00}{102.21}{\gS{R}{10}{c000000}{(W.) Thacker (\& Co.), London, U.K. }}
\gT{78.00}{120.21}{\gS{R}{10}{c000000}{Charles C. Thomas (Publishers), Springfield, IL, USA }}
\gT{78.00}{138.21}{\gS{R}{10}{c000000}{Thomas Law Book (Co.), St. Louis, MO, USA }}
\gT{78.00}{156.21}{\gS{R}{10}{c000000}{Thomas Publishing (Co.), New York, NY, USA }}
\gT{78.00}{174.21}{\gS{R}{10}{c000000}{Time-Warner (Inc.), New York, NY, USA }}
\gT{78.00}{192.21}{\gS{R}{10}{c000000}{Transactions Books, New Brunswick, NJ, USA }}
\gT{78.00}{210.21}{\gS{R}{10}{c000000}{Transatlantic Arts (Inc.), New York, NY, USA }}
\gT{78.00}{228.21}{\gS{R}{10}{c000000}{Triangle Books, New York, NY, USA }}
\gT{78.00}{246.21}{\gS{R}{10}{c000000}{Tudor (Publishing Co.), New York, NY, USA }}
\gT{78.00}{264.20}{\gS{R}{10}{c000000}{Tupper \& Love (Inc.), Doraville, GA, USA }}
\gT{78.00}{282.20}{\gS{R}{10}{c000000}{(Charles E.) Tuttle (Co., Inc.), Rutland, VT, USA }}
\gT{78.00}{300.20}{\gS{R}{10}{c000000}{Twayne (Publishers, Inc.), New York, NY, USA }}
\gT{78.00}{318.20}{\gS{R}{10}{c000000}{Twentieth Century (Fund, Inc.), New York, NY, USA }}
\gT{78.00}{354.20}{\gS{RB}{10}{c000000}{U}}
\gT{78.00}{372.20}{\gS{R}{10}{c000000}{Underwriter’s Lab., Chicago, IL, USA }}
\gT{78.00}{390.20}{\gS{R}{10}{c000000}{Frederick Unger (Publishing Co.), New York, NY, USA }}
\gT{78.00}{408.20}{\gS{R}{10}{c000000}{Univelt (Inc.), San Diego, CA }}
\gT{78.00}{426.20}{\gS{R}{10}{c000000}{Univ. of California Press, Berkeley, CA, USA }}
\gT{78.00}{444.20}{\gS{R}{10}{c000000}{Univ. of Chicago Press, Chicago, IL, USA }}
\gT{78.00}{462.20}{\gS{R}{10}{c000000}{Univ. of Connecticut Press, Storrs, CT, USA }}
\gT{78.00}{480.20}{\gS{R}{10}{c000000}{Univ. of Florida Press, Gainesville, FL, USA }}
\gT{78.00}{498.20}{\gS{R}{10}{c000000}{Univ. of Illinois Press, Urbana, IL, USA }}
\gT{78.00}{516.21}{\gS{R}{10}{c000000}{Univ. of Iowa Press, Iowa City, IA, USA }}
\gT{78.00}{534.21}{\gS{R}{10}{c000000}{Univ. of Kentucky Press, Lexington, KY, USA }}
\gT{78.00}{552.21}{\gS{R}{10}{c000000}{Univ. of London Press, (Ltd.), London, U.K. }}
\gT{78.00}{570.21}{\gS{R}{10}{c000000}{Univ. of Michigan Press, Ann Arbor, MI, USA }}
\gT{78.00}{588.21}{\gS{R}{10}{c000000}{Univ. of Minnesota Press, Minneapolis, MN, USA }}
\gT{78.00}{606.21}{\gS{R}{10}{c000000}{Univ. of Missouri Press, Columbia, MO, USA }}
\gT{78.00}{624.21}{\gS{R}{10}{c000000}{Univ. of Nebraska Press, Lincoln, NE, USA }}
\gT{78.00}{642.21}{\gS{R}{10}{c000000}{Univ. of New Mexico Press, Albuquerque, NM, USA }}
\gT{78.00}{660.21}{\gS{R}{10}{c000000}{Univ. of North Carolina Press, Chapel Hill, NC, USA }}
\gT{78.00}{678.21}{\gS{R}{10}{c000000}{Univ. of Notre Dame Press, Notre Dame, IN, USA }}
\gT{78.00}{696.21}{\gS{R}{10}{c000000}{Univ. of Oklahoma Press, Norman, OK, USA }}
\gT{78.00}{714.21}{\gS{R}{10}{c000000}{Univ. of Pennsylvania Press, Philadelphia, PA, USA }}
\end{guidepage}
\begin{guidepage}% page 43
\gF{cEFEFEF}{72.00}{216.00}{501.75}{234.00}
\gF{cEFEFEF}{72.00}{324.00}{501.75}{342.00}
\gL{c000000}{1.00}{78.00}{229.50}{85.00}{229.50}
\gL{c000000}{1.00}{78.00}{337.50}{88.00}{337.50}
\gL{c000000}{1.00}{72.50}{72.00}{72.50}{720.00}
\gL{c000000}{1.00}{501.50}{72.00}{501.50}{720.00}
\gL{c000000}{1.00}{72.00}{72.50}{502.00}{72.50}
\gL{c000000}{1.00}{72.00}{90.50}{502.00}{90.50}
\gL{c000000}{1.00}{72.00}{108.50}{502.00}{108.50}
\gL{c000000}{1.00}{72.00}{126.50}{502.00}{126.50}
\gL{c000000}{1.00}{72.00}{144.50}{502.00}{144.50}
\gL{c000000}{1.00}{72.00}{162.50}{502.00}{162.50}
\gL{c000000}{1.00}{72.00}{180.50}{502.00}{180.50}
\gL{c000000}{1.00}{72.00}{198.50}{502.00}{198.50}
\gL{c000000}{1.00}{72.00}{216.50}{502.00}{216.50}
\gL{c000000}{1.00}{72.00}{234.50}{502.00}{234.50}
\gL{c000000}{1.00}{72.00}{252.50}{502.00}{252.50}
\gL{c000000}{1.00}{72.00}{270.50}{502.00}{270.50}
\gL{c000000}{1.00}{72.00}{288.50}{502.00}{288.50}
\gL{c000000}{1.00}{72.00}{306.50}{502.00}{306.50}
\gL{c000000}{1.00}{72.00}{324.50}{502.00}{324.50}
\gL{c000000}{1.00}{72.00}{342.50}{502.00}{342.50}
\gL{c000000}{1.00}{72.00}{360.50}{502.00}{360.50}
\gL{c000000}{1.00}{72.00}{378.50}{502.00}{378.50}
\gL{c000000}{1.00}{72.00}{396.50}{502.00}{396.50}
\gL{c000000}{1.00}{72.00}{414.50}{502.00}{414.50}
\gL{c000000}{1.00}{72.00}{432.50}{502.00}{432.50}
\gL{c000000}{1.00}{72.00}{450.50}{502.00}{450.50}
\gL{c000000}{1.00}{72.00}{468.50}{502.00}{468.50}
\gL{c000000}{1.00}{72.00}{486.50}{502.00}{486.50}
\gL{c000000}{1.00}{72.00}{504.50}{502.00}{504.50}
\gL{c000000}{1.00}{72.00}{522.50}{502.00}{522.50}
\gL{c000000}{1.00}{72.00}{540.50}{502.00}{540.50}
\gL{c000000}{1.00}{72.00}{558.50}{502.00}{558.50}
\gL{c000000}{1.00}{72.00}{576.50}{502.00}{576.50}
\gL{c000000}{1.00}{72.00}{594.50}{502.00}{594.50}
\gL{c000000}{1.00}{72.00}{612.50}{502.00}{612.50}
\gL{c000000}{1.00}{72.00}{630.50}{502.00}{630.50}
\gL{c000000}{1.00}{72.00}{648.50}{502.00}{648.50}
\gL{c000000}{1.00}{72.00}{666.50}{502.00}{666.50}
\gL{c000000}{1.00}{72.00}{684.50}{502.00}{684.50}
\gL{c000000}{1.00}{72.00}{702.50}{502.00}{702.50}
\gL{c000000}{1.00}{72.00}{720.50}{502.00}{720.50}
\gT{532.01}{43.47}{\gS{R}{8}{c000000}{43 }}
\gT{78.00}{84.21}{\gS{R}{10}{c000000}{University Science (Books), Mill Valley, CA, USA }}
\gT{78.00}{102.21}{\gS{R}{10}{c000000}{Univ. of Pittsburgh Press, Pittsburgh, PA, USA }}
\gT{78.00}{120.21}{\gS{R}{10}{c000000}{Univ. of Texas Press, Austin, TX, USA }}
\gT{78.00}{138.21}{\gS{R}{10}{c000000}{Univ. of Washington Press, Seattle, WA, USA }}
\gT{78.00}{156.21}{\gS{R}{10}{c000000}{Univ. of Wisconsin Press, Madison, WI, USA }}
\gT{78.00}{174.21}{\gS{R}{10}{c000000}{Univ. of Virginia Press, Charlottesville, VA, USA }}
\gT{78.00}{192.21}{\gS{R}{10}{c000000}{(T. Fisher) Unwin, London, U.K. }}
\gT{78.00}{228.21}{\gS{RB}{10}{c000000}{V}}
\gT{78.00}{246.21}{\gS{R}{10}{c000000}{Vanguard (Press, Inc.), New York, NY, USA }}
\gT{78.00}{264.20}{\gS{R}{10}{c000000}{(D.) Van Nostrand (and Co., Inc.), New York, NY, USA }}
\gT{78.00}{282.20}{\gS{R}{10}{c000000}{VCH (Publishers, Inc.), New York, NY, USA }}
\gT{78.00}{300.20}{\gS{R}{10}{c000000}{Viking (Press), New York, NY, USA }}
\gT{78.00}{336.20}{\gS{RB}{10}{c000000}{W}}
\gT{78.00}{354.20}{\gS{R}{10}{c000000}{Wadsworth (Publishing Co.), Belmont, CA, USA }}
\gT{78.00}{372.20}{\gS{R}{10}{c000000}{J. F. Wagner (Inc.), New York, NY, USA }}
\gT{78.00}{390.20}{\gS{R}{10}{c000000}{(George) Wahr (Publishing Co.), Ann Arbor, MI, USA }}
\gT{78.00}{408.20}{\gS{R}{10}{c000000}{(Frederick) Warne (\& Co., Inc.), New York, NY, USA }}
\gT{78.00}{426.20}{\gS{R}{10}{c000000}{Warwick \& York (Inc.), Baltimore, MD, USA }}
\gT{78.00}{444.20}{\gS{R}{10}{c000000}{(Ives) Washburn (Inc.), New York, NY, USA }}
\gT{78.00}{462.20}{\gS{R}{10}{c000000}{Watson-Guptill (Publications, Inc.), New York, NY, USA }}
\gT{78.00}{480.20}{\gS{R}{10}{c000000}{(C. A.) Watts (\& Co., Ltd.), London, U.K. }}
\gT{78.00}{498.20}{\gS{R}{10}{c000000}{Franklin Watts (Inc.), New York, NY, USA }}
\gT{78.00}{516.21}{\gS{R}{10}{c000000}{Wayne State Univ. Press, Detroit, MI, USA }}
\gT{78.00}{534.21}{\gS{R}{10}{c000000}{Webb (Publishing Co.), St. Paul, MN, USA }}
\gT{78.00}{552.21}{\gS{R}{10}{c000000}{Webster (Publishing Co.), St. Louis, MO, USA }}
\gT{78.00}{570.21}{\gS{R}{10}{c000000}{Wellesley-Cambridge (Press), Cambridge, MA, USA }}
\gT{78.00}{588.21}{\gS{R}{10}{c000000}{West (Publishing Co.), St. Paul, MN, USA }}
\gT{78.00}{606.21}{\gS{R}{10}{c000000}{Westminster (Press), Philadelphia, PA, USA }}
\gT{78.00}{624.21}{\gS{R}{10}{c000000}{Westview (Press, Inc.), Boulder, CO, USA }}
\gT{78.00}{642.21}{\gS{R}{10}{c000000}{Whiteside (Inc.), New York, NY, USA }}
\gT{78.00}{660.21}{\gS{R}{10}{c000000}{A. Whitman (\& Co.), Chicago, IL, USA }}
\gT{78.00}{678.21}{\gS{R}{10}{c000000}{Whitman (Publishing Co.), Racine, WI, USA }}
\gT{78.00}{696.21}{\gS{R}{10}{c000000}{Whittlesey House, New York, NY, USA }}
\gT{78.00}{714.21}{\gS{R}{10}{c000000}{W. A. Wilde (Co.), Natick, MA, USA }}
\end{guidepage}
\begin{guidepage}% page 44
\gF{cEFEFEF}{72.00}{360.00}{501.75}{378.00}
\gF{cEFEFEF}{72.00}{468.00}{501.75}{486.00}
\gL{c000000}{1.00}{78.00}{373.50}{85.00}{373.50}
\gL{c000000}{1.00}{78.00}{481.50}{85.00}{481.50}
\gL{c000000}{1.00}{72.50}{72.00}{72.50}{504.00}
\gL{c000000}{1.00}{501.50}{72.00}{501.50}{504.00}
\gL{c000000}{1.00}{72.00}{72.50}{502.00}{72.50}
\gL{c000000}{1.00}{72.00}{90.50}{502.00}{90.50}
\gL{c000000}{1.00}{72.00}{108.50}{502.00}{108.50}
\gL{c000000}{1.00}{72.00}{126.50}{502.00}{126.50}
\gL{c000000}{1.00}{72.00}{144.50}{502.00}{144.50}
\gL{c000000}{1.00}{72.00}{162.50}{502.00}{162.50}
\gL{c000000}{1.00}{72.00}{180.50}{502.00}{180.50}
\gL{c000000}{1.00}{72.00}{198.50}{502.00}{198.50}
\gL{c000000}{1.00}{72.00}{216.50}{502.00}{216.50}
\gL{c000000}{1.00}{72.00}{234.50}{502.00}{234.50}
\gL{c000000}{1.00}{72.00}{252.50}{502.00}{252.50}
\gL{c000000}{1.00}{72.00}{270.50}{502.00}{270.50}
\gL{c000000}{1.00}{72.00}{288.50}{502.00}{288.50}
\gL{c000000}{1.00}{72.00}{306.50}{502.00}{306.50}
\gL{c000000}{1.00}{72.00}{324.50}{502.00}{324.50}
\gL{c000000}{1.00}{72.00}{342.50}{502.00}{342.50}
\gL{c000000}{1.00}{72.00}{360.50}{502.00}{360.50}
\gL{c000000}{1.00}{72.00}{378.50}{502.00}{378.50}
\gL{c000000}{1.00}{72.00}{396.50}{502.00}{396.50}
\gL{c000000}{1.00}{72.00}{414.50}{502.00}{414.50}
\gL{c000000}{1.00}{72.00}{432.50}{502.00}{432.50}
\gL{c000000}{1.00}{72.00}{450.50}{502.00}{450.50}
\gL{c000000}{1.00}{72.00}{468.50}{502.00}{468.50}
\gL{c000000}{1.00}{72.00}{486.50}{502.00}{486.50}
\gL{c000000}{1.00}{72.00}{504.50}{502.00}{504.50}
\gL{c000000}{1.00}{72.00}{626.50}{287.00}{626.50}
\gT{532.01}{43.47}{\gS{R}{8}{c000000}{44 }}
\gT{78.00}{84.21}{\gS{R}{10}{c000000}{(John) Wiley (\& Sons, Inc.), Hoboken, NJ, USA }}
\gT{78.00}{102.21}{\gS{R}{10}{c000000}{Willett, Clark (\& Co.), Chicago, IL, USA }}
\gT{78.00}{120.21}{\gS{R}{10}{c000000}{Williams \& Nargate (Ltd.), London, U.K. }}
\gT{78.00}{138.21}{\gS{R}{10}{c000000}{Williams \& Wilkins (Co.), Baltimore, MD, USA }}
\gT{78.00}{156.21}{\gS{R}{10}{c000000}{H. W. Wilson (Co.), New York, NY, USA }}
\gT{78.00}{174.21}{\gS{R}{10}{c000000}{(William H.) Wire (\& Co., Inc.), Union City, NJ, USA }}
\gT{78.00}{192.21}{\gS{R}{10}{c000000}{Wise (Book Co.), New York, NY, USA }}
\gT{78.00}{210.21}{\gS{R}{10}{c000000}{(M.) Witmark (\& Sons), New York, NY, USA }}
\gT{78.00}{228.21}{\gS{R}{10}{c000000}{(George) Wittenborn (Inc.), New York, NY, USA }}
\gT{78.00}{246.21}{\gS{R}{10}{c000000}{(William) Wood (\& Co.), Baltimore, MD, USA }}
\gT{78.00}{264.20}{\gS{R}{10}{c000000}{World Book (Co.), New York, NY, USA }}
\gT{78.00}{282.20}{\gS{R}{10}{c000000}{World (Publishing Co.), Cleveland, OH, USA }}
\gT{78.00}{300.20}{\gS{R}{10}{c000000}{World Scientific, Singapore }}
\gT{78.00}{318.20}{\gS{R}{10}{c000000}{WSES, Crete, Greece }}
\gT{78.00}{336.20}{\gS{R}{10}{c000000}{(A. A.) Wyn (Inc.), New York, NY, USA }}
\gT{78.00}{372.20}{\gS{RB}{10}{c000000}{Y}}
\gT{78.00}{390.20}{\gS{R}{10}{c000000}{Yale Univ. Press, New Haven, CT, USA }}
\gT{78.00}{408.20}{\gS{R}{10}{c000000}{Year (Inc.), New York, NY, USA }}
\gT{78.00}{426.20}{\gS{R}{10}{c000000}{Year Book Medical (Publishers, Inc.), Chicago, IL, USA }}
\gT{78.00}{444.20}{\gS{R}{10}{c000000}{(Thomas) Yoseloff (Inc.), New York, NY, USA }}
\gT{78.00}{480.20}{\gS{RB}{10}{c000000}{Z}}
\gT{78.00}{498.20}{\gS{R}{10}{c000000}{Ziff-Davis (Publishing Co.), Chicago, IL, USA }}
\gT{85.50}{578.75}{\gS{RB}{14}{c000000}{C.}\hspace*{4.40bp}\gS{R}{14}{c000000}{Abbreviations of Periodicals With Non-English Titles }}
\gT{176.02}{623.82}{\gS{RB}{10}{c000000}{A }}
\gT{78.00}{636.57}{\gS{R}{10}{c000000}{Acad. Serbe Sci. Arts Glas Cl. Sci. Tech. }}
\gT{78.00}{648.57}{\gS{R}{10}{c000000}{Acta Acust. }}
\gT{78.00}{660.57}{\gS{R}{10}{c000000}{Acta Astron. (Poland) }}
\gT{78.00}{672.57}{\gS{R}{10}{c000000}{Acta Astron. Sin. (China) }}
\gT{78.00}{684.57}{\gS{R}{10}{c000000}{Acta Astronaut. (U.K.) }}
\gT{78.00}{696.57}{\gS{R}{10}{c000000}{Acta Astrophys. Sin. (China) }}
\gT{78.00}{708.57}{\gS{R}{10}{c000000}{Acta Autom. Sin. }}
\gT{330.00}{610.02}{\gS{R}{10}{c000000}{Acta Cienc. Indica Math. }}
\gT{330.00}{622.02}{\gS{R}{10}{c000000}{Acta Cienc. Indica Phys. }}
\gT{330.00}{634.02}{\gS{R}{10}{c000000}{Acta Crystallogr. A, Found. Crystallogr. }}
\gT{330.00}{646.02}{\gS{R}{10}{c000000}{Acta Crystallogr. B, Struct. Sci. }}
\gT{330.00}{658.02}{\gS{R}{10}{c000000}{Acta Crystallogr. C, Cryst. Struct. Commun. }}
\gT{330.00}{670.02}{\gS{R}{10}{c000000}{Acta Electron. (France) }}
\gT{330.00}{682.02}{\gS{R}{10}{c000000}{Acta Cyberno }}
\gT{330.00}{694.02}{\gS{R}{10}{c000000}{Acta Electron. Sin. (China) }}
\gT{330.00}{706.02}{\gS{R}{10}{c000000}{Acta Geod. Geophys. Montan. Hung. }}
\end{guidepage}
\begin{guidepage}% page 45
\gL{c000000}{1.00}{324.00}{126.50}{539.00}{126.50}
\gL{c000000}{1.00}{324.00}{294.50}{539.00}{294.50}
\gL{c000000}{1.00}{324.00}{355.50}{539.00}{355.50}
\gL{c000000}{1.00}{324.00}{392.50}{539.00}{392.50}
\gT{532.01}{43.47}{\gS{R}{8}{c000000}{45 }}
\gT{78.00}{80.83}{\gS{R}{10}{c000000}{Acta Geophys. Pol. }}
\gT{78.00}{92.83}{\gS{R}{10}{c000000}{Acta Geophys. Sin. (China) }}
\gT{78.00}{104.83}{\gS{R}{10}{c000000}{Acta Geophys. Sin. (USA) }}
\gT{78.00}{116.83}{\gS{R}{10}{c000000}{Acta Metall. }}
\gT{78.00}{128.83}{\gS{R}{10}{c000000}{Acta Mex. Cienc. Tecnol. }}
\gT{78.00}{140.83}{\gS{R}{10}{c000000}{Acta Phys. Hung. }}
\gT{78.00}{152.83}{\gS{R}{10}{c000000}{Acta Phys. Pol. A }}
\gT{78.00}{164.83}{\gS{R}{10}{c000000}{Acta Phys. Pol. B }}
\gT{78.00}{176.83}{\gS{R}{10}{c000000}{Acta Phys. Sin. }}
\gT{78.00}{188.83}{\gS{R}{10}{c000000}{Acta Phys. Slovaca }}
\gT{78.00}{200.83}{\gS{R}{10}{c000000}{Acta Politec. Mex. }}
\gT{78.00}{212.83}{\gS{R}{10}{c000000}{Acta Polytech. Scand. Appl. Phys. Ser. }}
\gT{78.00}{224.83}{\gS{R}{10}{c000000}{Acta Polytech. Scand. Chem. Technol. }}
\gT{78.00}{236.83}{\gS{R}{10}{c000000}{Metall. Ser. }}
\gT{78.00}{248.83}{\gS{R}{10}{c000000}{Acta Polytech. Scand. Electr. Eng. Ser. }}
\gT{78.00}{260.83}{\gS{R}{10}{c000000}{Acta Polytech. Scand. Math. Comput. Sci. Ser. }}
\gT{78.00}{272.83}{\gS{R}{10}{c000000}{Acta Polytech. Scand. Mech. Eng. Ser. }}
\gT{78.00}{284.83}{\gS{R}{10}{c000000}{Acta Seismol. Sin. }}
\gT{78.00}{296.83}{\gS{R}{10}{c000000}{Acta Tech. Acad. Sci. Hung. }}
\gT{78.00}{308.83}{\gS{R}{10}{c000000}{Acta Tech. CSAV }}
\gT{78.00}{320.83}{\gS{R}{10}{c000000}{Acustica }}
\gT{78.00}{332.83}{\gS{R}{10}{c000000}{(AEU) Arch. Elektr. Ubertragung }}
\gT{78.00}{344.83}{\gS{R}{10}{c000000}{Akust. Zh. }}
\gT{78.00}{356.83}{\gS{R}{10}{c000000}{Algorithmica }}
\gT{78.00}{368.83}{\gS{R}{10}{c000000}{Alta Freq. }}
\gT{78.00}{380.83}{\gS{R}{10}{c000000}{An. Acad. Bras. Cienc. }}
\gT{78.00}{392.83}{\gS{R}{10}{c000000}{An. Fis. }}
\gT{78.00}{404.83}{\gS{R}{10}{c000000}{An. Mec. Electr. }}
\gT{78.00}{416.83}{\gS{R}{10}{c000000}{Angew. Inform. }}
\gT{78.00}{428.83}{\gS{R}{10}{c000000}{Ann. Inst. Henri Poincare Phys. Theor. }}
\gT{78.00}{440.83}{\gS{R}{10}{c000000}{Ann. Soc. Sci. Brux. I, Sci. Math. Astron. Phys. }}
\gT{78.00}{452.83}{\gS{R}{10}{c000000}{Arch. Elektr. Uebertrag. (AEU) (before 1971) }}
\gT{78.00}{464.83}{\gS{R}{10}{c000000}{Arch. Elektron. Uebertrag. Tech. (after 1971) }}
\gT{78.00}{476.83}{\gS{R}{10}{c000000}{Arch. Elektrotech. (Poland) }}
\gT{78.00}{488.83}{\gS{R}{10}{c000000}{Arch. Elektrotech. (Germany) }}
\gT{78.00}{500.83}{\gS{R}{10}{c000000}{Ark. Fys. Semin. Trondheim }}
\gT{78.00}{512.83}{\gS{R}{10}{c000000}{Astrofizika }}
\gT{78.00}{524.83}{\gS{R}{10}{c000000}{Astron. Nachr. (Germany) }}
\gT{78.00}{536.83}{\gS{R}{10}{c000000}{Astron. Tidsskr. }}
\gT{78.00}{548.83}{\gS{R}{10}{c000000}{Astron. Vestn. }}
\gT{78.00}{560.83}{\gS{R}{10}{c000000}{Astron. Zh. }}
\gT{78.00}{572.83}{\gS{R}{10}{c000000}{Atomwirtsch.-Atomtech. }}
\gT{78.00}{584.83}{\gS{R}{10}{c000000}{Atti Accad. Sci. Ist. Bologna CI. Sci. Fis. }}
\gT{78.00}{596.83}{\gS{R}{10}{c000000}{Rend. XIII }}
\gT{78.00}{608.83}{\gS{R}{10}{c000000}{Atti Accad. Sci. Torino I, CI. Sci. Fis. Math. Nat. }}
\gT{78.00}{620.83}{\gS{R}{10}{c000000}{Autom. Strum. }}
\gT{78.00}{632.83}{\gS{R}{10}{c000000}{Autom. Tech. Prax. }}
\gT{78.00}{644.83}{\gS{R}{10}{c000000}{Automatica }}
\gT{78.00}{656.83}{\gS{R}{10}{c000000}{Automatie }}
\gT{78.00}{668.83}{\gS{R}{10}{c000000}{Automatika }}
\gT{78.00}{680.83}{\gS{R}{10}{c000000}{Automatisierungstechnik }}
\gT{78.00}{692.83}{\gS{R}{10}{c000000}{Automatizace }}
\gT{78.00}{704.83}{\gS{R}{10}{c000000}{Automedica }}
\gT{78.00}{716.83}{\gS{R}{10}{c000000}{Avtom. Telemekh. }}
\gT{330.00}{80.83}{\gS{R}{10}{c000000}{Avtom. Vychisl. Tekh. }}
\gT{330.00}{92.83}{\gS{R}{10}{c000000}{Avtomatika }}
\gT{330.00}{104.83}{\gS{R}{10}{c000000}{Avtometriya }}
\gT{428.29}{119.83}{\gS{RB}{10}{c000000}{B }}
\gT{330.00}{135.58}{\gS{R}{10}{c000000}{Ber. Bunsenges. Phys. Chem. }}
\gT{330.00}{147.58}{\gS{R}{10}{c000000}{Biofizika }}
\gT{330.00}{159.58}{\gS{R}{10}{c000000}{Biometrika }}
\gT{330.00}{171.58}{\gS{R}{10}{c000000}{Boll. Geofis. Teor. Appl. }}
\gT{330.00}{183.58}{\gS{R}{10}{c000000}{Bull. Acad. Serbe Sci. Arts cl. Sci. Tech. }}
\gT{330.00}{195.58}{\gS{R}{10}{c000000}{Bull. Annu. Soc. Suisse Chronom. Lab. Suisse. }}
\gT{330.00}{207.58}{\gS{R}{10}{c000000}{Rech. Horlog. }}
\gT{330.00}{219.58}{\gS{R}{10}{c000000}{Bull. Cl. Sci. Acad. R. Belg. }}
\gT{330.00}{231.58}{\gS{R}{10}{c000000}{Bull. Dir. Etud. Rech. A }}
\gT{330.00}{243.58}{\gS{R}{10}{c000000}{Bull. Dir. Etud. Rech. B }}
\gT{330.00}{255.58}{\gS{R}{10}{c000000}{Bull. Dir. Etud. Rech. C }}
\gT{330.00}{267.58}{\gS{R}{10}{c000000}{Bull. Liaison Rech. Inform. Autom. }}
\gT{330.00}{279.58}{\gS{R}{10}{c000000}{Bur. Etud. Autom. }}
\gT{428.02}{291.58}{\gS{RB}{10}{c000000}{C }}
\gT{330.00}{304.33}{\gS{R}{10}{c000000}{CFI-Ceram. Forum Int.-Ber. Dtsch. Keram. Ges. }}
\gT{330.00}{316.33}{\gS{R}{10}{c000000}{Chem. Scr. }}
\gT{330.00}{328.33}{\gS{R}{10}{c000000}{Ciel Terre }}
\gT{330.00}{340.33}{\gS{R}{10}{c000000}{Cybernetica }}
\gT{428.02}{352.33}{\gS{RB}{10}{c000000}{D }}
\gT{330.00}{365.08}{\gS{R}{10}{c000000}{Deut. Hydrogr. Z. }}
\gT{330.00}{377.08}{\gS{R}{10}{c000000}{Dokl. Akad. Nauk SSSR }}
\gT{428.29}{389.08}{\gS{RB}{10}{c000000}{E }}
\gT{330.00}{401.83}{\gS{R}{10}{c000000}{Electroacoustique }}
\gT{330.00}{413.83}{\gS{R}{10}{c000000}{Electrochim. Acta }}
\gT{330.00}{425.83}{\gS{R}{10}{c000000}{Electrochim. Metal. }}
\gT{330.00}{437.83}{\gS{R}{10}{c000000}{Elektor Electron. }}
\gT{330.00}{449.83}{\gS{R}{10}{c000000}{Elektr. Bahnen }}
\gT{330.00}{461.83}{\gS{R}{10}{c000000}{Elektr. Energ.-Tech. }}
\gT{330.00}{473.83}{\gS{R}{10}{c000000}{Elektr. Masch. }}
\gT{330.00}{485.83}{\gS{R}{10}{c000000}{Elektr. Stn. }}
\gT{330.00}{497.83}{\gS{R}{10}{c000000}{Elektrichestvo }}
\gT{330.00}{509.83}{\gS{R}{10}{c000000}{Elektrie }}
\gT{330.00}{521.83}{\gS{R}{10}{c000000}{Elektrizitaetswirtschaft }}
\gT{330.00}{533.83}{\gS{R}{10}{c000000}{Elektro }}
\gT{330.00}{545.83}{\gS{R}{10}{c000000}{Elektro-Anz. }}
\gT{330.00}{557.83}{\gS{R}{10}{c000000}{Elektro-Jahr }}
\gT{330.00}{569.83}{\gS{R}{10}{c000000}{Elektrokhimiya }}
\gT{330.00}{581.83}{\gS{R}{10}{c000000}{Elektron }}
\gT{330.00}{593.83}{\gS{R}{10}{c000000}{Elektron Int. }}
\gT{330.00}{605.83}{\gS{R}{10}{c000000}{Elektron. Entwitkl. }}
\gT{330.00}{617.83}{\gS{R}{10}{c000000}{Elektron. Ind }}
\gT{330.00}{629.83}{\gS{R}{10}{c000000}{Elektron. J. }}
\gT{330.00}{641.83}{\gS{R}{10}{c000000}{Elektron. Prax. }}
\gT{330.00}{653.83}{\gS{R}{10}{c000000}{Elektron. Tekh. }}
\gT{330.00}{665.83}{\gS{R}{10}{c000000}{Elektronica }}
\gT{330.00}{677.83}{\gS{R}{10}{c000000}{Elektronik }}
\gT{330.00}{689.83}{\gS{R}{10}{c000000}{Elektronika }}
\gT{330.00}{701.83}{\gS{R}{10}{c000000}{Elektroniker }}
\gT{330.00}{713.83}{\gS{R}{10}{c000000}{Elektronikschau }}
\end{guidepage}
\begin{guidepage}% page 46
\gL{c000000}{1.00}{72.00}{455.50}{287.00}{455.50}
\gL{c000000}{1.00}{72.00}{648.50}{287.00}{648.50}
\gL{c000000}{1.00}{72.00}{697.50}{287.00}{697.50}
\gL{c000000}{1.00}{324.00}{156.50}{539.00}{156.50}
\gL{c000000}{1.00}{324.00}{444.50}{539.00}{444.50}
\gL{c000000}{1.00}{324.00}{469.50}{539.00}{469.50}
\gL{c000000}{1.00}{324.00}{602.50}{539.00}{602.50}
\gT{532.01}{43.47}{\gS{R}{8}{c000000}{46 }}
\gT{78.00}{80.83}{\gS{R}{10}{c000000}{Elektrosvyaz }}
\gT{78.00}{92.83}{\gS{R}{10}{c000000}{Elektrotech. Cas. }}
\gT{78.00}{104.83}{\gS{R}{10}{c000000}{Elektrotech. Inf. Tech. }}
\gT{78.00}{116.83}{\gS{R}{10}{c000000}{Elektrotech. Obz. }}
\gT{78.00}{128.83}{\gS{R}{10}{c000000}{Elektrotechniek }}
\gT{78.00}{140.83}{\gS{R}{10}{c000000}{Elektrotechnik (Czechoslovakia) }}
\gT{78.00}{152.83}{\gS{R}{10}{c000000}{Elektrotechnik (Switzerland) }}
\gT{78.00}{164.83}{\gS{R}{10}{c000000}{Elektrotechnik (Germany) }}
\gT{78.00}{176.83}{\gS{R}{10}{c000000}{Elektrotechnika }}
\gT{78.00}{188.83}{\gS{R}{10}{c000000}{Elektrotehnika, Zagreb }}
\gT{78.00}{200.83}{\gS{R}{10}{c000000}{Elektrotekhnika }}
\gT{78.00}{212.83}{\gS{R}{10}{c000000}{Elektroteknikeren }}
\gT{78.00}{224.83}{\gS{R}{10}{c000000}{Elektrowaerme Int. B. }}
\gT{78.00}{236.83}{\gS{R}{10}{c000000}{Elettrificazione }}
\gT{78.00}{248.83}{\gS{R}{10}{c000000}{Elettron. Oggi }}
\gT{78.00}{260.83}{\gS{R}{10}{c000000}{Elettron. Telecomun. }}
\gT{78.00}{272.83}{\gS{R}{10}{c000000}{Elettrotecnica }}
\gT{78.00}{284.83}{\gS{R}{10}{c000000}{Elek. Med Aktuell Elektron. }}
\gT{78.00}{296.83}{\gS{R}{10}{c000000}{Elteknik }}
\gT{78.00}{308.83}{\gS{R}{10}{c000000}{Energ. Atomtech. }}
\gT{78.00}{320.83}{\gS{R}{10}{c000000}{Energ. Elettr. }}
\gT{78.00}{332.83}{\gS{R}{10}{c000000}{Energetica }}
\gT{78.00}{344.83}{\gS{R}{10}{c000000}{Energetik }}
\gT{78.00}{356.83}{\gS{R}{10}{c000000}{Energetika }}
\gT{78.00}{368.83}{\gS{R}{10}{c000000}{Energetyka }}
\gT{78.00}{380.83}{\gS{R}{10}{c000000}{Energia Nuclear }}
\gT{78.00}{392.83}{\gS{R}{10}{c000000}{Energie Technik (Germany) }}
\gT{78.00}{404.83}{\gS{R}{10}{c000000}{Energie Technik (Switzerland) }}
\gT{78.00}{416.83}{\gS{R}{10}{c000000}{Entropie }}
\gT{78.00}{428.83}{\gS{R}{10}{c000000}{ETZ }}
\gT{78.00}{440.83}{\gS{R}{10}{c000000}{ETZ Arch. }}
\gT{176.57}{452.83}{\gS{RB}{10}{c000000}{F }}
\gT{78.00}{465.58}{\gS{R}{10}{c000000}{Feingeraetetechnik }}
\gT{78.00}{477.58}{\gS{R}{10}{c000000}{Feinw. Tech. Messtech. }}
\gT{78.00}{489.58}{\gS{R}{10}{c000000}{Fert. Tech. Betr. }}
\gT{78.00}{501.58}{\gS{R}{10}{c000000}{Fis. Tecnol. }}
\gT{78.00}{513.58}{\gS{R}{10}{c000000}{Fiz. Khim. Obrab. Mater. }}
\gT{78.00}{525.58}{\gS{R}{10}{c000000}{Fiz. Met. Metalloved }}
\gT{78.00}{537.58}{\gS{R}{10}{c000000}{Fiz. Nizk. Temp. }}
\gT{78.00}{549.58}{\gS{R}{10}{c000000}{Fiz. Plazmy }}
\gT{78.00}{561.58}{\gS{R}{10}{c000000}{Fiz. Tekh. Poluprovodn. }}
\gT{78.00}{573.58}{\gS{R}{10}{c000000}{Fiz. Tverd. Tela }}
\gT{78.00}{585.58}{\gS{R}{10}{c000000}{Fiz.-Khim. Mekh. Mater. }}
\gT{78.00}{597.58}{\gS{R}{10}{c000000}{Fizika }}
\gT{78.00}{609.58}{\gS{R}{10}{c000000}{Forsch.-Ber. Landes Nordrh.-Westfal. }}
\gT{78.00}{621.58}{\gS{R}{10}{c000000}{Frequenz }}
\gT{78.00}{633.58}{\gS{R}{10}{c000000}{Fys. Tidsskr. }}
\gT{175.74}{645.58}{\gS{RB}{10}{c000000}{G }}
\gT{78.00}{658.33}{\gS{R}{10}{c000000}{G. Fis. }}
\gT{78.00}{670.33}{\gS{R}{10}{c000000}{Geliotekhnika }}
\gT{78.00}{682.33}{\gS{R}{10}{c000000}{Geochim. Cosmochim. Acta }}
\gT{175.74}{694.33}{\gS{RB}{10}{c000000}{H }}
\gT{78.00}{707.08}{\gS{R}{10}{c000000}{Haerterei-Tech. Mitt. }}
\gT{330.00}{80.83}{\gS{R}{10}{c000000}{Helv. Chim. Acta }}
\gT{330.00}{92.83}{\gS{R}{10}{c000000}{Helv. Med. Acta }}
\gT{330.00}{104.83}{\gS{R}{10}{c000000}{Helv. Phys. Acta }}
\gT{330.00}{116.83}{\gS{R}{10}{c000000}{Hochfreq. Electroakust }}
\gT{330.00}{128.83}{\gS{R}{10}{c000000}{Hoppe-Seylers Z. Physiol. Chem. }}
\gT{429.68}{152.83}{\gS{RB}{10}{c000000}{I }}
\gT{330.00}{165.58}{\gS{R}{10}{c000000}{Inf. Elektron. }}
\gT{330.00}{177.58}{\gS{R}{10}{c000000}{Inf. Elettron. }}
\gT{330.00}{189.58}{\gS{R}{10}{c000000}{Inform. Forsch. Entwickl. }}
\gT{330.00}{201.58}{\gS{R}{10}{c000000}{Inform. Spektrum }}
\gT{330.00}{213.58}{\gS{R}{10}{c000000}{Inform.-Fachber. }}
\gT{330.00}{225.58}{\gS{R}{10}{c000000}{Informatie }}
\gT{330.00}{237.58}{\gS{R}{10}{c000000}{Informatik }}
\gT{330.00}{249.58}{\gS{R}{10}{c000000}{Informatologia Yugosl. }}
\gT{330.00}{261.58}{\gS{R}{10}{c000000}{Informatyka }}
\gT{330.00}{273.58}{\gS{R}{10}{c000000}{Infowelt }}
\gT{330.00}{285.58}{\gS{R}{10}{c000000}{Ing. Electr. Mec. }}
\gT{330.00}{297.58}{\gS{R}{10}{c000000}{Ing. Mec. Electr. }}
\gT{330.00}{309.58}{\gS{R}{10}{c000000}{Ing.-Arch. }}
\gT{330.00}{321.58}{\gS{R}{10}{c000000}{Inzh.-Fiz. Zh. }}
\gT{330.00}{333.58}{\gS{R}{10}{c000000}{Izmer. Tekh. }}
\gT{330.00}{345.58}{\gS{R}{10}{c000000}{Izv. Akad. Nauk Arm. SSR Ser. Tekh. Nauk }}
\gT{330.00}{357.58}{\gS{R}{10}{c000000}{Izv. Akad. Nauk SSSR Energ. Transp. }}
\gT{330.00}{369.58}{\gS{R}{10}{c000000}{Izv. Akad. Nauk SSSR Fiz. Atmos. Okeana }}
\gT{330.00}{381.58}{\gS{R}{10}{c000000}{Izv. Akad. Nauk SSSR Fiz. Zemli }}
\gT{330.00}{393.58}{\gS{R}{10}{c000000}{Izv. Akad. Nauk SSSR Ser. Fiz. }}
\gT{330.00}{405.58}{\gS{R}{10}{c000000}{Izv. Vyssh. Uchebn. Zaved. Elektromekh. }}
\gT{330.00}{417.58}{\gS{R}{10}{c000000}{Izv. Vyssh. Uchebn. Zaved. Radioelektron. }}
\gT{330.00}{429.58}{\gS{R}{10}{c000000}{Izv. Vyssh. Uchebn. Zaved. Radiofiz. }}
\gT{429.13}{441.58}{\gS{RB}{10}{c000000}{J }}
\gT{330.00}{454.33}{\gS{R}{10}{c222222}{J. Chim Phys. Phys.-Chim Biol.}}
\gT{427.74}{466.33}{\gS{RB}{10}{c222222}{K}}
\gT{330.00}{479.08}{\gS{R}{10}{c000000}{Kernenergie }}
\gT{330.00}{491.08}{\gS{R}{10}{c000000}{Kerntechnik }}
\gT{330.00}{503.08}{\gS{R}{10}{c000000}{Khim. Fiz. }}
\gT{330.00}{515.08}{\gS{R}{10}{c000000}{Kibern. Vychisl. Tekh. }}
\gT{330.00}{527.08}{\gS{R}{10}{c000000}{Kibernetika }}
\gT{330.00}{539.08}{\gS{R}{10}{c000000}{Kristallografiya }}
\gT{330.00}{551.08}{\gS{R}{10}{c000000}{Kvantovaya Elektron. Mosk. }}
\gT{330.00}{563.08}{\gS{R}{10}{c000000}{Kybernetes }}
\gT{330.00}{575.08}{\gS{R}{10}{c000000}{Kybernetika }}
\gT{426.91}{599.08}{\gS{RB}{10}{c000000}{M }}
\gT{330.00}{611.83}{\gS{R}{10}{c000000}{Med. Tek. }}
\gT{330.00}{623.83}{\gS{R}{10}{c000000}{Mekh. Avtom. Proizvod. }}
\gT{330.00}{635.83}{\gS{R}{10}{c000000}{Meres Autom. }}
\gT{330.00}{647.83}{\gS{R}{10}{c000000}{Mesures }}
\gT{330.00}{659.83}{\gS{R}{10}{c000000}{Metallofizika }}
\gT{330.00}{671.83}{\gS{R}{10}{c000000}{Metalloved. Term. Obrab. Met. }}
\gT{330.00}{683.83}{\gS{R}{10}{c000000}{Meteorol. Gidrol. }}
\gT{330.00}{695.83}{\gS{R}{10}{c000000}{Meteorol. Rundsch. }}
\gT{330.00}{707.83}{\gS{R}{10}{c000000}{Metrol. Apl. }}
\end{guidepage}
\begin{guidepage}% page 47
\gL{c000000}{1.00}{72.00}{108.50}{287.00}{108.50}
\gL{c000000}{1.00}{72.00}{252.50}{287.00}{252.50}
\gL{c000000}{1.00}{72.00}{313.50}{287.00}{313.50}
\gL{c000000}{1.00}{72.00}{554.50}{287.00}{554.50}
\gL{c000000}{1.00}{72.00}{711.50}{287.00}{711.50}
\gL{c000000}{1.00}{324.00}{72.50}{539.00}{72.50}
\gL{c000000}{1.00}{324.00}{240.50}{539.00}{240.50}
\gL{c000000}{1.00}{324.00}{529.50}{539.00}{529.50}
\gL{c000000}{1.00}{324.00}{566.50}{539.00}{566.50}
\gT{532.01}{43.47}{\gS{R}{8}{c000000}{47 }}
\gT{78.00}{80.83}{\gS{R}{10}{c000000}{Metrologia }}
\gT{78.00}{92.83}{\gS{R}{10}{c000000}{Medel. Simul. }}
\gT{176.02}{104.83}{\gS{RB}{10}{c000000}{N }}
\gT{78.00}{117.58}{\gS{R}{10}{c000000}{Nachr. Dok. }}
\gT{78.00}{129.58}{\gS{R}{10}{c000000}{Nachr.tech. Elektron. }}
\gT{78.00}{141.58}{\gS{R}{10}{c000000}{Naturwissentchaften }}
\gT{78.00}{153.58}{\gS{R}{10}{c000000}{Neue Tech. }}
\gT{78.00}{165.58}{\gS{R}{10}{c000000}{Neue Tech. Buero }}
\gT{78.00}{177.58}{\gS{R}{10}{c000000}{Nukleonika }}
\gT{78.00}{189.58}{\gS{R}{10}{c000000}{Numer. Math. }}
\gT{78.00}{201.58}{\gS{R}{10}{c000000}{Nuovo Cimento A }}
\gT{78.00}{213.58}{\gS{R}{10}{c000000}{Nuovo Cimento B }}
\gT{78.00}{225.58}{\gS{R}{10}{c000000}{Nuovo Cimento C }}
\gT{78.00}{237.58}{\gS{R}{10}{c000000}{Nuovo Cimento D }}
\gT{175.74}{249.58}{\gS{RB}{10}{c000000}{O }}
\gT{78.00}{262.33}{\gS{R}{10}{c000000}{Okeanologiya }}
\gT{78.00}{274.33}{\gS{R}{10}{c000000}{Opt. Spektrosk. }}
\gT{78.00}{286.33}{\gS{R}{10}{c000000}{Opt.-Mekh. Prom. }}
\gT{78.00}{298.33}{\gS{R}{10}{c000000}{Optik }}
\gT{176.57}{310.33}{\gS{RB}{10}{c000000}{P }}
\gT{78.00}{323.08}{\gS{R}{10}{c000000}{Photogrammetria }}
\gT{78.00}{335.08}{\gS{R}{10}{c000000}{Photonics Spectra }}
\gT{78.00}{347.08}{\gS{R}{10}{c000000}{Pis’ma Astron. Zh. }}
\gT{78.00}{359.08}{\gS{R}{10}{c000000}{Pis’ma Zh. Eksp. Teor. Fiz. }}
\gT{78.00}{371.08}{\gS{R}{10}{c000000}{Pis’ma Zh. Tekh. Fiz. }}
\gT{78.00}{383.08}{\gS{R}{10}{c000000}{Poverkhn., Fiz. Khim. Mekh. }}
\gT{78.00}{395.08}{\gS{R}{10}{c000000}{Pr. Inst. Elektrotech. }}
\gT{78.00}{407.08}{\gS{R}{10}{c000000}{Prib. Sist. Upr. }}
\gT{78.00}{419.08}{\gS{R}{10}{c000000}{Prib. Tekh. Eksp. }}
\gT{78.00}{431.08}{\gS{R}{10}{c000000}{Prikl. Mat. Mekh. }}
\gT{78.00}{443.08}{\gS{R}{10}{c000000}{Prikl. Mekh. }}
\gT{78.00}{455.08}{\gS{R}{10}{c000000}{Probl. Kibern. }}
\gT{78.00}{467.08}{\gS{R}{10}{c000000}{Probl. Peredachi Inf. }}
\gT{78.00}{479.08}{\gS{R}{10}{c000000}{Proc. K. Ned. Akad. Wet. B, Palaeontol. }}
\gT{78.00}{491.08}{\gS{R}{10}{c000000}{Anthropol.  }}
\gT{78.00}{503.08}{\gS{R}{10}{c000000}{Programmirovanie }}
\gT{78.00}{515.08}{\gS{R}{10}{c000000}{Prz. Elektrotech. }}
\gT{78.00}{527.08}{\gS{R}{10}{c000000}{Prz. Telekomun. }}
\gT{78.00}{539.08}{\gS{R}{10}{c000000}{PT/Elektrotech. Elektron. }}
\gT{176.02}{551.08}{\gS{RB}{10}{c000000}{R }}
\gT{78.00}{563.83}{\gS{R}{10}{c000000}{Radio Fernsehen Elektron. }}
\gT{78.00}{575.83}{\gS{R}{10}{c000000}{Radiotekh. Elektron. }}
\gT{78.00}{587.83}{\gS{R}{10}{c000000}{Radiotekhnika Mosk. }}
\gT{78.00}{599.83}{\gS{R}{10}{c000000}{Rev. Acad. Cienc. Zaragoza }}
\gT{78.00}{611.83}{\gS{R}{10}{c000000}{Rev. Electrotec. (Argentina) }}
\gT{78.00}{623.83}{\gS{R}{10}{c000000}{Rev. Electrotec. (Spain) }}
\gT{78.00}{635.83}{\gS{R}{10}{c000000}{Rev. Energ. }}
\gT{78.00}{647.83}{\gS{R}{10}{c000000}{Rev. Esp. Electron. Rev. Geofis. }}
\gT{78.00}{659.83}{\gS{R}{10}{c000000}{Ric. Autom. }}
\gT{78.00}{671.83}{\gS{R}{10}{c000000}{Ric. Spettrosc. }}
\gT{78.00}{683.83}{\gS{R}{10}{c000000}{Robotersysteme }}
\gT{78.00}{695.83}{\gS{R}{10}{c000000}{Rozpr. Electrotech. }}
\gT{176.85}{707.83}{\gS{RB}{10}{c000000}{S }}
\gT{330.00}{81.58}{\gS{R}{10}{c000000}{Sadhana }}
\gT{330.00}{93.58}{\gS{R}{10}{c000000}{Schweiz. Tech. Z. }}
\gT{330.00}{105.58}{\gS{R}{10}{c000000}{Scientia }}
\gT{330.00}{117.58}{\gS{R}{10}{c000000}{Siemens Forsch. Entwickl. Ber. }}
\gT{330.00}{129.58}{\gS{R}{10}{c000000}{Sist. Autom. }}
\gT{330.00}{141.58}{\gS{R}{10}{c000000}{Sitzungsber. Oester. Akad. Wiss. Math.- }}
\gT{330.00}{153.58}{\gS{R}{10}{c000000}{Naturwiss. Kl. Abt. II (Austria) }}
\gT{330.00}{165.58}{\gS{R}{10}{c000000}{Spectrochim. Acta A, Mol. Spectrosc. }}
\gT{330.00}{177.58}{\gS{R}{10}{c000000}{Spectrochim. Acta B, At. Spectrosc. }}
\gT{330.00}{189.58}{\gS{R}{10}{c000000}{Sprache Datenverarb. }}
\gT{330.00}{201.58}{\gS{R}{10}{c000000}{Stanki Instrum. }}
\gT{330.00}{213.58}{\gS{R}{10}{c000000}{Steklo Keram. }}
\gT{330.00}{225.58}{\gS{R}{10}{c000000}{Svetotekhnika }}
\gT{428.29}{237.58}{\gS{RB}{10}{c000000}{T }}
\gT{330.00}{250.33}{\gS{R}{10}{c000000}{TE Int. }}
\gT{330.00}{262.33}{\gS{R}{10}{c000000}{Tech. Bull. Vevey }}
\gT{330.00}{274.33}{\gS{R}{10}{c000000}{Tech. Mitt. Krupp (Engl. Ed.) }}
\gT{330.00}{286.33}{\gS{R}{10}{c000000}{Tech. Mitt. PTT }}
\gT{330.00}{298.33}{\gS{R}{10}{c000000}{Tech. Mitt. RFZ }}
\gT{330.00}{310.33}{\gS{R}{10}{c000000}{Technica }}
\gT{330.00}{322.33}{\gS{R}{10}{c000000}{Tecnica }}
\gT{330.00}{334.33}{\gS{R}{10}{c000000}{Teh. Fiz. }}
\gT{330.00}{346.33}{\gS{R}{10}{c000000}{Tehnika }}
\gT{330.00}{358.33}{\gS{R}{10}{c000000}{Tekh. Elektrodin. }}
\gT{330.00}{370.33}{\gS{R}{10}{c000000}{Tekh. Kibern. }}
\gT{330.00}{382.33}{\gS{R}{10}{c000000}{Tekh. Kino Telev. }}
\gT{330.00}{394.33}{\gS{R}{10}{c000000}{Tekh. Misul }}
\gT{330.00}{406.33}{\gS{R}{10}{c000000}{Telekomunikacije }}
\gT{330.00}{418.33}{\gS{R}{10}{c000000}{Telektronikk }}
\gT{330.00}{430.33}{\gS{R}{10}{c000000}{Teleteknik }}
\gT{330.00}{442.33}{\gS{R}{10}{c000000}{Teor. Mat. Fiz. }}
\gT{330.00}{454.33}{\gS{R}{10}{c000000}{Teploeoergetika }}
\gT{330.00}{466.33}{\gS{R}{10}{c000000}{Teplofiz. Vvs. Temp. }}
\gT{330.00}{478.33}{\gS{R}{10}{c000000}{Tidskr. Dok. }}
\gT{330.00}{490.33}{\gS{R}{10}{c000000}{TN Nachr. }}
\gT{330.00}{502.33}{\gS{R}{10}{c000000}{Toute Electron. }}
\gT{330.00}{514.33}{\gS{R}{10}{c000000}{Tr. Inst. Teor. Astron. }}
\gT{428.02}{526.33}{\gS{RB}{10}{c000000}{U }}
\gT{330.00}{539.08}{\gS{R}{10}{c000000}{Ukr. Fiz. Zh. }}
\gT{330.00}{551.08}{\gS{R}{10}{c000000}{Usp. Fiz. Nauk }}
\gT{428.02}{563.08}{\gS{RB}{10}{c000000}{V }}
\gT{330.00}{575.83}{\gS{R}{10}{c000000}{Vak.-Tech. }}
\gT{330.00}{587.83}{\gS{R}{10}{c000000}{VDE Fachiber. }}
\gT{330.00}{599.83}{\gS{R}{10}{c000000}{VDI Z. }}
\gT{330.00}{611.83}{\gS{R}{10}{c000000}{Vestn. Mashinostr. }}
\gT{330.00}{623.83}{\gS{R}{10}{c000000}{Vestn. Mosk. Univ. 15, Vychisl. Mat. Kibern. }}
\gT{330.00}{635.83}{\gS{R}{10}{c000000}{Vestn. Mosk. Univ. 3, Fiz. Astron. }}
\gT{330.00}{647.83}{\gS{R}{10}{c000000}{Vesti Akad. Navuk BSSR Ser. Fiz. Energ. Navuk }}
\gT{330.00}{659.83}{\gS{R}{10}{c000000}{VGB Kraftwerkstech. (Ger. Ed.) }}
\gT{330.00}{671.83}{\gS{R}{10}{c000000}{Vide Couches Minces }}
\gT{330.00}{683.83}{\gS{R}{10}{c000000}{Vistas Astron. }}
\gT{330.00}{695.83}{\gS{R}{10}{c000000}{Vopr. At. Nauki Tekh. Ser., Fiz. Radiats. }}
\gT{330.00}{707.83}{\gS{R}{10}{c000000}{Povrezhdenii Radiats. Materialoved. }}
\end{guidepage}
\begin{guidepage}% page 48
\gL{c000000}{1.00}{72.00}{108.50}{287.00}{108.50}
\gL{c000000}{1.00}{72.00}{204.50}{287.00}{204.50}
\gL{c000000}{1.00}{72.00}{229.50}{287.00}{229.50}
\gT{532.01}{43.47}{\gS{R}{8}{c000000}{48 }}
\gT{78.00}{80.83}{\gS{R}{10}{c000000}{Vopr. At. Nauki Tekh. Ser., Obshch. Yad. Fiz. }}
\gT{78.00}{92.83}{\gS{R}{10}{c000000}{Vuoto Sci. Tecnol. }}
\gT{174.63}{104.83}{\gS{RB}{10}{c000000}{W }}
\gT{78.00}{117.58}{\gS{R}{10}{c000000}{Wiss. Z. Friedrich-Schiller-Univ. Jena }}
\gT{78.00}{129.58}{\gS{R}{10}{c000000}{Nat.wiss. Reihe }}
\gT{78.00}{141.58}{\gS{R}{10}{c000000}{Wiss. Z. Karl-Marx-Univ. Leipz. Math.- }}
\gT{78.00}{153.58}{\gS{R}{10}{c000000}{Nat.wiss. Reihe }}
\gT{78.00}{165.58}{\gS{R}{10}{c000000}{Wiss. Z. Tech. Hochsch. Ilmenau }}
\gT{78.00}{177.58}{\gS{R}{10}{c000000}{Wiss. Z. Tech. Univ. Dresd. }}
\gT{78.00}{189.58}{\gS{R}{10}{c000000}{Wiss. Z. Tech. Univ. Karl-Marx-Stadt }}
\gT{176.02}{201.58}{\gS{RB}{10}{c000000}{Y }}
\gT{78.00}{214.33}{\gS{R}{10}{c000000}{Yad. Fiz. }}
\gT{176.29}{226.33}{\gS{RB}{10}{c000000}{Z }}
\gT{78.00}{239.08}{\gS{R}{10}{c000000}{Z. Angew. Math. Mech. }}
\gT{78.00}{251.08}{\gS{R}{10}{c000000}{Z. Angew. Math. Phys. }}
\gT{78.00}{263.08}{\gS{R}{10}{c000000}{Z. Met.kd. }}
\gT{78.00}{275.08}{\gS{R}{10}{c000000}{Z. Nat. Forsch. A, Phys. Phys. Chem. Kosmophys. }}
\gT{78.00}{287.08}{\gS{R}{10}{c000000}{Z. Oper. Res. A, Theor. }}
\gT{330.00}{80.83}{\gS{R}{10}{c000000}{Z. Oper. Res. B, Prax. }}
\gT{330.00}{92.83}{\gS{R}{10}{c000000}{Z. Phys. }}
\gT{330.00}{104.83}{\gS{R}{10}{c000000}{Z. Phys. A, At. Nuclei }}
\gT{330.00}{116.83}{\gS{R}{10}{c000000}{Z. Phys. B, Condens. Matter }}
\gT{330.00}{128.83}{\gS{R}{10}{c000000}{Z. Phys. C, Part Fields }}
\gT{330.00}{140.83}{\gS{R}{10}{c000000}{Z. Phys. Chem. Neue Folge }}
\gT{330.00}{152.83}{\gS{R}{10}{c000000}{Z. Phys. Chem., Leipz. }}
\gT{330.00}{164.83}{\gS{R}{10}{c000000}{Z. Phys. D, At. Mol. Clusters }}
\gT{330.00}{176.83}{\gS{R}{10}{c000000}{Zavod. Lab. }}
\gT{330.00}{188.83}{\gS{R}{10}{c000000}{Zh. Eksp. Teor. Fiz. }}
\gT{330.00}{200.83}{\gS{R}{10}{c000000}{Zh. Fiz. Khim. }}
\gT{330.00}{212.83}{\gS{R}{10}{c000000}{Zh. Prikl. Mekh. Tekh. Fiz. }}
\gT{330.00}{224.83}{\gS{R}{10}{c000000}{Zh. Prikl. Spektrosk. }}
\gT{330.00}{236.83}{\gS{R}{10}{c000000}{Zh. Tekh. Fiz. }}
\gT{330.00}{248.83}{\gS{R}{10}{c000000}{Zh. Vychisl. Mat. Mat. Fiz. }}
\gT{330.00}{260.83}{\gS{R}{10}{c000000}{Zisin, J. Seismol. Soc. Jpn. }}
\end{guidepage}

\end{document}
